
\Data\ID{1003163401}
\Updated{2021-12-18 10:12:35}
\Level{A}
\SubArea{Volume and surface area formulas}
\LANGen{
\Question{
Find the volume and the surface area of a cube with the edge length of \( 5\,\mathrm{cm} \).}
{\( V=125\,\mathrm{cm}^3 \), \( S=150\,\mathrm{cm}^2 \)}{\( V=15\,\mathrm{cm}^3 \), \( S=25\,\mathrm{cm}^2 \)}{\( V=75\,\mathrm{cm}^3 \), \( S=150\,\mathrm{cm}^2 \)}{\( V=125\,\mathrm{cm}^3 \), \( S=30\,\mathrm{cm}^2 \)}{}{}}
\LANGpl{
\Question{
Oblicz  objętość i pole powierzchni sześcianu, długość krawędzi sześcianu jest równa  \( 5\,\mathrm{cm} \).}
{\( V=125\,\mathrm{cm}^3 \), \( S=150\,\mathrm{cm}^2 \)}{\( V=15\,\mathrm{cm}^3 \), \( S=25\,\mathrm{cm}^2 \)}{\( V=75\,\mathrm{cm}^3 \), \( S=150\,\mathrm{cm}^2 \)}{\( V=125\,\mathrm{cm}^3 \), \( S=30\,\mathrm{cm}^2 \)}{}{}}
\LANGcs{
\Question{
Vypočítejte objem a povrch krychle s délkou hrany \( 5\,\mathrm{cm} \).}
{\( V=125\,\mathrm{cm}^3 \), \( S=150\,\mathrm{cm}^2 \)}{\( V=15\,\mathrm{cm}^3 \), \( S=25\,\mathrm{cm}^2 \)}{\( V=75\,\mathrm{cm}^3 \), \( S=150\,\mathrm{cm}^2 \)}{\( V=125\,\mathrm{cm}^3 \), \( S=30\,\mathrm{cm}^2 \)}{}{}}
\LANGsk{
\Question{
Vypočítajte objem a povrch kocky s dĺžkou hrany \( 5\,\mathrm{cm} \).}
{\( V=125\,\mathrm{cm}^3 \), \( S=150\,\mathrm{cm}^2 \)}{\( V=15\,\mathrm{cm}^3 \), \( S=25\,\mathrm{cm}^2 \)}{\( V=75\,\mathrm{cm}^3 \), \( S=150\,\mathrm{cm}^2 \)}{\( V=125\,\mathrm{cm}^3 \), \( S=30\,\mathrm{cm}^2 \)}{}{}}
\LANGes{
\Question{
Calcula el volumen y la superficie de un cubo cuyo lado es de \( 5\,\mathrm{cm} \).}
{\( V=125\,\mathrm{cm}^3 \), \( S=150\,\mathrm{cm}^2 \)}{\( V=15\,\mathrm{cm}^3 \), \( S=25\,\mathrm{cm}^2 \)}{\( V=75\,\mathrm{cm}^3 \), \( S=150\,\mathrm{cm}^2 \)}{\( V=125\,\mathrm{cm}^3 \), \( S=30\,\mathrm{cm}^2 \)}{}{}}
\End

\Data\ID{1003163402}
\Updated{2021-01-15 10:01:46}
\Level{A}
\SubArea{Volume and surface area formulas}
\LANGen{
\Question{
The volume of a cube is \( 64\,\mathrm{cm}^3 \). Find the surface area of a cube.}
{\( 96\,\mathrm{cm}^2 \)}{\( 384\,\mathrm{cm}^2 \)}{\( 24\,\mathrm{cm}^2 \)}{\( 192\,\mathrm{cm}^2 \)}{}{}}
\LANGpl{
\Question{
Objętość sześcianu jest równa \( 64\,\mathrm{cm}^3 \). Oblicz  pole powierzchni sześcianu.}
{\( 96\,\mathrm{cm}^2 \)}{\( 384\,\mathrm{cm}^2 \)}{\( 24\,\mathrm{cm}^2 \)}{\( 192\,\mathrm{cm}^2 \)}{}{}}
\LANGcs{
\Question{
Objem krychle je \( 64\,\mathrm{cm}^3 \). Určete její povrch.}
{\( 96\,\mathrm{cm}^2 \)}{\( 384\,\mathrm{cm}^2 \)}{\( 24\,\mathrm{cm}^2 \)}{\( 192\,\mathrm{cm}^2 \)}{}{}}
\LANGsk{
\Question{
Objem kocky je \( 64\,\mathrm{cm}^3 \). Určte jej povrch.}
{\( 96\,\mathrm{cm}^2 \)}{\( 384\,\mathrm{cm}^2 \)}{\( 24\,\mathrm{cm}^2 \)}{\( 192\,\mathrm{cm}^2 \)}{}{}}
\LANGes{
\Question{
El volumen del cubo es \( 64\,\mathrm{cm}^3 \). Calcula su superficie.}
{\( 96\,\mathrm{cm}^2 \)}{\( 384\,\mathrm{cm}^2 \)}{\( 24\,\mathrm{cm}^2 \)}{\( 192\,\mathrm{cm}^2 \)}{}{}}
\End

\Data\ID{1003163404}
\Updated{2021-12-18 10:12:56}
\Level{A}
\SubArea{Volume and surface area formulas}
\LANGen{
\Question{
The numerical value of the volume of a cube equals to the numerical value of the surface area of the cube. What is the numerical value of the length of the cube’s edge?}
{\( 6 \)}{\( \sqrt6 \)}{\( 6\sqrt6 \)}{\( 1 \)}{}{}}
\LANGpl{
\Question{
Wartość liczbowa objętości sześcianu jest równa wartości liczbowej pola powierzchni sześcianu. Jaka jest wartość liczbowa długości krawędzi sześcianu?}
{\( 6 \)}{\( \sqrt6 \)}{\( 6\sqrt6 \)}{\( 1 \)}{}{}}
\LANGcs{
\Question{
Numericky (číselně) je objem krychle roven jejímu povrchu. Jaká je délka její hrany (v jednotkách délky)?}
{\( 6 \)}{\( \sqrt6 \)}{\( 6\sqrt6 \)}{\( 1 \)}{}{}}
\LANGsk{
\Question{
Numericky (číselne) je objem kocky rovný jej povrchu. Aká je dĺžka jej hrany (v jednotkách dĺžky)?}
{\( 6 \)}{\( \sqrt6 \)}{\( 6\sqrt6 \)}{\( 1 \)}{}{}}
\LANGes{
\Question{
El valor numérico de un cubo es igual a su superficie. ¿Cuánto mide su lado?}
{\( 6 \)}{\( \sqrt6 \)}{\( 6\sqrt6 \)}{\( 1 \)}{}{}}
\End

\Data\ID{1003163406}
\Updated{2021-01-15 10:01:59}
\Level{A}
\SubArea{Volume and surface area formulas}
\LANGen{
\Question{
The volume of a cube is \( 1\,\mathrm{l} \). Find the edge length of this cube.}
{\( 10\,\mathrm{cm} \)}{\( \sqrt{10}\,\mathrm{cm} \)}{\( 10\sqrt{10}\,\mathrm{cm} \)}{\( 1\,\mathrm{cm} \)}{}{}}
\LANGpl{
\Question{
Objętość sześcianu jest równa \( 1\,\mathrm{l} \). Wskaż długość jego krawędzi.}
{\( 10\,\mathrm{cm} \)}{\( \sqrt{10}\,\mathrm{cm} \)}{\( 10\sqrt{10}\,\mathrm{cm} \)}{\( 1\,\mathrm{cm} \)}{}{}}
\LANGcs{
\Question{
Objem krychle je \( 1\,\mathrm{l} \). Určete délku její hrany.}
{\( 10\,\mathrm{cm} \)}{\( \sqrt{10}\,\mathrm{cm} \)}{\( 10\sqrt{10}\,\mathrm{cm} \)}{\( 1\,\mathrm{cm} \)}{}{}}
\LANGsk{
\Question{
Objem kocky je \( 1\,\mathrm{l} \). Určte dĺžku jej hrany.}
{\( 10\,\mathrm{cm} \)}{\( \sqrt{10}\,\mathrm{cm} \)}{\( 10\sqrt{10}\,\mathrm{cm} \)}{\( 1\,\mathrm{cm} \)}{}{}}
\LANGes{
\Question{
El volumen del cubo es \( 1\,\mathrm{l} \). Averigua la longitud de su lado.}
{\( 10\,\mathrm{cm} \)}{\( \sqrt{10}\,\mathrm{cm} \)}{\( 10\sqrt{10}\,\mathrm{cm} \)}{\( 1\,\mathrm{cm} \)}{}{}}
\End

\Data\ID{1003163407}
\Updated{2021-12-18 10:12:05}
\Level{A}
\SubArea{Volume and surface area formulas}
\LANGen{
\Question{
The volume of a cube is \( 8\,\mathrm{l} \). Three-quarters of the cube are filled with water. What is the height of the water level in the cube?}
{\( 15\,\mathrm{cm} \)}{\( 7.5\,\mathrm{cm} \)}{\( 16\,\mathrm{cm} \)}{\( 24\,\mathrm{cm} \)}{}{}}
\LANGpl{
\Question{
Objętość sześcianu wynosi \( 8\,\mathrm{l} \). Trzy czwarte sześcianu wypełnione jest wodą. Jaka jest wysokość poziomu wody w sześcianie?}
{\( 15\,\mathrm{cm} \)}{\( 7{,}5\,\mathrm{cm} \)}{\( 16\,\mathrm{cm} \)}{\( 24\,\mathrm{cm} \)}{}{}}
\LANGcs{
\Question{
Krychle mající objem \( 8\,\mathrm{l} \) je naplněna ze tří čtvrtin vodou. Určete výšku hladiny vody v krychli.}
{\( 15\,\mathrm{cm} \)}{\( 7{,}5\,\mathrm{cm} \)}{\( 16\,\mathrm{cm} \)}{\( 24\,\mathrm{cm} \)}{}{}}
\LANGsk{
\Question{
Kocka, ktorá má objem \( 8\,\mathrm{l} \) je naplnená z troch štvrtín vodou. Určte výšku hladiny vody v kocke.}
{\( 15\,\mathrm{cm} \)}{\( 7{,}5\,\mathrm{cm} \)}{\( 16\,\mathrm{cm} \)}{\( 24\,\mathrm{cm} \)}{}{}}
\LANGes{
\Question{
Un cubo tiene un volumen de \( 8\,\mathrm{l} \) y tres cuartas partes están llenas de agua. ¿Qué altura alcanza el agua?}
{\( 15\,\mathrm{cm} \)}{\( 7{,}5\,\mathrm{cm} \)}{\( 16\,\mathrm{cm} \)}{\( 24\,\mathrm{cm} \)}{}{}}
\End

\Data\ID{1003163701}
\Updated{2022-04-23 05:04:11}
\Level{A}
\SubArea{Volume and surface area formulas}
\LANGen{
\Question{
Find the volume and the surface area of a rectangular prism with the edges of lengths \( 8\,\mathrm{cm} \), \( 6\,\mathrm{cm} \), and \( 4\,\mathrm{cm} \).}
{\( V= 192\,\mathrm{cm}^3 \), \( S= 208\,\mathrm{cm}^2 \)}{\( V= 192\,\mathrm{cm}^3 \), \( S= 104\,\mathrm{cm}^2 \)}{\( V= 208\,\mathrm{cm}^3 \), \( S= 192\,\mathrm{cm}^2 \)}{\( V= 192\,\mathrm{cm}^3 \), \( S= 416\,\mathrm{cm}^2 \)}{}{}}
\LANGpl{
\Question{
Oblicz  objętość i pole powierzchni prostopadłościanu o krawędziach \( 8\,\mathrm{cm} \), \( 6\,\mathrm{cm} \), i \( 4\,\mathrm{cm} \).}
{\( V= 192\,\mathrm{cm}^3 \), \( S= 208\,\mathrm{cm}^2 \)}{\( V= 192\,\mathrm{cm}^3 \), \( S= 104\,\mathrm{cm}^2 \)}{\( V= 208\,\mathrm{cm}^3 \), \( S= 192\,\mathrm{cm}^2 \)}{\( V= 192\,\mathrm{cm}^3 \), \( S= 416\,\mathrm{cm}^2 \)}{}{}}
\LANGcs{
\Question{
Vypočítejte objem a povrch kvádru s hranami o délce \( 8\,\mathrm{cm} \), \( 6\,\mathrm{cm} \) a \( 4\,\mathrm{cm} \).}
{\( V= 192\,\mathrm{cm}^3 \), \( S= 208\,\mathrm{cm}^2 \)}{\( V= 192\,\mathrm{cm}^3 \), \( S= 104\,\mathrm{cm}^2 \)}{\( V= 208\,\mathrm{cm}^3 \), \( S= 192\,\mathrm{cm}^2 \)}{\( V= 192\,\mathrm{cm}^3 \), \( S= 416\,\mathrm{cm}^2 \)}{}{}}
\LANGsk{
\Question{
Vypočítajte objem a povrch kvádra s hranami dĺžky \( 8\,\mathrm{cm} \), \( 6\,\mathrm{cm} \) a \( 4\,\mathrm{cm} \).}
{\( V= 192\,\mathrm{cm}^3 \), \( S= 208\,\mathrm{cm}^2 \)}{\( V= 192\,\mathrm{cm}^3 \), \( S= 104\,\mathrm{cm}^2 \)}{\( V= 208\,\mathrm{cm}^3 \), \( S= 192\,\mathrm{cm}^2 \)}{\( V= 192\,\mathrm{cm}^3 \), \( S= 416\,\mathrm{cm}^2 \)}{}{}}
\LANGes{
\Question{
Calcula el volumen y la superficie de un prisma cuyos lados son de \( 8\,\mathrm{cm} \), \( 6\,\mathrm{cm} \) y \( 4\,\mathrm{cm} \).}
{\( V= 192\,\mathrm{cm}^3 \), \( S= 208\,\mathrm{cm}^2 \)}{\( V= 192\,\mathrm{cm}^3 \), \( S= 104\,\mathrm{cm}^2 \)}{\( V= 208\,\mathrm{cm}^3 \), \( S= 192\,\mathrm{cm}^2 \)}{\( V= 192\,\mathrm{cm}^3 \), \( S= 416\,\mathrm{cm}^2 \)}{}{}}
\End

\Data\ID{1003163702}
\Updated{2021-01-15 11:01:08}
\Level{A}
\SubArea{Volume and surface area formulas}
\LANGen{
\Question{
Consider a rectangular prism with the volume of \( 30\,\mathrm{dm}^3 \). Its length is \( 2\,\mathrm{dm} \), its width \( 3\,\mathrm{dm} \). What is its height?}
{\( 5\,\mathrm{dm} \)}{\( 2.5\,\mathrm{dm} \)}{\( 6\,\mathrm{dm} \)}{\( 10\,\mathrm{dm} \)}{}{}}
\LANGpl{
\Question{
Dany jest prostopadłościan o objętości równej \( 30\,\mathrm{dm}^3 \). Długość krawędzi wynosi \( 2\,\mathrm{dm} \), a szerokość  \( 3\,\mathrm{dm} \). Jaka jest jego wysokość?}
{\( 5\,\mathrm{dm} \)}{\( 2{,}5\,\mathrm{dm} \)}{\( 6\,\mathrm{dm} \)}{\( 10\,\mathrm{dm} \)}{}{}}
\LANGcs{
\Question{
Objem kvádru je \( 30\,\mathrm{dm}^3 \). Určete jeho výšku, je-li jeho délka \( 2\,\mathrm{dm} \) a šířka \( 3\,\mathrm{dm} \).}
{\( 5\,\mathrm{dm} \)}{\( 2{,}5\,\mathrm{dm} \)}{\( 6\,\mathrm{dm} \)}{\( 10\,\mathrm{dm} \)}{}{}}
\LANGsk{
\Question{
Objem kvádra je \( 30\,\mathrm{dm}^3 \). Určte jeho výšku, ak je jeho dĺžka \( 2\,\mathrm{dm} \) a šírka \( 3\,\mathrm{dm} \).}
{\( 5\,\mathrm{dm} \)}{\( 2{,}5\,\mathrm{dm} \)}{\( 6\,\mathrm{dm} \)}{\( 10\,\mathrm{dm} \)}{}{}}
\LANGes{
\Question{
El volumen de un prisma es \( 30\,\mathrm{dm}^3 \). Averigua su altura si su longitud es \( 2\,\mathrm{dm} \) y su anchura \( 3\,\mathrm{dm} \).}
{\( 5\,\mathrm{dm} \)}{\( 2{,}5\,\mathrm{dm} \)}{\( 6\,\mathrm{dm} \)}{\( 10\,\mathrm{dm} \)}{}{}}
\End

\Data\ID{1003163703}
\Updated{2021-12-18 10:12:38}
\Level{A}
\SubArea{Volume and surface area formulas}
\LANGen{
\Question{
The volume of a rectangular prism is \( 392\,\mathrm{cm}^3 \), the length of its square base is \( 7\,\mathrm{cm} \). Find the surface area of this prism.}
{\( 322\,\mathrm{cm}^2 \)}{\( 105\,\mathrm{cm}^2 \)}{\( 784\,\mathrm{cm}^2 \)}{\( 210\,\mathrm{cm}^2 \)}{}{}}
\LANGpl{
\Question{
Objętość prostopadłościanu wynosi  \( 392\,\mathrm{cm}^3 \), długość jego podstawy kwadratowej jest równa  \( 7\,\mathrm{cm} \). Oblicz pole powierzchni prostopadłościanu.}
{\( 322\,\mathrm{cm}^2 \)}{\( 105\,\mathrm{cm}^2 \)}{\( 784\,\mathrm{cm}^2 \)}{\( 210\,\mathrm{cm}^2 \)}{}{}}
\LANGcs{
\Question{
Kvádr se čtvercovou podstavou o straně délky \( 7\,\mathrm{cm} \) má objem \( 392\,\mathrm{cm}^3 \). Určete jeho povrch.}
{\( 322\,\mathrm{cm}^2 \)}{\( 105\,\mathrm{cm}^2 \)}{\( 784\,\mathrm{cm}^2 \)}{\( 210\,\mathrm{cm}^2 \)}{}{}}
\LANGsk{
\Question{
Kváder so štvorcovou podstavou so stranou dĺžky \( 7\,\mathrm{cm} \) má objem \( 392\,\mathrm{cm}^3 \). Určte jeho povrch.}
{\( 322\,\mathrm{cm}^2 \)}{\( 105\,\mathrm{cm}^2 \)}{\( 784\,\mathrm{cm}^2 \)}{\( 210\,\mathrm{cm}^2 \)}{}{}}
\LANGes{
\Question{
Un prisma con base cuadrada de lado de \( 7\,\mathrm{cm} \) tiene un volumen de \( 392\,\mathrm{cm}^3 \). Averigua su superficie.}
{\( 322\,\mathrm{cm}^2 \)}{\( 105\,\mathrm{cm}^2 \)}{\( 784\,\mathrm{cm}^2 \)}{\( 210\,\mathrm{cm}^2 \)}{}{}}
\End

\Data\ID{1003163704}
\Updated{2022-04-23 05:04:57}
\Level{A}
\SubArea{Volume and surface area formulas}
\LANGen{
\Question{
A rectangular aquarium has length of \( 50\,\mathrm{cm} \) and width of \( 30\,\mathrm{cm} \). Suppose  we place a decorative stone in the aquarium and the water level rises by \( 4\,\mathrm{cm} \). What is the volume of the stone?}
{\( 6\,\mathrm{dm}^3 \)}{\( 60\,\mathrm{dm}^3 \)}{\( 1.5\,\mathrm{dm}^3 \)}{\( 150\,\mathrm{dm}^3 \)}{}{}}
\LANGpl{
\Question{
Prostokątne akwarium ma długość \( 50\,\mathrm{cm} \) i szerokość \( 30\,\mathrm{cm} \). Jeśli umieścimy  kamień dekoracyjny w akwarium,  poziom wody wzrośnie wtedy o \( 4\,\mathrm{cm} \). Oblicz objętość kamienia dekoracyjnego?}
{\( 6\,\mathrm{dm}^3 \)}{\( 60\,\mathrm{dm}^3 \)}{\( 1{,}5\,\mathrm{dm}^3 \)}{\( 150\,\mathrm{dm}^3 \)}{}{}}
\LANGcs{
\Question{
Akvárium má rozměr dna \( 50\,\mathrm{cm} \) a \( 30\,\mathrm{cm} \). Vložíme-li do něj dekorační kameny, stoupne v něm hladina vody o \( 4\,\mathrm{cm} \). Určete objem vložených kamenů.}
{\( 6\,\mathrm{dm}^3 \)}{\( 60\,\mathrm{dm}^3 \)}{\( 1{,}5\,\mathrm{dm}^3 \)}{\( 150\,\mathrm{dm}^3 \)}{}{}}
\LANGsk{
\Question{
Akvárium má rozmer dna \( 50\,\mathrm{cm} \) a \( 30\,\mathrm{cm} \). Ak do neho vložíme dekoračné kamene, stupne v ňom hladina vody o \( 4\,\mathrm{cm} \). Určte objem vložených kameňov.}
{\( 6\,\mathrm{dm}^3 \)}{\( 60\,\mathrm{dm}^3 \)}{\( 1{,}5\,\mathrm{dm}^3 \)}{\( 150\,\mathrm{dm}^3 \)}{}{}}
\LANGes{
\Question{
La base de un acuario tiene unas medidas de \( 50\,\mathrm{cm} \) y \( 30\,\mathrm{cm} \). Si insertamos piedras decorativas dentro del acuario, el nivel de agua sube \( 4\,\mathrm{cm} \). Averigua el volumen de las piedras decorativas.}
{\( 6\,\mathrm{dm}^3 \)}{\( 60\,\mathrm{dm}^3 \)}{\( 1{,}5\,\mathrm{dm}^3 \)}{\( 150\,\mathrm{dm}^3 \)}{}{}}
\End

\Data\ID{1003163705}
\Updated{2022-11-01 09:11:48}
\Level{A}
\SubArea{Volume and surface area formulas}
\LANGen{
\Question{
Consider  a carton storage box in the shape of a cube with the edge length of \( 60\,\mathrm{cm} \). Suppose  we want to fill this carton box with small paper  boxes  of the dimensions:  \( 20\,\mathrm{cm} \), \( 5\,\mathrm{cm} \), \( 5\,\mathrm{cm} \). How many of small boxes do we  need to fill the big box completely?}
{\( 432 \)}{\( 72 \)}{\( 216 \)}{\( 75 \)}{}{}}
\LANGpl{
\Question{
Posiadamy  pudło do przechowywania w kształcie sześcianu o długości krawędzi \( 60\,\mathrm{cm} \). Chcemy wypełnić to pudło  małymi papierowymi pudełeczkami o wymiarach:  \( 20\,\mathrm{cm} \), \( 5\,\mathrm{cm} \), \( 5\,\mathrm{cm} \). Ile małych pudełek musimy użyć, aby całkowicie wypełnić pudło?}
{\( 432 \)}{\( 72 \)}{\( 216 \)}{\( 75 \)}{}{}}
\LANGcs{
\Question{
Papírová krabice tvaru krychle má délku hrany \( 60\,\mathrm{cm} \). Kolik krabiček o rozměrech \( 20\,\mathrm{cm} \), \( 5\,\mathrm{cm} \) a \( 5\,\mathrm{cm} \) potřebujeme k jejímu úplnému zaplnění?}
{\( 432 \)}{\( 72 \)}{\( 216 \)}{\( 75 \)}{}{}}
\LANGsk{
\Question{
Papierová krabica tvaru kocky má dĺžku hrany \( 60\,\mathrm{cm} \). Koľko krabičiek s rozmermi \( 20\,\mathrm{cm} \), \( 5\,\mathrm{cm} \) a \( 5\,\mathrm{cm} \) potrebujeme k jej úplnému zaplneniu?}
{\( 432 \)}{\( 72 \)}{\( 216 \)}{\( 75 \)}{}{}}
\LANGes{
\Question{
La longitud del lado de un cubo de papel es de \( 60\,\mathrm{cm} \). ¿Cuántas cajas pequeñas de medidas de \( 20\,\mathrm{cm} \), \( 5\,\mathrm{cm} \) y \( 5\,\mathrm{cm} \) necesitamos para llenarlo?}
{\( 432 \)}{\( 72 \)}{\( 216 \)}{\( 75 \)}{}{}}
\End

\Data\ID{1003163706}
\Updated{2022-01-06 08:01:00}
\Level{A}
\SubArea{Volume and surface area formulas}
\LANGen{
\Question{
Let there be a rectangular prism with the length of \( 8\,\mathrm{cm} \), the width of \( 6\,\mathrm{cm} \) and the length of a space diagonal of \( 10\sqrt2\,\mathrm{cm} \). Find the surface area of this prism.}
{\( 376\,\mathrm{cm}^2 \)}{\( 480\,\mathrm{cm}^2 \)}{\( \left(96+280\cdot\sqrt2\right)\,\mathrm{cm}^2 \)}{\( 480\sqrt2\,\mathrm{cm}^2 \)}{}{}}
\LANGpl{
\Question{
Dany jest prostopadłościan o długości \( 8\,\mathrm{cm} \), szerokości \( 6\,\mathrm{cm} \), długość jego przekątnej wynosi \( 10\sqrt2\,\mathrm{cm} \). Oblicz pole powierzchni prostopadłościanu.}
{\( 376\,\mathrm{cm}^2 \)}{\( 480\,\mathrm{cm}^2 \)}{\( \left(96+280\cdot\sqrt2\right)\,\mathrm{cm}^2 \)}{\( 480\sqrt2\,\mathrm{cm}^2 \)}{}{}}
\LANGcs{
\Question{
Kvádr má délku \( 8\,\mathrm{cm} \), šířku \( 6\,\mathrm{cm} \) a délku tělesové úhlopříčky  \( 10\sqrt2\,\mathrm{cm} \). Určete jeho povrch.}
{\( 376\,\mathrm{cm}^2 \)}{\( 480\,\mathrm{cm}^2 \)}{\( \left(96+280\cdot\sqrt2\right)\,\mathrm{cm}^2 \)}{\( 480\sqrt2\,\mathrm{cm}^2 \)}{}{}}
\LANGsk{
\Question{
Kváder má dĺžku \( 8\,\mathrm{cm} \), šírku \( 6\,\mathrm{cm} \) a dĺžku telesovej uhlopriečky  \( 10\sqrt2\,\mathrm{cm} \). Určte jeho povrch.}
{\( 376\,\mathrm{cm}^2 \)}{\( 480\,\mathrm{cm}^2 \)}{\( \left(96+280\cdot\sqrt2\right)\,\mathrm{cm}^2 \)}{\( 480\sqrt2\,\mathrm{cm}^2 \)}{}{}}
\LANGes{
\Question{
Un prisma rectangular tiene \( 8\,\mathrm{cm} \) de longitud ,\( 6\,\mathrm{cm} \) de anchura y la longitud de su diagonal espacial es  \( 10\sqrt2\,\mathrm{cm} \). Averigua su superficie.}
{\( 376\,\mathrm{cm}^2 \)}{\( 480\,\mathrm{cm}^2 \)}{\( \left(96+280\cdot\sqrt2\right)\,\mathrm{cm}^2 \)}{\( 480\sqrt2\,\mathrm{cm}^2 \)}{}{}}
\End

\Data\ID{1003163801}
\Updated{2022-04-23 05:04:57}
\Level{A}
\SubArea{Volume and surface area formulas}
\LANGen{
\Question{
The volume of a rectangular prism is \( 1620\,\mathrm{cm}^3 \), the ratio of its length, width and height is \( 3:4:5 \). Find the surface area of this prism.}
{\( 846\,\mathrm{cm}^2 \)}{\( 423\,\mathrm{cm}^2 \)}{\( 7614\,\mathrm{cm}^2 \)}{\( 3807\,\mathrm{cm}^2 \)}{}{}}
\LANGpl{
\Question{
Objętość prostopadłościanu wynosi \( 1620\,\mathrm{cm}^3 \), stosunek jego długości, szerokości i wysokości wynosi  \( 3:4:5 \). Oblicz pole powierzchni prostopadłościanu.}
{\( 846\,\mathrm{cm}^2 \)}{\( 423\,\mathrm{cm}^2 \)}{\( 7614\,\mathrm{cm}^2 \)}{\( 3807\,\mathrm{cm}^2 \)}{}{}}
\LANGcs{
\Question{
Kvádr má objem \( 1620\,\mathrm{cm}^3 \), jeho rozměry jsou v poměru  \( 3:4:5 \). Určete jeho povrch.}
{\( 846\,\mathrm{cm}^2 \)}{\( 423\,\mathrm{cm}^2 \)}{\( 7614\,\mathrm{cm}^2 \)}{\( 3807\,\mathrm{cm}^2 \)}{}{}}
\LANGsk{
\Question{
Kváder má objem \( 1620\,\mathrm{cm}^3 \), jeho rozmery sú v pomere  \( 3:4:5 \). Určte jeho povrch.}
{\( 846\,\mathrm{cm}^2 \)}{\( 423\,\mathrm{cm}^2 \)}{\( 7614\,\mathrm{cm}^2 \)}{\( 3807\,\mathrm{cm}^2 \)}{}{}}
\LANGes{
\Question{
El volumen de un prisma rectangular es \( 1620\,\mathrm{cm}^3 \), sus medidas están en proporciones  \( 3:4:5 \). Averigua su superficie.}
{\( 846\,\mathrm{cm}^2 \)}{\( 423\,\mathrm{cm}^2 \)}{\( 7614\,\mathrm{cm}^2 \)}{\( 3807\,\mathrm{cm}^2 \)}{}{}}
\End

\Data\ID{1003163802}
\Updated{2022-01-06 08:01:39}
\Level{A}
\SubArea{Volume and surface area formulas}
\LANGen{
\Question{
The surface area of a rectangular prism is \( 832\,\mathrm{cm}^2 \), and the ratio of its length, width and height is \( 2:3:4 \). Find the volume of this prism.}
{\( 1536\,\mathrm{cm}^3 \)}{\( 3072\sqrt2\,\mathrm{cm}^3 \)}{\( 1536\sqrt3\,\mathrm{cm}^3 \)}{\( 98304\,\mathrm{cm}^3 \)}{}{}}
\LANGpl{
\Question{
Pole powierzchni prostopadłościanu wynosi \( 832\,\mathrm{cm}^2 \), a stosunek jego długości, szerokości i wysokości to \( 2:3:4 \). Oblicz objętość prostopadłościanu.}
{\( 1536\,\mathrm{cm}^3 \)}{\( 3072\sqrt2\,\mathrm{cm}^3 \)}{\( 1536\sqrt3\,\mathrm{cm}^3 \)}{\( 98304\,\mathrm{cm}^3 \)}{}{}}
\LANGcs{
\Question{
Kvádr má povrch  \( 832\,\mathrm{cm}^2 \),  jeho rozměry jsou v poměru \( 2:3:4 \). Určete jeho objem.}
{\( 1536\,\mathrm{cm}^3 \)}{\( 3072\sqrt2\,\mathrm{cm}^3 \)}{\( 1536\sqrt3\,\mathrm{cm}^3 \)}{\( 98304\,\mathrm{cm}^3 \)}{}{}}
\LANGsk{
\Question{
Kváder má povrch  \( 832\,\mathrm{cm}^2 \),  jeho rozmery sú v pomere \( 2:3:4 \). Určte jeho objem.}
{\( 1536\,\mathrm{cm}^3 \)}{\( 3072\sqrt2\,\mathrm{cm}^3 \)}{\( 1536\sqrt3\,\mathrm{cm}^3 \)}{\( 98304\,\mathrm{cm}^3 \)}{}{}}
\LANGes{
\Question{
Un prisma rectangular tiene una superficie de  \( 832\,\mathrm{cm}^2 \),  sus medidas están en proporciones \( 2:3:4 \). Averigua su volumen.}
{\( 1536\,\mathrm{cm}^3 \)}{\( 3072\sqrt2\,\mathrm{cm}^3 \)}{\( 1536\sqrt3\,\mathrm{cm}^3 \)}{\( 98304\,\mathrm{cm}^3 \)}{}{}}
\End

\Data\ID{1003163803}
\Updated{2021-01-20 11:01:15}
\Level{A}
\SubArea{Volume and surface area formulas}
\LANGen{
\Question{
How much does the volume of a cube increase if we lengthen its edge two times?}
{\( 8\times \)}{\( 4\times \)}{\( 2\times \)}{\( 16\times \)}{}{}}
\LANGpl{
\Question{
O ile zwiększy się objętość sześcianu, jeśli dwa razy wydłużymy jego krawędź?}
{\( 8\times \)}{\( 4\times \)}{\( 2\times \)}{\( 16\times \)}{}{}}
\LANGcs{
\Question{
Hranu krychle zvětšíme dvakrát. Kolikrát vzroste její objem?}
{\( 8\times \)}{\( 4\times \)}{\( 2\times \)}{\( 16\times \)}{}{}}
\LANGsk{
\Question{
Hranu kocky zväčšíme dvakrát. Koľkokrát sa zväčší jej objem?}
{\( 8\times \)}{\( 4\times \)}{\( 2\times \)}{\( 16\times \)}{}{}}
\LANGes{
\Question{
Vamos a duplicar el lado de un cubo. ¿Cuántas veces aumenta su volumen?}
{\( 8\times \)}{\( 4\times \)}{\( 2\times \)}{\( 16\times \)}{}{}}
\End

\Data\ID{1003163804}
\Updated{2022-01-06 08:01:30}
\Level{A}
\SubArea{Volume and surface area formulas}
\LANGen{
\Question{
By what percentage does the volume of a cube decrease if we shorten its edge by \( 50\,\% \)?}
{\( 87.5\,\% \)}{\( 12.5\,\% \)}{\( 50\,\% \)}{\( 25\,\% \)}{}{}}
\LANGpl{
\Question{
O ile procent zmniejszy się objętość sześcianu jeśli skrócimy jego krawędź o \( 50\,\% \)?}
{\( 87{,}5\,\% \)}{\( 12{,}5\,\% \)}{\( 50\,\% \)}{\( 25\,\% \)}{}{}}
\LANGcs{
\Question{
Hranu krychle zmenšíme o \( 50\,\% \). O kolik procent se zmenší její objem?}
{\( 87{,}5\,\% \)}{\( 12{,}5\,\% \)}{\( 50\,\% \)}{\( 25\,\% \)}{}{}}
\LANGsk{
\Question{
Hranu kocky zmenšíme o \( 50\,\% \). O koľko percent sa zmenší jej objem?}
{\( 87{,}5\,\% \)}{\( 12{,}5\,\% \)}{\( 50\,\% \)}{\( 25\,\% \)}{}{}}
\LANGes{
\Question{
Vamos a disminuir la longitud del lado de un cubo un \( 50\,\% \). ¿En qué porcentaje disminuye su volumen?}
{\( 87{,}5\,\% \)}{\( 12{,}5\,\% \)}{\( 50\,\% \)}{\( 25\,\% \)}{}{}}
\End

\Data\ID{1003163805}
\Updated{2021-01-20 11:01:29}
\Level{A}
\SubArea{Volume and surface area formulas}
\LANGen{
\Question{
How much does the surface area of a cube increase if we lengthen its edge three times?}
{\( 9\times \)}{\( 3\times \)}{\( 54\times \)}{\( 18\times \)}{}{}}
\LANGpl{
\Question{
Ile razy wzrośnie powierzchnia sześcianu jeśli zwiększymy jego krawędź trzykrotnie?}
{\( 9\times \)}{\( 3\times \)}{\( 54\times\)}{\( 18\times \)}{}{}}
\LANGcs{
\Question{
Hranu krychle zvětšíme třikrát. Kolikrát vzroste její povrch?}
{\( 9\times \)}{\( 3\times \)}{\( 54\times \)}{\( 18\times \)}{}{}}
\LANGsk{
\Question{
Hranu kocky zväčšíme trikrát. Koľkokrát sa zväčší jej povrch?}
{\( 9\times \)}{\( 3\times \)}{\( 54\times \)}{\( 18\times \)}{}{}}
\LANGes{
\Question{
Aumentamos el lado de un cubo por tres. ¿Cuántas veces aumenta su superficie?}
{\( 9\times \)}{\( 3\times \)}{\( 54\times \)}{\( 18\times \)}{}{}}
\End

\Data\ID{1003163806}
\Updated{2022-01-06 08:01:12}
\Level{A}
\SubArea{Volume and surface area formulas}
\LANGen{
\Question{
Dimensions of a rectangular prism are in ratio \( 3:4:7 \) and the volume of the prism is \( 672\,\mathrm{cm}^3 \). Find the length of its longest side diagonal.}
{\( 2\sqrt{65}\,\mathrm{cm} \)}{\( 13\sqrt{10}\,\mathrm{cm} \)}{\( 4\sqrt{65}\,\mathrm{cm} \)}{\( 260\,\mathrm{cm} \)}{}{}}
\LANGpl{
\Question{
Wymiary prostopadłościanu są podane w stosunku \( 3:4:7 \) jego objętość jest równa \( 672\,\mathrm{cm}^3 \). Oblicz długość najdłuższej przekątnej ściany bocznej.}
{\( 2\sqrt{65}\,\mathrm{cm} \)}{\( 13\sqrt{10}\,\mathrm{cm} \)}{\( 4\sqrt{65}\,\mathrm{cm} \)}{\( 260\,\mathrm{cm} \)}{}{}}
\LANGcs{
\Question{
Rozměry kvádru jsou v poměru \( 3:4:7 \), jeho objem je \( 672\,\mathrm{cm}^3 \). Určete velikost jeho nejdelší stěnové úhlopříčky.}
{\( 2\sqrt{65}\,\mathrm{cm} \)}{\( 13\sqrt{10}\,\mathrm{cm} \)}{\( 4\sqrt{65}\,\mathrm{cm} \)}{\( 260\,\mathrm{cm} \)}{}{}}
\LANGsk{
\Question{
Rozmery kvádra sú v pomere \( 3:4:7 \), jeho objem je \( 672\,\mathrm{cm}^3 \). Určte veľkosť jeho najdlhšej stenovej uhlopriečky.}
{\( 2\sqrt{65}\,\mathrm{cm} \)}{\( 13\sqrt{10}\,\mathrm{cm} \)}{\( 4\sqrt{65}\,\mathrm{cm} \)}{\( 260\,\mathrm{cm} \)}{}{}}
\LANGes{
\Question{
Las medidas de un prisma rectangular están en proporción \( 3:4:7 \), su volumen es \( 672\,\mathrm{cm}^3 \). Averigua el tamaño de la diagonal de su cara más grande.}
{\( 2\sqrt{65}\,\mathrm{cm} \)}{\( 13\sqrt{10}\,\mathrm{cm} \)}{\( 4\sqrt{65}\,\mathrm{cm} \)}{\( 260\,\mathrm{cm} \)}{}{}}
\End

\Data\ID{1003163807}
\Updated{2022-04-23 05:04:57}
\Level{A}
\SubArea{Volume and surface area formulas}
\LANGen{
\Question{
A rectangular prism has edges of the lengths of \( 2\,\mathrm{dm} \), \( 4\,\mathrm{dm} \) and \( 5\,\mathrm{dm} \). What is the maximum number of cubes with the edge length of \( 5\,\mathrm{cm} \) that can fit into this prism?}
{\( 320 \)}{\( 1600 \)}{\( 32 \)}{\( 160 \)}{}{}}
\LANGpl{
\Question{
Dany jest prostopadłościan o krawędziach  \( 2\,\mathrm{dm} \), \( 4\,\mathrm{dm} \) i \( 5\,\mathrm{dm} \). Ile sześcianów o krawędzi  \( 5\,\mathrm{cm} \) można maksymalnie umieścić w prostopadłościanie?}
{\( 320 \)}{\( 1600 \)}{\( 32 \)}{\( 160 \)}{}{}}
\LANGcs{
\Question{
Jaký je maximální počet krychlí o velikosti hrany  \( 5\,\mathrm{cm} \), které se vejdou do kvádru o rozměrech \( 2\,\mathrm{dm} \), \( 4\,\mathrm{dm} \) a \( 5\,\mathrm{dm} \)?}
{\( 320 \)}{\( 1600 \)}{\( 32 \)}{\( 160 \)}{}{}}
\LANGsk{
\Question{
Aký je maximálny počet kociek s veľkosťou hrany  \( 5\,\mathrm{cm} \), ktoré sa vojdú do kvádra s rozmermi \( 2\,\mathrm{dm} \), \( 4\,\mathrm{dm} \) a \( 5\,\mathrm{dm} \)?}
{\( 320 \)}{\( 1600 \)}{\( 32 \)}{\( 160 \)}{}{}}
\LANGes{
\Question{
¿Cuál es la cantidad máxima de cubos de lado \( 5\,\mathrm{cm} \), que caben dentro de un prisma rectángular de medidas \( 2\,\mathrm{dm} \), \( 4\,\mathrm{dm} \) y \( 5\,\mathrm{dm} \)?}
{\( 320 \)}{\( 1600 \)}{\( 32 \)}{\( 160 \)}{}{}}
\End

\Data\ID{1003163808}
\Updated{2021-01-20 12:01:27}
\Level{A}
\SubArea{Volume and surface area formulas}
\LANGen{
\Question{
A rectangular pool is filled with water \( 120\,\mathrm{cm} \) deep. The volume of water in the pool is \( 24\,\mathrm{m}^3 \). Find the dimensions of the pool bottom, if you know  that its length is by \( 1\,\mathrm{m} \) larger than its width.}
{\( 5\,\mathrm{m} \), \( 4\,\mathrm{m} \)}{\( 4\,\mathrm{m} \), \( 3\,\mathrm{m} \)}{\( 6\,\mathrm{m} \), \( 5\,\mathrm{m} \)}{\( 3.5\,\mathrm{m} \), \( 4.5\,\mathrm{m} \)}{}{}}
\LANGpl{
\Question{
Prostokątny basen wypełniony jest wodą do  \( 120\,\mathrm{cm} \). Objętość wody w basenie wynosi  \( 24\,\mathrm{m}^3 \). Wskaż  wymiary dna basenu, jeśli wiesz, że jego długość jest o  \( 1\,\mathrm{m} \) większa niż jego  szerokość.}
{\( 5\,\mathrm{m} \), \( 4\,\mathrm{m} \)}{\( 4\,\mathrm{m} \), \( 3\,\mathrm{m} \)}{\( 6\,\mathrm{m} \), \( 5\,\mathrm{m} \)}{\( 3{,}5\,\mathrm{m} \), \( 4{,}5\,\mathrm{m} \)}{}{}}
\LANGcs{
\Question{
Bazén tvaru kvádru je naplněný vodou. Objem vody je \( 24\,\mathrm{m}^3 \) a hloubka je \( 120\,\mathrm{cm} \). Určete rozměry dna bazénu, je-li jeho délka o  \( 1\,\mathrm{m} \) větší než jeho šířka.}
{\( 5\,\mathrm{m} \), \( 4\,\mathrm{m} \)}{\( 4\,\mathrm{m} \), \( 3\,\mathrm{m} \)}{\( 6\,\mathrm{m} \), \( 5\,\mathrm{m} \)}{\( 3{,}5\,\mathrm{m} \), \( 4{,}5\,\mathrm{m} \)}{}{}}
\LANGsk{
\Question{
Bazén tvaru kvádra je naplnený vodou. Objem vody je \( 24\,\mathrm{m}^3 \) a hĺbka je \( 120\,\mathrm{cm} \). Určte rozmery dna bazénu, ak je jeho dĺžka o  \( 1\,\mathrm{m} \) väčšia než jeho šírka.}
{\( 5\,\mathrm{m} \), \( 4\,\mathrm{m} \)}{\( 4\,\mathrm{m} \), \( 3\,\mathrm{m} \)}{\( 6\,\mathrm{m} \), \( 5\,\mathrm{m} \)}{\( 3{,}5\,\mathrm{m} \), \( 4{,}5\,\mathrm{m} \)}{}{}}
\LANGes{
\Question{
Una piscina rectangular está llena de agua. El volumen de agua es \( 24\,\mathrm{m}^3 \) y la profundidad es \( 120\,\mathrm{cm} \). Averigua las medidas del fondo de la piscina si su longitud es \( 1\,\mathrm{m} \) mayor que su anchura.}
{\( 5\,\mathrm{m} \), \( 4\,\mathrm{m} \)}{\( 4\,\mathrm{m} \), \( 3\,\mathrm{m} \)}{\( 6\,\mathrm{m} \), \( 5\,\mathrm{m} \)}{\( 3{,}5\,\mathrm{m} \), \( 4{,}5\,\mathrm{m} \)}{}{}}
\End

\Data\ID{1003163809}
\Updated{2022-01-06 08:01:43}
\Level{A}
\SubArea{Volume and surface area formulas}
\LANGen{
\Question{
The meteo stats show that during the last storm there was a rainfall of \( 5\,\mathrm{mm} \). How much water rained down on a rectangular garden sized \( 10\,\mathrm{m} \) by \( 7\,\mathrm{m} \)?}
{\( 350\,\mathrm{l} \)}{\( 0.35\,\mathrm{l} \)}{\( 3.50\,\mathrm{l} \)}{\( 35\,\mathrm{l} \)}{}{}}
\LANGpl{
\Question{
Deszczomierze wskazały, że w czasie ostatniej burzy opady wynosiły \( 5\,\mathrm{mm} \). Ile wody spadło  na prostokątny ogród o wymiarach \( 10\,\mathrm{m} \) i \( 7\,\mathrm{m} \)?}
{\( 350\,\mathrm{l} \)}{\( 0{,}35\,\mathrm{l} \)}{\( 3{,}50\,\mathrm{l} \)}{\( 35\,\mathrm{l} \)}{}{}}
\LANGcs{
\Question{
Meteorologická předpověď oznámila, že v době bouřky napršelo \( 5\,\mathrm{mm} \) srážek. Kolik litrů vody napršelo na zahradu tvaru obdélníku o rozměrech \( 10\,\mathrm{m} \) a \( 7\,\mathrm{m} \)?}
{\( 350\,\mathrm{l} \)}{\( 0{,}35\,\mathrm{l} \)}{\( 3{,}50\,\mathrm{l} \)}{\( 35\,\mathrm{l} \)}{}{}}
\LANGsk{
\Question{
Meteorologická predpoveď oznámila, že cez búrku napršalo \( 5\,\mathrm{mm} \) zrážok. Koľko litrov vody napršalo na záhradu tvaru obdĺžnika s rozmermi \( 10\,\mathrm{m} \) a \( 7\,\mathrm{m} \)?}
{\( 350\,\mathrm{l} \)}{\( 0{,}35\,\mathrm{l} \)}{\( 3{,}50\,\mathrm{l} \)}{\( 35\,\mathrm{l} \)}{}{}}
\LANGes{
\Question{
El pronóstico de tiempo anunció que durante la tormenta llovieron \( 5\,\mathrm{mm} \) . ¿Cuántos litros de agua llovió en un jardín rectangular cuyos lados miden \( 10\,\mathrm{m} \) y \( 7\,\mathrm{m} \)?}
{\( 350\,\mathrm{l} \)}{\( 0{,}35\,\mathrm{l} \)}{\( 3{,}50\,\mathrm{l} \)}{\( 35\,\mathrm{l} \)}{}{}}
\End

\Data\ID{1103163403}
\Updated{2022-01-06 08:01:31}
\Level{A}
\SubArea{Volume and surface area formulas}
\IMG{1103163403.pdf}
\LANGen{
\Question{
The length of a face diagonal of a cube is \( 6\sqrt2\,\mathrm{cm} \). Find the surface area of the cube.}
{\( 216\,\mathrm{cm}^2 \)}{\( 36\,\mathrm{cm}^2 \)}{\( 96\,\mathrm{cm}^2 \)}{\( 72\sqrt2\,\mathrm{cm}^2 \)}{}{}}
\LANGpl{
\Question{
Długość przekątnej ściany sześcianu jest równa \( 6\sqrt2\,\mathrm{cm} \). Oblicz pole powierzchni sześcianu.}
{\( 216\,\mathrm{cm}^2 \)}{\( 36\,\mathrm{cm}^2 \)}{\( 96\,\mathrm{cm}^2 \)}{\( 72\sqrt2\,\mathrm{cm}^2 \)}{}{}}
\LANGcs{
\Question{
Délka stěnové úhlopříčky krychle je \( 6\sqrt2\,\mathrm{cm} \) (viz obrázek). Vypočítejte povrch krychle.}
{\( 216\,\mathrm{cm}^2 \)}{\( 36\,\mathrm{cm}^2 \)}{\( 96\,\mathrm{cm}^2 \)}{\( 72\sqrt2\,\mathrm{cm}^2 \)}{}{}}
\LANGsk{
\Question{
Dĺžka stenovej uhlopriečky kocky je \( 6\sqrt2\,\mathrm{cm} \) (viď obrázok). Vypočítajte povrch kocky.}
{\( 216\,\mathrm{cm}^2 \)}{\( 36\,\mathrm{cm}^2 \)}{\( 96\,\mathrm{cm}^2 \)}{\( 72\sqrt2\,\mathrm{cm}^2 \)}{}{}}
\LANGes{
\Question{
La longitud de la diagonal de una cara lateral es de \( 6\sqrt2\,\mathrm{cm} \) (mira el dibujo). Calcula la superficie del cubo.}
{\( 216\,\mathrm{cm}^2 \)}{\( 36\,\mathrm{cm}^2 \)}{\( 96\,\mathrm{cm}^2 \)}{\( 72\sqrt2\,\mathrm{cm}^2 \)}{}{}}
\End

\Data\ID{1103163405}
\Updated{2022-01-06 08:01:51}
\Level{A}
\SubArea{Volume and surface area formulas}
\IMG{1103163405.pdf}
\LANGen{
\Question{
The length of a space diagonal of a cube is \( 9\,\mathrm{cm} \). Find the volume of the cube.}
{\( 81\sqrt3\,\mathrm{cm}^3 \)}{\( 9\sqrt3\,\mathrm{cm}^3 \)}{\( 27\sqrt3\,\mathrm{cm}^3 \)}{\( 81\,\mathrm{cm}^3 \)}{}{}}
\LANGpl{
\Question{
Długość przekątnej sześcianu to  \( 9\,\mathrm{cm} \). Oblicz objętość sześcianu.}
{\( 81\sqrt3\,\mathrm{cm}^3 \)}{\( 9\sqrt3\,\mathrm{cm}^3 \)}{\( 27\sqrt3\,\mathrm{cm}^3 \)}{\( 81\,\mathrm{cm}^3 \)}{}{}}
\LANGcs{
\Question{
Délka tělesové úhlopříčky krychle je \( 9\,\mathrm{cm} \) (viz obrázek). Určete její objem.}
{\( 81\sqrt3\,\mathrm{cm}^3 \)}{\( 9\sqrt3\,\mathrm{cm}^3 \)}{\( 27\sqrt3\,\mathrm{cm}^3 \)}{\( 81\,\mathrm{cm}^3 \)}{}{}}
\LANGsk{
\Question{
Dĺžka telesovej uhlopriečky kocky je \( 9\,\mathrm{cm} \) (viď obrázok). Určte jej objem.}
{\( 81\sqrt3\,\mathrm{cm}^3 \)}{\( 9\sqrt3\,\mathrm{cm}^3 \)}{\( 27\sqrt3\,\mathrm{cm}^3 \)}{\( 81\,\mathrm{cm}^3 \)}{}{}}
\LANGes{
\Question{
La longitud de la diagonal interior es de \( 9\,\mathrm{cm} \) (miraa el dibujo). Averigua el volumen del cubo.}
{\( 81\sqrt3\,\mathrm{cm}^3 \)}{\( 9\sqrt3\,\mathrm{cm}^3 \)}{\( 27\sqrt3\,\mathrm{cm}^3 \)}{\( 81\,\mathrm{cm}^3 \)}{}{}}
\End

\Data\ID{2010016501}
\Updated{2022-11-02 09:11:35}
\Level{A}
\SubArea{Volume and surface area formulas}
\LANGen{
\Question{
Find the volume and the surface area of a rectangular prism with the edges of lengths \( 3\,\mathrm{cm} \), \( 9\,\mathrm{cm} \) and \( 15\,\mathrm{cm} \).}
{\( V= 405\,\mathrm{cm}^3 \), \( S= 414\,\mathrm{cm}^2 \)}{\( V= 414\,\mathrm{cm}^3 \), \( S= 405\,\mathrm{cm}^2 \)}{\( V= 415\,\mathrm{cm}^3 \), \( S= 404\,\mathrm{cm}^2 \)}{\( V= 42\,\mathrm{cm}^3 \), \( S= 84\,\mathrm{cm}^2 \)}{}{}}
\LANGpl{
\Question{
Znajdź objętość i pole powierzchni prostopadłościanu o krawędziach długości \( 3\,\mathrm{cm} \), \( 9\,\mathrm{cm} \) i \( 15\,\mathrm{cm} \).}
{\( V= 405\,\mathrm{cm}^3 \), \( S= 414\,\mathrm{cm}^2 \)}{\( V= 414\,\mathrm{cm}^3 \), \( S= 405\,\mathrm{cm}^2 \)}{\( V= 415\,\mathrm{cm}^3 \), \( S= 404\,\mathrm{cm}^2 \)}{\( V= 42\,\mathrm{cm}^3 \), \( S= 84\,\mathrm{cm}^2 \)}{}{}}
\LANGcs{
\Question{
Vypočtěte objem a povrch kvádru s hranami o délce \( 3\,\mathrm{cm} \), \( 9\,\mathrm{cm} \) a \( 15\,\mathrm{cm} \).}
{\( V= 405\,\mathrm{cm}^3 \), \( S= 414\,\mathrm{cm}^2 \)}{\( V= 414\,\mathrm{cm}^3 \), \( S= 405\,\mathrm{cm}^2 \)}{\( V= 415\,\mathrm{cm}^3 \), \( S= 404\,\mathrm{cm}^2 \)}{\( V= 42\,\mathrm{cm}^3 \), \( S= 84\,\mathrm{cm}^2 \)}{}{}}
\LANGsk{
\Question{
Vypočítajte objem a povrch kvádra s hranami dĺžky \( 3\,\mathrm{cm} \), \( 9\,\mathrm{cm} \) a \( 15\,\mathrm{cm} \).}
{\( V= 405\,\mathrm{cm}^3 \), \( S= 414\,\mathrm{cm}^2 \)}{\( V= 414\,\mathrm{cm}^3 \), \( S= 405\,\mathrm{cm}^2 \)}{\( V= 415\,\mathrm{cm}^3 \), \( S= 404\,\mathrm{cm}^2 \)}{\( V= 42\,\mathrm{cm}^3 \), \( S= 84\,\mathrm{cm}^2 \)}{}{}}
\LANGes{
\Question{
Halla el volumen y el área total de un prisma rectangular cuyas aristas tienen longitudes \( 3\,\mathrm{cm} \), \( 9\,\mathrm{cm} \) y \( 15\,\mathrm{cm} \).}
{\( V= 405\,\mathrm{cm}^3 \), \( S= 414\,\mathrm{cm}^2 \)}{\( V= 414\,\mathrm{cm}^3 \), \( S= 405\,\mathrm{cm}^2 \)}{\( V= 415\,\mathrm{cm}^3 \), \( S= 404\,\mathrm{cm}^2 \)}{\( V= 42\,\mathrm{cm}^3 \), \( S= 84\,\mathrm{cm}^2 \)}{}{}}
\End

\Data\ID{9000120301}
\Updated{2022-04-23 05:04:57}
\Level{A}
\SubArea{Volume and surface area formulas}
\IMG{001203_01-figure0.pdf}
\LANGen{
\Question{
The length of a solid diagonal in a cube is
\(u = 2\sqrt{6}\, \mathrm{cm}\). Find
the surface area.}
{\(48\, \mathrm{cm}^{2}\)}{\(24\, \mathrm{cm}^{2}\)}{\(24\sqrt{2}\, \mathrm{cm}^{2}\)}{\(16\sqrt{2}\, \mathrm{cm}^{2}\)}{\(12\sqrt{6}\, \mathrm{cm}^{2}\)}{}}
\LANGpl{
\Question{
Długość przekątnej sześcianu jest równa
\(u = 2\sqrt{6}\, \mathrm{cm}\). Oblicz pole powierzchni sześcianu.}
{\(48\, \mathrm{cm}^{2}\)}{\(24\, \mathrm{cm}^{2}\)}{\(24\sqrt{2}\, \mathrm{cm}^{2}\)}{\(16\sqrt{2}\, \mathrm{cm}^{2}\)}{\(12\sqrt{6}\, \mathrm{cm}^{2}\)}{}}
\LANGcs{
\Question{
Délka tělesové úhlopříčky krychle je
\(2\sqrt{6}\, \mathrm{cm}\).
Povrch této krychle je:}
{\(48\, \mathrm{cm}^{2}\)}{\(24\, \mathrm{cm}^{2}\)}{\(24\sqrt{2}\, \mathrm{cm}^{2}\)}{\(16\sqrt{2}\, \mathrm{cm}^{2}\)}{\(12\sqrt{6}\, \mathrm{cm}^{2}\)}{}}
\LANGsk{
\Question{
Dĺžka telesovej uhlopriečky kocky je
\(2\sqrt{6}\, \mathrm{cm}\).
Povrch tejto kocky je:}
{\(48\, \mathrm{cm}^{2}\)}{\(24\, \mathrm{cm}^{2}\)}{\(24\sqrt{2}\, \mathrm{cm}^{2}\)}{\(16\sqrt{2}\, \mathrm{cm}^{2}\)}{\(12\sqrt{6}\, \mathrm{cm}^{2}\)}{}}
\LANGes{
\Question{
La longitud de la diagonal interior de un cubo es de \(2\sqrt{6}\, \mathrm{cm}\).
La superficie de este cubo es:}
{\(48\, \mathrm{cm}^{2}\)}{\(24\, \mathrm{cm}^{2}\)}{\(24\sqrt{2}\, \mathrm{cm}^{2}\)}{\(16\sqrt{2}\, \mathrm{cm}^{2}\)}{\(12\sqrt{6}\, \mathrm{cm}^{2}\)}{}}
\End

\Data\ID{9000120306}
\Updated{2022-04-23 05:04:57}
\Level{A}
\SubArea{Volume and surface area formulas}
\IMG{001203_06-figure0.pdf}
\LANGen{
\Question{
The lengths of a side, face diagonal and solid diagonal through the vertex
\(A\) in a rectangular
box \(ABCDEFGH\) 
are \(|AB| = 6\, \mathrm{cm}\),
\(|AC| = 10\, \mathrm{cm}\),
\(|AG| = 15\, \mathrm{cm}\). Find
the surface area.}
{\(\left (96 + 140\sqrt{5}\right )\, \mathrm{cm}^{2}\)}{\(600\, \mathrm{cm}^{2}\)}{\(236\sqrt{5}\, \mathrm{cm}^{2}\)}{\(\left (48 + 70\sqrt{5}\right )\, \mathrm{cm}^{2}\)}{\(240\sqrt{5}\, \mathrm{cm}^{2}\)}{}}
\LANGpl{
\Question{
Długość boku, przekątnej podstawy prostopadłościanu  i przekątnej prostopadłościanu przechodzącej przez wierzchołek 
\(A\) w prostopadłościanie 
 \(ABCDEFGH\) 
są odpowiednio  równe \(|AB| = 6\, \mathrm{cm}\),
\(|AC| = 10\, \mathrm{cm}\),
\(|AG| = 15\, \mathrm{cm}\). Oblicz pole powierzchni prostopadłościanu.}
{\(\left (96 + 140\sqrt{5}\right )\, \mathrm{cm}^{2}\)}{\(600\, \mathrm{cm}^{2}\)}{\(236\sqrt{5}\, \mathrm{cm}^{2}\)}{\(\left (48 + 70\sqrt{5}\right )\, \mathrm{cm}^{2}\)}{\(240\sqrt{5}\, \mathrm{cm}^{2}\)}{}}
\LANGcs{
\Question{
V kvádru \(ABCDEFGH\) 
platí: \(|AB| = 6\, \mathrm{cm};\ |AC| = 10\, \mathrm{cm};\ |AG| = 15\, \mathrm{cm}\). Povrch tohoto kvádru je:}
{\(96 + 140\sqrt{5}\, \mathrm{cm}^{2}\)}{\(600\, \mathrm{cm}^{2}\)}{\(236\sqrt{5}\, \mathrm{cm}^{2}\)}{\(48 + 70\sqrt{5}\, \mathrm{cm}^{2}\)}{\(240\sqrt{5}\, \mathrm{cm}^{2}\)}{}}
\LANGsk{
\Question{
V kvádri \(ABCDEFGH\) 
platí: \(|AB| = 6\, \mathrm{cm};\ |AC| = 10\, \mathrm{cm};\ |AG| = 15\, \mathrm{cm}\). Povrch tohoto kvádra je:}
{\(96 + 140\sqrt{5}\, \mathrm{cm}^{2}\)}{\(600\, \mathrm{cm}^{2}\)}{\(236\sqrt{5}\, \mathrm{cm}^{2}\)}{\(48 + 70\sqrt{5}\, \mathrm{cm}^{2}\)}{\(240\sqrt{5}\, \mathrm{cm}^{2}\)}{}}
\LANGes{
\Question{
Un prisma rectangular  \(ABCDEFGH\) 
satisface \(|AB| = 6\, \mathrm{cm};\ |AC| = 10\, \mathrm{cm};\ |AG| = 15\, \mathrm{cm}\). La superficie de este prisma es:}
{\(96 + 140\sqrt{5}\, \mathrm{cm}^{2}\)}{\(600\, \mathrm{cm}^{2}\)}{\(236\sqrt{5}\, \mathrm{cm}^{2}\)}{\(48 + 70\sqrt{5}\, \mathrm{cm}^{2}\)}{\(240\sqrt{5}\, \mathrm{cm}^{2}\)}{}}
\End

\Data\ID{9000120307}
\Updated{2021-12-15 01:12:41}
\Level{A}
\SubArea{Volume and surface area formulas}
\IMG{001203_07-figure0.pdf}
\LANGen{
\Question{
The lengths of a side, base diagonal and solid diagonal through the vertex
\(A\) in a rectangular
box \(ABCDEFGH\) 
are \(|AB| = 6\, \mathrm{cm}\),
\(|AC| = 10\, \mathrm{cm}\),
\(|AG| = 15\, \mathrm{cm}\). Find
the volume of the box.}
{\(240\sqrt{5}\, \mathrm{cm}^{3}\)}{\(900\, \mathrm{cm}^{3}\)}{\(300\sqrt{5}\, \mathrm{cm}^{3}\)}{\(600\sqrt{2}\, \mathrm{cm}^{3}\)}{\(240\sqrt{2}\, \mathrm{cm}^{3}\)}{}}
\LANGpl{
\Question{
Długość boku, przekątnej podstawy i przekątnej prostopadłościanu przechodzącej przez wierzchołek 
\(A\) w prostopadłościanie
 \(ABCDEFGH\) 
są równe  \(|AB| = 6\, \mathrm{cm}\),
\(|AC| = 10\, \mathrm{cm}\),
\(|AG| = 15\, \mathrm{cm}\). Oblicz objętość prostopadłościanu.}
{\(240\sqrt{5}\, \mathrm{cm}^{3}\)}{\(900\, \mathrm{cm}^{3}\)}{\(300\sqrt{5}\, \mathrm{cm}^{3}\)}{\(600\sqrt{2}\, \mathrm{cm}^{3}\)}{\(240\sqrt{2}\, \mathrm{cm}^{3}\)}{}}
\LANGcs{
\Question{
V kvádru \(ABCDEFGH\) 
platí: \(|AB| = 6\, \mathrm{cm};\ |AC| = 10\, \mathrm{cm};\ |AG| = 15\, \mathrm{cm}\). Objem tohoto kvádru je:}
{\(240\sqrt{5}\, \mathrm{cm}^{3}\)}{\(900\, \mathrm{cm}^{3}\)}{\(300\sqrt{5}\, \mathrm{cm}^{3}\)}{\(600\sqrt{2}\, \mathrm{cm}^{3}\)}{\(240\sqrt{2}\, \mathrm{cm}^{3}\)}{}}
\LANGsk{
\Question{
V kvádri \(ABCDEFGH\) 
platí: \(|AB| = 6\, \mathrm{cm};\ |AC| = 10\, \mathrm{cm};\ |AG| = 15\, \mathrm{cm}\). Objem tohoto kvádra je:}
{\(240\sqrt{5}\, \mathrm{cm}^{3}\)}{\(900\, \mathrm{cm}^{3}\)}{\(300\sqrt{5}\, \mathrm{cm}^{3}\)}{\(600\sqrt{2}\, \mathrm{cm}^{3}\)}{\(240\sqrt{2}\, \mathrm{cm}^{3}\)}{}}
\LANGes{
\Question{
Un prisma rectangular  \(ABCDEFGH\) 
satisface \(|AB| = 6\, \mathrm{cm};\ |AC| = 10\, \mathrm{cm};\ |AG| = 15\, \mathrm{cm}\). El volumen de este prisma es:}
{\(240\sqrt{5}\, \mathrm{cm}^{3}\)}{\(900\, \mathrm{cm}^{3}\)}{\(300\sqrt{5}\, \mathrm{cm}^{3}\)}{\(600\sqrt{2}\, \mathrm{cm}^{3}\)}{\(240\sqrt{2}\, \mathrm{cm}^{3}\)}{}}
\End

\Data\ID{9000120310}
\Updated{2021-01-21 08:01:19}
\Level{A}
\SubArea{Volume and surface area formulas}
\IMG{001203_10-figure0.pdf}
\LANGen{
\Question{
The base of a rectangular box \(ABCDEFGH\) 
has sides \(|AB| = 6\, \mathrm{cm}\) and
\(|BC| = 8\, \mathrm{cm}\). The angle between
the solid diagonal \(AG\) 
and the base \(ABC\) 
is \(60^{\circ }\).
Find the volume of the box.}
{\(480\sqrt{3}\, \mathrm{cm}^{3}\)}{\(960\, \mathrm{cm}^{3}\)}{\(288\sqrt{3}\, \mathrm{cm}^{3}\)}{\(160\sqrt{3}\, \mathrm{cm}^{3}\)}{\(240\, \mathrm{cm}^{3}\)}{}}
\LANGpl{
\Question{
Dany jest prostopadłościan o bokach \(ABCDEFGH\) 
 \(|AB| = 6\, \mathrm{cm}\) i
\(|BC| = 8\, \mathrm{cm}\). Kąt między
przekątną \(AG\) 
a podstawą \(ABC\) 
jest równy \(60^{\circ }\).
Oblicz objętość prostopadłościanu.}
{\(480\sqrt{3}\, \mathrm{cm}^{3}\)}{\(960\, \mathrm{cm}^{3}\)}{\(288\sqrt{3}\, \mathrm{cm}^{3}\)}{\(160\sqrt{3}\, \mathrm{cm}^{3}\)}{\(240\, \mathrm{cm}^{3}\)}{}}
\LANGcs{
\Question{
V kvádru \(ABCDEFGH\) 
(\(|AB| = 6\, \mathrm{cm}\),
\(|BC| = 8\, \mathrm{cm}\)) je odchylka
úhlopříčky \(AG\) 
od roviny \(ABC\) 
rovna \(60^{\circ }\). Objem tohoto tělesa je roven:}
{\(480\sqrt{3}\, \mathrm{cm}^{3}\)}{\(960\, \mathrm{cm}^{3}\)}{\(288\sqrt{3}\, \mathrm{cm}^{3}\)}{\(160\sqrt{3}\, \mathrm{cm}^{3}\)}{\(240\, \mathrm{cm}^{3}\)}{}}
\LANGsk{
\Question{
V kvádri \(ABCDEFGH\) 
(\(|AB| = 6\, \mathrm{cm}\),
\(|BC| = 8\, \mathrm{cm}\)) je odchýlka
uhlopriečky \(AG\) 
od roviny \(ABC\) 
rovná \(60^{\circ }\). Objem tohoto telesa je rovný:}
{\(480\sqrt{3}\, \mathrm{cm}^{3}\)}{\(960\, \mathrm{cm}^{3}\)}{\(288\sqrt{3}\, \mathrm{cm}^{3}\)}{\(160\sqrt{3}\, \mathrm{cm}^{3}\)}{\(240\, \mathrm{cm}^{3}\)}{}}
\LANGes{
\Question{
Sea un prisma rectangular \(ABCDEFGH\) 
(\(|AB| = 6\, \mathrm{cm}\),
\(|BC| = 8\, \mathrm{cm}\)) El ángulo entre la diagonal \(AG\) 
y el plano \(ABC\) 
es \(60^{\circ }\). El volumen de este prisma es:}
{\(480\sqrt{3}\, \mathrm{cm}^{3}\)}{\(960\, \mathrm{cm}^{3}\)}{\(288\sqrt{3}\, \mathrm{cm}^{3}\)}{\(160\sqrt{3}\, \mathrm{cm}^{3}\)}{\(240\, \mathrm{cm}^{3}\)}{}}
\End

\Data\ID{1003165902}
\Updated{2022-01-06 08:01:13}
\Level{B}
\SubArea{Volume and surface area formulas}
\LANGen{
\Question{
Find the capacity of a garden pool in the shape of a cylinder with the diameter of \( 366\,\mathrm{cm} \) and the height of \( 0.91\,\mathrm{m} \). Round your result to \( 2 \) decimal places.}
{\( 9.57\,\mathrm{m}^3 \)}{\( 38.30\,\mathrm{m}^3 \)}{\( 957.74\,\mathrm{m}^3 \)}{\( 19.15\,\mathrm{m}^3 \)}{}{}}
\LANGpl{
\Question{
Oblicz pojemność basenu ogrodowego w kształcie walca o średnicy \( 366\,\mathrm{cm} \) i wysokości  \( 0{,}91\,\mathrm{m} \). Zaokrągli wynik do \( 2 \) miejsc po przecinku.}
{\( 9{,}57\,\mathrm{m}^3 \)}{\( 38{,}30\,\mathrm{m}^3 \)}{\( 957{,}74\,\mathrm{m}^3 \)}{\( 19{,}15\,\mathrm{m}^3 \)}{}{}}
\LANGcs{
\Question{
Vypočítejte objem zahradního bazénu tvaru válce s průměrem dna  \( 366\,\mathrm{cm} \) a výškou \( 0{,}91\,\mathrm{m} \). Výsledek uveďte s přesností na \( 2 \) desetinná místa.}
{\( 9{,}57\,\mathrm{m}^3 \)}{\( 38{,}30\,\mathrm{m}^3 \)}{\( 957{,}74\,\mathrm{m}^3 \)}{\( 19{,}15\,\mathrm{m}^3 \)}{}{}}
\LANGsk{
\Question{
Vypočítajte objem záhradného bazénu tvaru valca s priemerom dna  \( 366\,\mathrm{cm} \) a výškou \( 0{,}91\,\mathrm{m} \). Výsledok uveďte s presnosťou na \( 2 \) desatinné miesta.}
{\( 9{,}57\,\mathrm{m}^3 \)}{\( 38{,}30\,\mathrm{m}^3 \)}{\( 957{,}74\,\mathrm{m}^3 \)}{\( 19{,}15\,\mathrm{m}^3 \)}{}{}}
\LANGes{
\Question{
Calcula el volumen de una piscina cilíndrica cuyo radio de la base es  \( 366\,\mathrm{cm} \) y cuya altura es \( 0{,}91\,\mathrm{m} \). Para el resultado usa \( 2 \) cifras decimales.}
{\( 9{,}57\,\mathrm{m}^3 \)}{\( 38{,}30\,\mathrm{m}^3 \)}{\( 957{,}74\,\mathrm{m}^3 \)}{\( 19{,}15\,\mathrm{m}^3 \)}{}{}}
\End

\Data\ID{1003165903}
\Updated{2022-04-23 05:04:57}
\Level{B}
\SubArea{Volume and surface area formulas}
\LANGen{
\Question{
The volume of a cylinder with diameter of \( 20\,\mathrm{cm} \) is \( 5\,\mathrm{l} \). Find the height of the cylinder. Round your result to \( 2 \) decimal places.}
{\( 15.92\,\mathrm{cm} \)}{\( 3.98\,\mathrm{cm} \)}{\( 79.58\,\mathrm{cm} \)}{\( 159.92\,\mathrm{cm} \)}{}{}}
\LANGpl{
\Question{
Objętość walca o średnicy równej  \( 20\,\mathrm{cm} \) wynosi \( 5\,\mathrm{l} \). Oblicz wysokość walca. Zaokrągli wynik do  \( 2 \) miejsc po przecinku.}
{\( 15{,}92\,\mathrm{cm} \)}{\( 3{,}98\,\mathrm{cm} \)}{\( 79{,}58\,\mathrm{cm} \)}{\( 159{,}92\,\mathrm{cm} \)}{}{}}
\LANGcs{
\Question{
Určete výšku válce o objemu  \( 5\,\mathrm{l} \), jehož podstava má průměr \( 20\,\mathrm{cm} \). Výsledek uveďte s přesností na \( 2 \) desetinná místa.}
{\( 15{,}92\,\mathrm{cm} \)}{\( 3{,}98\,\mathrm{cm} \)}{\( 79{,}58\,\mathrm{cm} \)}{\( 159{,}92\,\mathrm{cm} \)}{}{}}
\LANGsk{
\Question{
Určte výšku valca s objemom  \( 5\,\mathrm{l} \), ktorého podstava má priemer \( 20\,\mathrm{cm} \). Výsledok uveďte s presnosťou na \( 2 \) desatinné miesta.}
{\( 15{,}92\,\mathrm{cm} \)}{\( 3{,}98\,\mathrm{cm} \)}{\( 79{,}58\,\mathrm{cm} \)}{\( 159{,}92\,\mathrm{cm} \)}{}{}}
\LANGes{
\Question{
Averigua la altura de un cilindro cuyo volumen es \( 5\,\mathrm{l} \), y cuya base tiene un radio de \( 20\,\mathrm{cm} \). Expresa el resultado con exactitud a \( 2 \) cifras decimales.}
{\( 15{,}92\,\mathrm{cm} \)}{\( 3{,}98\,\mathrm{cm} \)}{\( 79{,}58\,\mathrm{cm} \)}{\( 159{,}92\,\mathrm{cm} \)}{}{}}
\End

\Data\ID{1003165904}
\Updated{2022-01-06 08:01:09}
\Level{B}
\SubArea{Volume and surface area formulas}
\LANGen{
\Question{
How many litres of water can a cylinder-shaped plastic barrel with diameter of \( 30.48\,\mathrm{cm} \) and height of \( 51\,\mathrm{cm} \) hold? Round your result to \( 1 \) decimal place.}
{\( 37.2\,\mathrm{l} \)}{\( 148.9\,\mathrm{l} \)}{\( 372.1\,\mathrm{l} \)}{\( 62.3\,\mathrm{l} \)}{}{}}
\LANGpl{
\Question{
Ile litrów wody zmieści się w beczce w kształcie walca o średnicy równej \( 30{,}48\,\mathrm{cm} \) i wysokości  \( 51\,\mathrm{cm} \)? Zaokrągli wynik do  \( 1 \) miejsca po przecinku.}
{\( 37{,}2\,\mathrm{l} \)}{\( 148{,}9\,\mathrm{l} \)}{\( 372{,}1\,\mathrm{l} \)}{\( 62{,}3\,\mathrm{l} \)}{}{}}
\LANGcs{
\Question{
Kolik litrů vody se vejde do plastového barelu na vodu tvaru válce s průměrem dna \( 30{,}48\,\mathrm{cm} \) a výškou \( 51\,\mathrm{cm} \)? Výsledek uveďte s přesností na \( 1 \) desetinné místo.}
{\( 37{,}2\,\mathrm{l} \)}{\( 148{,}9\,\mathrm{l} \)}{\( 372{,}1\,\mathrm{l} \)}{\( 62{,}3\,\mathrm{l} \)}{}{}}
\LANGsk{
\Question{
Koľko litrov vody sa vojde do plastového barelu na vodu tvaru valca s priemerom dna \( 30{,}48\,\mathrm{cm} \) a výškou \( 51\,\mathrm{cm} \)? Výsledok uveďte s presnosťou na \( 1 \) desatinné miesto.}
{\( 37{,}2\,\mathrm{l} \)}{\( 148{,}9\,\mathrm{l} \)}{\( 372{,}1\,\mathrm{l} \)}{\( 62{,}3\,\mathrm{l} \)}{}{}}
\LANGes{
\Question{
¿Cuántos litros de agua caben en un barril cilínrico de diámetro \( 30{,}48\,\mathrm{cm} \) y  altura \( 51\,\mathrm{cm} \)? Expresa el resultado cn exactitud a \( 1 \) cifra decimal.}
{\( 37{,}2\,\mathrm{l} \)}{\( 148{,}9\,\mathrm{l} \)}{\( 372{,}1\,\mathrm{l} \)}{\( 62{,}3\,\mathrm{l} \)}{}{}}
\End

\Data\ID{1003170501}
\Updated{2021-01-20 01:01:49}
\Level{B}
\SubArea{Volume and surface area formulas}
\LANGen{
\Question{
Find the volume and the surface area of a sphere with radius of \( 6\,\mathrm{cm} \). Leave your answer in terms of \( \pi \).}
{\( V=288\pi\,\mathrm{cm}^3 \), \( S=144\pi\,\mathrm{cm}^2 \)}{\( V=144\pi\,\mathrm{cm}^3 \), \( S=288\pi\,\mathrm{cm}^2 \)}{\( V=1728\pi\,\mathrm{cm}^3 \), \( S=144\pi\,\mathrm{cm}^2 \)}{\( V=36\pi\,\mathrm{cm}^3 \), \( S=36\pi\,\mathrm{cm}^2 \)}{}{}}
\LANGpl{
\Question{
Oblicz objętość i pole powierzchni kuli o promieniu \( 6\,\mathrm{cm} \). Podaj wynik jako wielokrotność \( \pi \).}
{\( V=288\pi\,\mathrm{cm}^3 \), \( S=144\pi\,\mathrm{cm}^2 \)}{\( V=144\pi\,\mathrm{cm}^3 \), \( S=288\pi\,\mathrm{cm}^2 \)}{\( V=1728\pi\,\mathrm{cm}^3 \), \( S=144\pi\,\mathrm{cm}^2 \)}{\( V=36\pi\,\mathrm{cm}^3 \), \( S=36\pi\,\mathrm{cm}^2 \)}{}{}}
\LANGcs{
\Question{
Určete objem a povrch koule o poloměru \( 6\,\mathrm{cm} \). Výsledek vyjádřete jako násobek \( \pi \).}
{\( V=288\pi\,\mathrm{cm}^3 \), \( S=144\pi\,\mathrm{cm}^2 \)}{\( V=144\pi\,\mathrm{cm}^3 \), \( S=288\pi\,\mathrm{cm}^2 \)}{\( V=1728\pi\,\mathrm{cm}^3 \), \( S=144\pi\,\mathrm{cm}^2 \)}{\( V=36\pi\,\mathrm{cm}^3 \), \( S=36\pi\,\mathrm{cm}^2 \)}{}{}}
\LANGsk{
\Question{
Určte objem a povrch gule s polomerom \( 6\,\mathrm{cm} \). Výsledok vyjadrite ako násobok \( \pi \).}
{\( V=288\pi\,\mathrm{cm}^3 \), \( S=144\pi\,\mathrm{cm}^2 \)}{\( V=144\pi\,\mathrm{cm}^3 \), \( S=288\pi\,\mathrm{cm}^2 \)}{\( V=1728\pi\,\mathrm{cm}^3 \), \( S=144\pi\,\mathrm{cm}^2 \)}{\( V=36\pi\,\mathrm{cm}^3 \), \( S=36\pi\,\mathrm{cm}^2 \)}{}{}}
\LANGes{
\Question{
Calcula el volumen y la superficie de una esfera de radio de \( 6\,\mathrm{cm} \). Expresa el resultado como multiplicación de \( \pi \).}
{\( V=288\pi\,\mathrm{cm}^3 \), \( S=144\pi\,\mathrm{cm}^2 \)}{\( V=144\pi\,\mathrm{cm}^3 \), \( S=288\pi\,\mathrm{cm}^2 \)}{\( V=1728\pi\,\mathrm{cm}^3 \), \( S=144\pi\,\mathrm{cm}^2 \)}{\( V=36\pi\,\mathrm{cm}^3 \), \( S=36\pi\,\mathrm{cm}^2 \)}{}{}}
\End

\Data\ID{1003170502}
\Updated{2022-04-23 05:04:57}
\Level{B}
\SubArea{Volume and surface area formulas}
\LANGen{
\Question{
A bowl in the shape of a hemisphere has diameter of \( 21\,\mathrm{cm} \). Find the volume (in liters) of water it can  hold.  Round your result  to  \( 1 \) decimal place.}
{\( 2.4\,\mathrm{l} \)}{\( 4.8\,\mathrm{l} \)}{\( 19.4\,\mathrm{l} \)}{\( 38.8\,\mathrm{l} \)}{}{}}
\LANGpl{
\Question{
Miska w kształcie półkuli ma średnicę równą  \( 21\,\mathrm{cm} \). Oblicz objętość (w litrach) wody jaką może pomieścić miska.  Zaokrągli wynik do  \( 1 \) miejsca po przecinku.}
{\( 2{,}4\,\mathrm{l} \)}{\( 4{,}8\,\mathrm{l} \)}{\( 19{,}4\,\mathrm{l} \)}{\( 38{,}8\,\mathrm{l} \)}{}{}}
\LANGcs{
\Question{
Určete objem vody v polokouli o průměru \( 21\,\mathrm{cm} \). Výsledek vyjádřete v litrech s přesností na \( 1 \) desetinné místo.}
{\( 2{,}4\,\mathrm{l} \)}{\( 4{,}8\,\mathrm{l} \)}{\( 19{,}4\,\mathrm{l} \)}{\( 38{,}8\,\mathrm{l} \)}{}{}}
\LANGsk{
\Question{
Určte objem vody v pologuli s priemerom \( 21\,\mathrm{cm} \). Výsledok vyjadrite v litroch s presnosťou na \( 1 \) desatinné miesto.}
{\( 2{,}4\,\mathrm{l} \)}{\( 4{,}8\,\mathrm{l} \)}{\( 19{,}4\,\mathrm{l} \)}{\( 38{,}8\,\mathrm{l} \)}{}{}}
\LANGes{
\Question{
Calcula el volumen del aguade una semiesfera de diámetro \( 21\,\mathrm{cm} \). Expresa el resultado en litros con exactitud de \( 1 \) cifra decimal.}
{\( 2{,}4\,\mathrm{l} \)}{\( 4{,}8\,\mathrm{l} \)}{\( 19{,}4\,\mathrm{l} \)}{\( 38{,}8\,\mathrm{l} \)}{}{}}
\End

\Data\ID{1003170503}
\Updated{2022-01-06 08:01:04}
\Level{B}
\SubArea{Volume and surface area formulas}
\LANGen{
\Question{
Find the volume (in liters) and the surface area (in \( \mathrm{dm}^2 \)) of a beach ball with diameter of \( 200\,\mathrm{mm} \). Round your result  to \( 1 \) decimal place.}
{\( V=4.2\,\mathrm{l} \), \( S=12.6\,\mathrm{dm}^2 \)}{\( V=42\,\mathrm{l} \), \( S=1.3\,\mathrm{dm}^2 \)}{\( V=33.5\,\mathrm{l} \), \( S=12.6\,\mathrm{dm}^2 \)}{\( V=4.2\,\mathrm{l} \), \( S=50.3\,\mathrm{dm}^2 \)}{}{}}
\LANGpl{
\Question{
Oblicz objętość (w litrach) i pole powierzchni (w \( \mathrm{dm}^2 \)) piłki plażowej o średnicy  \( 200\,\mathrm{mm} \). Zaokrągli wynik do  \( 1 \) miejsca po przecinku.}
{\( V=4{,}2\,\mathrm{l} \), \( S=12{,}6\,\mathrm{dm}^2 \)}{\( V=42\,\mathrm{l} \), \( S=1{,}3\,\mathrm{dm}^2 \)}{\( V=33{,}5\,\mathrm{l} \), \( S=12{,}6\,\mathrm{dm}^2 \)}{\( V=4{,}2\,\mathrm{l} \), \( S=50{,}3\,\mathrm{dm}^2 \)}{}{}}
\LANGcs{
\Question{
Určete objem a povrch volejbalového balónu o průměru \( 200\,\mathrm{mm} \). Výsledek objemu v litrech a povrchu v \( \mathrm{dm}^2 \) uveďte s přesností na \( 1 \) desetinné místo.}
{\( V=4{,}2\,\mathrm{l} \), \( S=12{,}6\,\mathrm{dm}^2 \)}{\( V=42\,\mathrm{l} \), \( S=1{,}3\,\mathrm{dm}^2 \)}{\( V=33{,}5\,\mathrm{l} \), \( S=12{,}6\,\mathrm{dm}^2 \)}{\( V=4{,}2\,\mathrm{l} \), \( S=50{,}3\,\mathrm{dm}^2 \)}{}{}}
\LANGsk{
\Question{
Určte objem a povrch volejbalovej lopty s priemerom \( 200\,\mathrm{mm} \). Výsledok objemu v litroch a povrchu v \( \mathrm{dm}^2 \) uveďte s presnosťou na \( 1 \) desatinné miesto.}
{\( V=4{,}2\,\mathrm{l} \), \( S=12{,}6\,\mathrm{dm}^2 \)}{\( V=42\,\mathrm{l} \), \( S=1{,}3\,\mathrm{dm}^2 \)}{\( V=33{,}5\,\mathrm{l} \), \( S=12{,}6\,\mathrm{dm}^2 \)}{\( V=4{,}2\,\mathrm{l} \), \( S=50{,}3\,\mathrm{dm}^2 \)}{}{}}
\LANGes{
\Question{
Averigua el volumen y la superficie de una pelota de volleyball cuyo radio es de \( 200\,\mathrm{mm} \). Expresa el resultado del volumen en litros y de la superficie en\( \mathrm{dm}^2 \) y con una exactitud de \( 1 \) cifra decimal}
{\( V=4{,}2\,\mathrm{l} \), \( S=12{,}6\,\mathrm{dm}^2 \)}{\( V=42\,\mathrm{l} \), \( S=1{,}3\,\mathrm{dm}^2 \)}{\( V=33{,}5\,\mathrm{l} \), \( S=12{,}6\,\mathrm{dm}^2 \)}{\( V=4{,}2\,\mathrm{l} \), \( S=50{,}3\,\mathrm{dm}^2 \)}{}{}}
\End

\Data\ID{1003170504}
\Updated{2022-01-06 08:01:34}
\Level{B}
\SubArea{Volume and surface area formulas}
\LANGen{
\Question{
Find the volume of a sphere with surface area of \( 200\,\mathrm{cm}^2 \). Round your result  to \( 1 \) decimal place.}
{\( 266\,\mathrm{cm}^3 \)}{\( 265.9\,\mathrm{cm}^3 \)}{\( 16886.9\,\mathrm{cm}^3 \)}{\( 797.9\,\mathrm{cm}^3 \)}{}{}}
\LANGpl{
\Question{
Oblicz objętość kuli, której pole powierzchni jest równe \( 200\,\mathrm{cm}^2 \). Zaokrągli wynik do \( 1 \) miejsca po przecinku.}
{\( 266\,\mathrm{cm}^3 \)}{\( 265{,}9\,\mathrm{cm}^3 \)}{\( 16886{,}9\,\mathrm{cm}^3 \)}{\( 797{,}9\,\mathrm{cm}^3 \)}{}{}}
\LANGcs{
\Question{
Určete objem koule, jejíž povrch je \( 200\,\mathrm{cm}^2 \). Výsledek uveďte s přesností na \( 1 \) desetinné místo.}
{\( 266\,\mathrm{cm}^3 \)}{\( 265{,}9\,\mathrm{cm}^3 \)}{\( 16886{,}9\,\mathrm{cm}^3 \)}{\( 797{,}9\,\mathrm{cm}^3 \)}{}{}}
\LANGsk{
\Question{
Určte objem gule, ktorej povrch je \( 200\,\mathrm{cm}^2 \). Výsledok uveďte s presnosťou na \( 1 \) desatinné miesto.}
{\( 266\,\mathrm{cm}^3 \)}{\( 265{,}9\,\mathrm{cm}^3 \)}{\( 16886{,}9\,\mathrm{cm}^3 \)}{\( 797{,}9\,\mathrm{cm}^3 \)}{}{}}
\LANGes{
\Question{
Averigua el volumen de una esfera cuya superficie es \( 200\,\mathrm{cm}^2 \). Expresa el resultado con una exactitud de \( 1 \) cifra decimal.}
{\( 266\,\mathrm{cm}^3 \)}{\( 265{,}9\,\mathrm{cm}^3 \)}{\( 16886{,}9\,\mathrm{cm}^3 \)}{\( 797{,}9\,\mathrm{cm}^3 \)}{}{}}
\End

\Data\ID{1003170505}
\Updated{2022-01-06 08:01:25}
\Level{B}
\SubArea{Volume and surface area formulas}
\LANGen{
\Question{
The total surface area of a hemisphere is \( 154\,\mathrm{cm}^2 \). Find the radius of this hemisphere. Round your result  to \( 1 \) decimal place.}
{\( 4.0\,\mathrm{cm} \)}{\( 5.0\,\mathrm{cm} \)}{\( 16.3\,\mathrm{cm} \)}{\( 4.2\,\mathrm{cm} \)}{}{}}
\LANGpl{
\Question{
Całkowite pole powierzchni półkuli jest równe \( 154\,\mathrm{cm}^2 \). Oblicz promień półkuli. Zaokrągli wynik do \( 1 \) miejsca po przecinku.}
{\( 4{,}0\,\mathrm{cm} \)}{\( 5{,}0\,\mathrm{cm} \)}{\( 16{,}3\,\mathrm{cm} \)}{\( 4{,}2\,\mathrm{cm} \)}{}{}}
\LANGcs{
\Question{
Povrch polokoule je \( 154\,\mathrm{cm}^2 \). Určete její poloměr. Výsledek uveďte s přesností na \( 1 \) desetinné místo.}
{\( 4{,}0\,\mathrm{cm} \)}{\( 5{,}0\,\mathrm{cm} \)}{\( 16{,}3\,\mathrm{cm} \)}{\( 4{,}2\,\mathrm{cm} \)}{}{}}
\LANGsk{
\Question{
Povrch pologule je \( 154\,\mathrm{cm}^2 \). Určte jej polomer. Výsledok uveďte s presnosťou na \( 1 \) desatinné miesto.}
{\( 4{,}0\,\mathrm{cm} \)}{\( 5{,}0\,\mathrm{cm} \)}{\( 16{,}3\,\mathrm{cm} \)}{\( 4{,}2\,\mathrm{cm} \)}{}{}}
\LANGes{
\Question{
La superficie de una semiesfera es de \( 154\,\mathrm{cm}^2 \). Averigua su radio. Expresa el resultado con exactitud de \( 1 \) cifra decimal.}
{\( 4{,}0\,\mathrm{cm} \)}{\( 5{,}0\,\mathrm{cm} \)}{\( 16{,}3\,\mathrm{cm} \)}{\( 4{,}2\,\mathrm{cm} \)}{}{}}
\End

\Data\ID{1003170506}
\Updated{2021-01-20 02:01:18}
\Level{B}
\SubArea{Volume and surface area formulas}
\LANGen{
\Question{
Find the total surface area of a hemisphere with radius of \( 0.8\,\mathrm{m} \). Leave your result in terms of \( \pi \).}
{\( S = 1.28\pi\,\mathrm{m}^2 \)}{\( S = 1.92\pi\,\mathrm{m}^2 \)}{\( S = 1.54\pi\,\mathrm{m}^2 \)}{\( S = 1.8\pi\,\mathrm{m}^2 \)}{}{}}
\LANGpl{
\Question{
Oblicz całkowite pole powierzchni półkuli o promieniu \( 0{,}8\,\mathrm{m} \). Podaj wynik jako wielokrotność  \( \pi \).}
{\( S = 1{,}28\pi\,\mathrm{m}^2 \)}{\( S = 1{,}92\pi\,\mathrm{m}^2 \)}{\( S = 1{,}54\pi\,\mathrm{m}^2 \)}{\( S = 1{,}8\pi\,\mathrm{m}^2 \)}{}{}}
\LANGcs{
\Question{
Určete povrch polokoule s poloměrem \( 0{,}8\,\mathrm{m} \). Výsledek vyjádřete jako násobek \( \pi \).}
{\( S = 1{,}28\pi\,\mathrm{m}^2 \)}{\( S = 1{,}92\pi\,\mathrm{m}^2 \)}{\( S = 1{,}54\pi\,\mathrm{m}^2 \)}{\( S = 1{,}8\pi\,\mathrm{m}^2 \)}{}{}}
\LANGsk{
\Question{
Určte povrch pologule s polomerom \( 0{,}8\,\mathrm{m} \). Výsledok vyjadrite ako násobok \( \pi \).}
{\( S = 1{,}28\pi\,\mathrm{m}^2 \)}{\( S = 1{,}92\pi\,\mathrm{m}^2 \)}{\( S = 1{,}54\pi\,\mathrm{m}^2 \)}{\( S = 1{,}8\pi\,\mathrm{m}^2 \)}{}{}}
\LANGes{
\Question{
Averigua la superficie de una semisfera con radio de \( 0{,}8\,\mathrm{m} \). Expresa el resultado como múltiplo de \( \pi \).}
{\( S = 1{,}28\pi\,\mathrm{m}^2 \)}{\( S = 1{,}92\pi\,\mathrm{m}^2 \)}{\( S = 1{,}54\pi\,\mathrm{m}^2 \)}{\( S = 1{,}8\pi\,\mathrm{m}^2 \)}{}{}}
\End

\Data\ID{1003170706}
\Updated{2022-01-06 08:01:00}
\Level{B}
\SubArea{Volume and surface area formulas}
\LANGen{
\Question{
A snack bar sells popcorn in cone-shaped containers. One such container has the diameter of \( 20.32\,\mathrm{cm} \) and the height of \( 25.4\,\mathrm{cm} \). Find the volume of this container and leave your answer in litres.}
{\( 2.75\,\mathrm{l} \)}{\( 8.24\,\mathrm{l} \)}{\( 10.98\,\mathrm{l} \)}{\( 0.54\,\mathrm{l} \)}{}{}}
\LANGpl{
\Question{
Na stoisku można kupić popcorn w pojemnikach w kształcie stożka.  Średnica pojemnika jest równa \( 20{,}32\,\mathrm{cm} \), a wysokość \( 25{,}4\,\mathrm{cm} \).  Oblicz objętość pojemnika. Podaj odpowiedź w litrach.}
{\( 2{,}75\,\mathrm{l} \)}{\( 8{,}24\,\mathrm{l} \)}{\( 10{,}98\,\mathrm{l} \)}{\( 0{,}54\,\mathrm{l} \)}{}{}}
\LANGcs{
\Question{
Ve snack baru dávají popcorn do kornoutu tvaru rotačního kužele. Průměr podstavy je \( 20{,}32\,\mathrm{cm} \) a výška \( 25{,}4\,\mathrm{cm} \). Určete v litrech jeho objem.}
{\( 2{,}75\,\mathrm{l} \)}{\( 8{,}24\,\mathrm{l} \)}{\( 10{,}98\,\mathrm{l} \)}{\( 0{,}54\,\mathrm{l} \)}{}{}}
\LANGsk{
\Question{
V snack bare dávajú popcorn do kornútu tvaru rotačného kužele. Priemer podstavy je \( 20{,}32\,\mathrm{cm} \) a výška \( 25{,}4\,\mathrm{cm} \). Určte v litroch jeho objem.}
{\( 2{,}75\,\mathrm{l} \)}{\( 8{,}24\,\mathrm{l} \)}{\( 10{,}98\,\mathrm{l} \)}{\( 0{,}54\,\mathrm{l} \)}{}{}}
\LANGes{
\Question{
En un snack bar venden palomitas en un recipiente con forma de cono. El diámetro de la base es de \( 20{,}32\,\mathrm{cm} \) y la altura es de \( 25{,}4\,\mathrm{cm} \). Averigua su volumen en litros,}
{\( 2{,}75\,\mathrm{l} \)}{\( 8{,}24\,\mathrm{l} \)}{\( 10{,}98\,\mathrm{l} \)}{\( 0{,}54\,\mathrm{l} \)}{}{}}
\End

\Data\ID{1103164601}
\Updated{2022-01-06 08:01:53}
\Level{B}
\SubArea{Volume and surface area formulas}
\IMG{1103164601.pdf}
\LANGen{
\Question{
Let there be a triangular prism with the base area of \( 8\,\mathrm{cm}^2 \) and the height of \( 10\,\mathrm{cm} \) (see the picture). The volume of the prism is:}
{\( 80\,\mathrm{cm}^3 \)}{\( 40\,\mathrm{cm}^3 \)}{\( \frac{80}3\,\mathrm{cm}^3 \)}{\( 20\,\mathrm{cm}^3 \)}{}{}}
\LANGpl{
\Question{
Dany jest graniastosłup trójkątny, pole powierzchni jego podstawy jest równe  \( 8\,\mathrm{cm}^2 \) a wysokość wynosi \( 10\,\mathrm{cm} \) (jak na rysunku). Objętość graniastosłupa wynosi:}
{\( 80\,\mathrm{cm}^3 \)}{\( 40\,\mathrm{cm}^3 \)}{\( \frac{80}3\,\mathrm{cm}^3 \)}{\( 20\,\mathrm{cm}^3 \)}{}{}}
\LANGcs{
\Question{
Trojboký hranol má obsah podstavy  \( 8\,\mathrm{cm}^2 \) a výšku \( 10\,\mathrm{cm} \) (viz obrázek). Objem hranolu je:}
{\( 80\,\mathrm{cm}^3 \)}{\( 40\,\mathrm{cm}^3 \)}{\( \frac{80}3\,\mathrm{cm}^3 \)}{\( 20\,\mathrm{cm}^3 \)}{}{}}
\LANGsk{
\Question{
Trojboký hranol má obsah podstavy  \( 8\,\mathrm{cm}^2 \) a výšku \( 10\,\mathrm{cm} \) (viď obrázok). Objem hranola je:}
{\( 80\,\mathrm{cm}^3 \)}{\( 40\,\mathrm{cm}^3 \)}{\( \frac{80}3\,\mathrm{cm}^3 \)}{\( 20\,\mathrm{cm}^3 \)}{}{}}
\LANGes{
\Question{
Sea un prisma triangular cuya base tiene una superficie de  \( 8\,\mathrm{cm}^2 \) y cuya altura es de \( 10\,\mathrm{cm} \) (mira el dibujo). El volumen del prisma es:}
{\( 80\,\mathrm{cm}^3 \)}{\( 40\,\mathrm{cm}^3 \)}{\( \frac{80}3\,\mathrm{cm}^3 \)}{\( 20\,\mathrm{cm}^3 \)}{}{}}
\End

\Data\ID{1103164602}
\Updated{2022-01-07 08:01:03}
\Level{B}
\SubArea{Volume and surface area formulas}
\IMG{1103164602.pdf}
\LANGen{
\Question{
Let there be a triangular prism with the base area of \( 50\,\mathrm{cm}^2 \) and the volume of \( 400\,\mathrm{cm}^3 \) (see the picture). The height of the prism is:}
{\( 8\,\mathrm{cm} \)}{\( 16\,\mathrm{cm} \)}{\( 4\,\mathrm{cm} \)}{\( 24\,\mathrm{cm} \)}{}{}}
\LANGpl{
\Question{
Dany jest graniastosłup trójkątny, pole podstawy graniastosłupa jest równe \( 50\,\mathrm{cm}^2 \), objętość wynosi \( 400\,\mathrm{cm}^3 \) (spójrz na rysunek). Wysokość graniastosłupa wynosi:}
{\( 8\,\mathrm{cm} \)}{\( 16\,\mathrm{cm} \)}{\( 4\,\mathrm{cm} \)}{\( 24\,\mathrm{cm} \)}{}{}}
\LANGcs{
\Question{
Trojboký hranol má obsah podstavy  \( 50\,\mathrm{cm}^2 \) a objem  \( 400\,\mathrm{cm}^3 \) (viz obrázek). Výška hranolu je:}
{\( 8\,\mathrm{cm} \)}{\( 16\,\mathrm{cm} \)}{\( 4\,\mathrm{cm} \)}{\( 24\,\mathrm{cm} \)}{}{}}
\LANGsk{
\Question{
Trojboký hranol má obsah podstavy  \( 50\,\mathrm{cm}^2 \) a objem  \( 400\,\mathrm{cm}^3 \) (viď obrázok). Výška hranola je:}
{\( 8\,\mathrm{cm} \)}{\( 16\,\mathrm{cm} \)}{\( 4\,\mathrm{cm} \)}{\( 24\,\mathrm{cm} \)}{}{}}
\LANGes{
\Question{
Tenemos un prisma triangular. La superficie de su base es \( 50\,\mathrm{cm}^2 \) y su volumen es \( 400\,\mathrm{cm}^3 \) (observa el dibujo). La altura del prisma es:}
{\( 8\,\mathrm{cm} \)}{\( 16\,\mathrm{cm} \)}{\( 4\,\mathrm{cm} \)}{\( 24\,\mathrm{cm} \)}{}{}}
\End

\Data\ID{1103164603}
\Updated{2022-01-07 08:01:31}
\Level{B}
\SubArea{Volume and surface area formulas}
\IMG{1103164603_0.pdf}
\LANGen{
\Question{
Find the volume of the right triangular prism from the picture.}
{\( 60\,\mathrm{cm}^3 \)}{\( 120\,\mathrm{cm}^3 \)}{\( 40\,\mathrm{cm}^3 \)}{\( 80\,\mathrm{cm}^3 \)}{}{}}
\LANGpl{
\Question{
Oblicz objętość graniastosłupa trójkątnego prostokątnego zamieszczonego na rysunku.}
{\( 60\,\mathrm{cm}^3 \)}{\( 120\,\mathrm{cm}^3 \)}{\( 40\,\mathrm{cm}^3 \)}{\( 80\,\mathrm{cm}^3 \)}{}{}}
\LANGcs{
\Question{
Určete objem trojbokého hranolu s podstavou pravoúhlého trojúhelníku, jehož rozměry jsou zakresleny v obrázku.}
{\( 60\,\mathrm{cm}^3 \)}{\( 120\,\mathrm{cm}^3 \)}{\( 40\,\mathrm{cm}^3 \)}{\( 80\,\mathrm{cm}^3 \)}{}{}}
\LANGsk{
\Question{
Určte objem trojbokého hranola s podstavou pravouhlého trojuholníka, ktorého rozmery sú zakreslené v obrázku.}
{\( 60\,\mathrm{cm}^3 \)}{\( 120\,\mathrm{cm}^3 \)}{\( 40\,\mathrm{cm}^3 \)}{\( 80\,\mathrm{cm}^3 \)}{}{}}
\LANGes{
\Question{
Averigua el volumen del prisma triangular cuya base es un  triángulo rectángulo de medidas marcadas en el dibujo.}
{\( 60\,\mathrm{cm}^3 \)}{\( 120\,\mathrm{cm}^3 \)}{\( 40\,\mathrm{cm}^3 \)}{\( 80\,\mathrm{cm}^3 \)}{}{}}
\End

\Data\ID{1103164604}
\Updated{2021-01-20 02:01:36}
\Level{B}
\SubArea{Volume and surface area formulas}
\IMG{1103164604_0.pdf}
\LANGen{
\Question{
Find the surface area of the  right triangular prism from the picture.}
{\( 120\,\mathrm{cm}^2 \)}{\( 96\,\mathrm{cm}^2 \)}{\( 240\,\mathrm{cm}^2 \)}{\( 60\,\mathrm{cm}^2 \)}{}{}}
\LANGpl{
\Question{
Oblicz pole powierzchni graniastosłupa trójkątnego prostokątnego zamieszczonego na rysunku.}
{\( 120\,\mathrm{cm}^2 \)}{\( 96\,\mathrm{cm}^2 \)}{\( 240\,\mathrm{cm}^2 \)}{\( 60\,\mathrm{cm}^2 \)}{}{}}
\LANGcs{
\Question{
Povrch trojbokého hranolu, jehož rozměry jsou zakresleny v obrázku, je:}
{\( 120\,\mathrm{cm}^2 \)}{\( 96\,\mathrm{cm}^2 \)}{\( 240\,\mathrm{cm}^2 \)}{\( 60\,\mathrm{cm}^2 \)}{}{}}
\LANGsk{
\Question{
Povrch trojbokého hranola, ktorého rozmery sú zakreslené v obrázku, je:}
{\( 120\,\mathrm{cm}^2 \)}{\( 96\,\mathrm{cm}^2 \)}{\( 240\,\mathrm{cm}^2 \)}{\( 60\,\mathrm{cm}^2 \)}{}{}}
\LANGes{
\Question{
La superficie del prisma triangular, cuyas medidas están marcadas en el dibujo, es:}
{\( 120\,\mathrm{cm}^2 \)}{\( 96\,\mathrm{cm}^2 \)}{\( 240\,\mathrm{cm}^2 \)}{\( 60\,\mathrm{cm}^2 \)}{}{}}
\End

\Data\ID{1103164605}
\Updated{2022-01-07 08:01:52}
\Level{B}
\SubArea{Volume and surface area formulas}
\IMG{1103164605.pdf}
\LANGen{
\Question{
A triangular prism has a triangular base with the side \( a \) of \( 6\,\mathrm{dm} \) and with the height \( v_a \) of \( 4\,\mathrm{dm} \). The height \( h \) of the prism is \( 10\,\mathrm{dm} \) (see the picture). Find the volume of the prism.}
{\( 120\,\mathrm{dm}^3 \)}{\( 240\,\mathrm{dm}^3 \)}{\( 60\,\mathrm{dm}^3 \)}{\( 80\,\mathrm{dm}^3 \)}{}{}}
\LANGpl{
\Question{
Graniastosłup trójkątny o podstawie trójkąta o boku \( a \) równym \( 6\,\mathrm{dm} \) i wysokości  \( v_a \) równej \( 4\,\mathrm{dm} \). Wysokość  \( h \) graniastosłupa jest równa \( 10\,\mathrm{dm} \) (spójrz na rysunek). Oblicz objętość graniastosłupa.}
{\( 120\,\mathrm{dm}^3 \)}{\( 240\,\mathrm{dm}^3 \)}{\( 60\,\mathrm{dm}^3 \)}{\( 80\,\mathrm{dm}^3 \)}{}{}}
\LANGcs{
\Question{
Podstavu trojbokého hranolu tvoří trojúhelník se stranou \( a \) délky \( 6\,\mathrm{dm} \) a příslušnou výškou  \( v_a \) délky \( 4\,\mathrm{dm} \). Výška hranolu \( h \) má délku \( 10\,\mathrm{dm} \) (viz obrázek). Jeho objem je:}
{\( 120\,\mathrm{dm}^3 \)}{\( 240\,\mathrm{dm}^3 \)}{\( 60\,\mathrm{dm}^3 \)}{\( 80\,\mathrm{dm}^3 \)}{}{}}
\LANGsk{
\Question{
Podstavu trojbokého hranola tvorí trojuholník so stranou \( a \) dĺžky \( 6\,\mathrm{dm} \) a príslušnou výškou  \( v_a \) dĺžky \( 4\,\mathrm{dm} \). Výška hranola \( h \) má dĺžku \( 10\,\mathrm{dm} \) (viď obrázok). Jeho objem je:}
{\( 120\,\mathrm{dm}^3 \)}{\( 240\,\mathrm{dm}^3 \)}{\( 60\,\mathrm{dm}^3 \)}{\( 80\,\mathrm{dm}^3 \)}{}{}}
\LANGes{
\Question{
La base de un prisma es un triángulo con un lado \( a \) de longitud \( 6\,\mathrm{dm} \) y la altura \( v_a \) de longitud \( 4\,\mathrm{dm} \). La altura del prisma \( h \) es \( 10\,\mathrm{dm} \) (observa el dibujo). El volumen del prisma es:}
{\( 120\,\mathrm{dm}^3 \)}{\( 240\,\mathrm{dm}^3 \)}{\( 60\,\mathrm{dm}^3 \)}{\( 80\,\mathrm{dm}^3 \)}{}{}}
\End

\Data\ID{1103164606}
\Updated{2022-01-07 08:01:33}
\Level{B}
\SubArea{Volume and surface area formulas}
\IMG{1103164606.pdf}
\LANGen{
\Question{
Let there be a trapezoidal prism with the base area of \( 20\,\mathrm{cm}^2 \) and the volume of \( 60\,\mathrm{cm}^3 \) (see the picture). The height of the prism is:}
{\( 3\,\mathrm{cm} \)}{\( 6\,\mathrm{cm} \)}{\( 1.5\,\mathrm{cm} \)}{\( 9\,\mathrm{cm} \)}{}{}}
\LANGpl{
\Question{
Dany jest graniastosłup trapezowy, którego pole powierzchni wynosi  \( 20\,\mathrm{cm}^2 \), a objętość jest równa  \( 60\,\mathrm{cm}^3 \) (spójrz na rysunek). Wysokość graniastosłupa jest równa:}
{\( 3\,\mathrm{cm} \)}{\( 6\,\mathrm{cm} \)}{\( 1{,}5\,\mathrm{cm} \)}{\( 9\,\mathrm{cm} \)}{}{}}
\LANGcs{
\Question{
Čtyřboký hranol s podstavou lichoběžníku o obsahu \( 20\,\mathrm{cm}^2 \) má objem  \( 60\,\mathrm{cm}^3 \) (viz obrázek). Výška hranolu je:}
{\( 3\,\mathrm{cm} \)}{\( 6\,\mathrm{cm} \)}{\( 1{,}5\,\mathrm{cm} \)}{\( 9\,\mathrm{cm} \)}{}{}}
\LANGsk{
\Question{
Štvorboký hranol s podstavou lichobežníka s obsahom \( 20\,\mathrm{cm}^2 \) má objem  \( 60\,\mathrm{cm}^3 \) (viď obrázok). Výška hranola je:}
{\( 3\,\mathrm{cm} \)}{\( 6\,\mathrm{cm} \)}{\( 1{,}5\,\mathrm{cm} \)}{\( 9\,\mathrm{cm} \)}{}{}}
\LANGes{
\Question{
Un prisma cuya base es un trapecio de superficie \( 20\,\mathrm{cm}^2 \) tiene un volumen de \( 60\,\mathrm{cm}^3 \) (observa el dibujo). La altura del prisma es:}
{\( 3\,\mathrm{cm} \)}{\( 6\,\mathrm{cm} \)}{\( 1{,}5\,\mathrm{cm} \)}{\( 9\,\mathrm{cm} \)}{}{}}
\End

\Data\ID{1103164607}
\Updated{2022-01-07 08:01:32}
\Level{B}
\SubArea{Volume and surface area formulas}
\IMG{1103164607.pdf}
\LANGen{
\Question{
Find the volume of the  trapezoidal prism from the picture.}
{\( 700\,\mathrm{cm}^3 \)}{\( 1400\,\mathrm{cm}^3 \)}{\( 800\,\mathrm{cm}^3 \)}{\( 2400\,\mathrm{cm}^3 \)}{}{}}
\LANGpl{
\Question{
Oblicz objętość graniastosłupa trapezowego zamieszczonego na rysunku.}
{\( 700\,\mathrm{cm}^3 \)}{\( 1400\,\mathrm{cm}^3 \)}{\( 800\,\mathrm{cm}^3 \)}{\( 2400\,\mathrm{cm}^3 \)}{}{}}
\LANGcs{
\Question{
Určete objem čtyřbokého hranolu s podstavou lichoběžníka, jehož rozměry jsou zakresleny v obrázku.}
{\( 700\,\mathrm{cm}^3 \)}{\( 1400\,\mathrm{cm}^3 \)}{\( 800\,\mathrm{cm}^3 \)}{\( 2400\,\mathrm{cm}^3 \)}{}{}}
\LANGsk{
\Question{
Určte objem štvorbokého hranola s podstavou lichobežníka, ktorého rozmery sú zakreslené v obrázku.}
{\( 700\,\mathrm{cm}^3 \)}{\( 1400\,\mathrm{cm}^3 \)}{\( 800\,\mathrm{cm}^3 \)}{\( 2400\,\mathrm{cm}^3 \)}{}{}}
\LANGes{
\Question{
Dado un prisma con un trapecio como base cuyas medidas están marcadas en el dibujo. Averigua su volumen.}
{\( 700\,\mathrm{cm}^3 \)}{\( 1400\,\mathrm{cm}^3 \)}{\( 800\,\mathrm{cm}^3 \)}{\( 2400\,\mathrm{cm}^3 \)}{}{}}
\End

\Data\ID{1103164608}
\Updated{2022-04-23 05:04:57}
\Level{B}
\SubArea{Volume and surface area formulas}
\IMG{1103164608.pdf}
\LANGen{
\Question{
Find the surface area of the  trapezoidal prism from the picture.}
{\( 316\,\mathrm{cm}^2 \)}{\( 368\,\mathrm{cm}^2 \)}{\( 264\,\mathrm{cm}^2 \)}{\( 184\,\mathrm{cm}^2 \)}{}{}}
\LANGpl{
\Question{
Oblicz pole powierzchni graniastosłupa trapezowego zamieszczonego na rysunku.}
{\( 316\,\mathrm{cm}^2 \)}{\( 368\,\mathrm{cm}^2 \)}{\( 264\,\mathrm{cm}^2 \)}{\( 184\,\mathrm{cm}^2 \)}{}{}}
\LANGcs{
\Question{
Určete povrch hranolu s podstavou pravoúhlého lichoběžníku, jehož rozměry jsou zakreslené v obrázku.}
{\( 316\,\mathrm{cm}^2 \)}{\( 368\,\mathrm{cm}^2 \)}{\( 264\,\mathrm{cm}^2 \)}{\( 184\,\mathrm{cm}^2 \)}{}{}}
\LANGsk{
\Question{
Určte povrch hranola s podstavou pravouhlého lichobežníka, ktorého rozmery sú zakreslené v obrázku.}
{\( 316\,\mathrm{cm}^2 \)}{\( 368\,\mathrm{cm}^2 \)}{\( 264\,\mathrm{cm}^2 \)}{\( 184\,\mathrm{cm}^2 \)}{}{}}
\LANGes{
\Question{
Sea un prisma con un trapezóide rectángulo, cuyas medidas están marcadas en el dibujo. Averigua su volumen.}
{\( 316\,\mathrm{cm}^2 \)}{\( 368\,\mathrm{cm}^2 \)}{\( 264\,\mathrm{cm}^2 \)}{\( 184\,\mathrm{cm}^2 \)}{}{}}
\End

\Data\ID{1103164609}
\Updated{2021-12-15 01:12:07}
\Level{B}
\SubArea{Volume and surface area formulas}
\IMG{1103164609.pdf}
\LANGen{
\Question{
Find the volume of the  trapezoidal prism from the picture.}
{\( 4140\,\mathrm{cm}^3 \)}{\( 2070\,\mathrm{cm}^3 \)}{\( 3680\,\mathrm{cm}^3 \)}{\( 1840\,\mathrm{cm}^3 \)}{}{}}
\LANGpl{
\Question{
Oblicz objętość graniastosłupa trapezowego zamieszczonego na rysunku.}
{\( 4140\,\mathrm{cm}^3 \)}{\( 2070\,\mathrm{cm}^3 \)}{\( 3680\,\mathrm{cm}^3 \)}{\( 1840\,\mathrm{cm}^3 \)}{}{}}
\LANGcs{
\Question{
Určete objem čtyřbokého hranolu, jehož rozměry jsou zobrazeny na obrázku.}
{\( 4140\,\mathrm{cm}^3 \)}{\( 2070\,\mathrm{cm}^3 \)}{\( 3680\,\mathrm{cm}^3 \)}{\( 1840\,\mathrm{cm}^3 \)}{}{}}
\LANGsk{
\Question{
Určte objem štvorbokého hranola, ktorého rozmery sú zakreslené v obrázku.}
{\( 4140\,\mathrm{cm}^3 \)}{\( 2070\,\mathrm{cm}^3 \)}{\( 3680\,\mathrm{cm}^3 \)}{\( 1840\,\mathrm{cm}^3 \)}{}{}}
\LANGes{
\Question{
Averigua el volumen de un prisma cuatrilátero, cuyas medidas están marcadas en el dibujo.}
{\( 4140\,\mathrm{cm}^3 \)}{\( 2070\,\mathrm{cm}^3 \)}{\( 3680\,\mathrm{cm}^3 \)}{\( 1840\,\mathrm{cm}^3 \)}{}{}}
\End

\Data\ID{1103165901}
\Updated{2022-01-06 08:01:49}
\Level{B}
\SubArea{Volume and surface area formulas}
\IMG{1103165901.pdf}
\LANGen{
\Question{
Find the volume and the surface area of a cylinder with the radius \( 3\,\mathrm{cm} \) and the height \( 8\,\mathrm{cm} \) (see the picture). Give your result in terms of \( \pi \).}
{\( V=72\pi\,\mathrm{cm}^3 \), \( S=66\pi\,\mathrm{cm}^2 \)}{\( V=144\pi\,\mathrm{cm}^3 \), \( S=198\pi\,\mathrm{cm}^2 \)}{\( V=144\pi\,\mathrm{cm}^3 \), \( S=66\pi\,\mathrm{cm}^2 \)}{\( V=72\pi\,\mathrm{cm}^3 \), \( S=198\pi\,\mathrm{cm}^2 \)}{}{}}
\LANGpl{
\Question{
Oblicz objętość i  pole powierzchni walca o promieniu równym \( 3\,\mathrm{cm} \) i wysokości \( 8\,\mathrm{cm} \) (spójrz na rysunek). Podaj wynik jako wielokrotność \( \pi \).}
{\( V=72\pi\,\mathrm{cm}^3 \), \( S=66\pi\,\mathrm{cm}^2 \)}{\( V=144\pi\,\mathrm{cm}^3 \), \( S=198\pi\,\mathrm{cm}^2 \)}{\( V=144\pi\,\mathrm{cm}^3 \), \( S=66\pi\,\mathrm{cm}^2 \)}{\( V=72\pi\,\mathrm{cm}^3 \), \( S=198\pi\,\mathrm{cm}^2 \)}{}{}}
\LANGcs{
\Question{
Vypočítejte objem a povrch válce s poloměrem podstavy \( 3\,\mathrm{cm} \) a výškou \( 8\,\mathrm{cm} \) (viz obrázek). Výsledek uveďte jako násobek \( \pi \).}
{\( V=72\pi\,\mathrm{cm}^3 \), \( S=66\pi\,\mathrm{cm}^2 \)}{\( V=144\pi\,\mathrm{cm}^3 \), \( S=198\pi\,\mathrm{cm}^2 \)}{\( V=144\pi\,\mathrm{cm}^3 \), \( S=66\pi\,\mathrm{cm}^2 \)}{\( V=72\pi\,\mathrm{cm}^3 \), \( S=198\pi\,\mathrm{cm}^2 \)}{}{}}
\LANGsk{
\Question{
Vypočítajte objem a povrch valca s polomerom podstavy \( 3\,\mathrm{cm} \) a výškou \( 8\,\mathrm{cm} \) (viď obrázok). Výsledok uveďte ako násobok \( \pi \).}
{\( V=72\pi\,\mathrm{cm}^3 \), \( S=66\pi\,\mathrm{cm}^2 \)}{\( V=144\pi\,\mathrm{cm}^3 \), \( S=198\pi\,\mathrm{cm}^2 \)}{\( V=144\pi\,\mathrm{cm}^3 \), \( S=66\pi\,\mathrm{cm}^2 \)}{\( V=72\pi\,\mathrm{cm}^3 \), \( S=198\pi\,\mathrm{cm}^2 \)}{}{}}
\LANGes{
\Question{
Calcula el volumen y la superficie de un cílindro sabiendo que el radio de su base es \( 3\,\mathrm{cm} \) y su altura es \( 8\,\mathrm{cm} \) (observa el dibujo). Expresa el resutado como múltiplo de \( \pi \).}
{\( V=72\pi\,\mathrm{cm}^3 \), \( S=66\pi\,\mathrm{cm}^2 \)}{\( V=144\pi\,\mathrm{cm}^3 \), \( S=198\pi\,\mathrm{cm}^2 \)}{\( V=144\pi\,\mathrm{cm}^3 \), \( S=66\pi\,\mathrm{cm}^2 \)}{\( V=72\pi\,\mathrm{cm}^3 \), \( S=198\pi\,\mathrm{cm}^2 \)}{}{}}
\End

\Data\ID{1103165905}
\Updated{2022-11-12 10:11:46}
\Level{B}
\SubArea{Volume and surface area formulas}
\IMG{1103165905.pdf}
\LANGen{
\Question{
How much paper do we need to label the can of peas with diameter of \( 10\,\mathrm{cm} \) and height of \( 20\,\mathrm{cm} \)? (Label covers the side of the can completely, the bottom and the top base are not labelled.) Round your result to \( 1 \) decimal place.}
{\( 628.3\,\mathrm{cm}^2 \)}{\( 1256.6\,\mathrm{cm}^2 \)}{\( 314.2\,\mathrm{cm}^2 \)}{\( 785.4\,\mathrm{cm}^2 \)}{}{}}
\LANGpl{
\Question{
Ile potrzebujemy papieru, aby owinąć puszkę groszku o średnicy \( 10\,\mathrm{cm} \) i wysokości \( 20\,\mathrm{cm} \)? (Papier musi pokryć całkowicie pole powierzchni bocznej puszki, papier nie pokrywa górnej i dolnej podstawy puszki.) Zaokrągli wynik do \( 1 \) miejsca po przecinku.}
{\( 628{,}3\,\mathrm{cm}^2 \)}{\( 1256{,}6\,\mathrm{cm}^2 \)}{\( 314{,}2\,\mathrm{cm}^2 \)}{\( 785{,}4\,\mathrm{cm}^2 \)}{}{}}
\LANGcs{
\Question{
Kolik papíru potřebujeme k výrobě etikety na plášť konzervy s hráškem tvaru válce o průměru \( 10\,\mathrm{cm} \) a výšce \( 20\,\mathrm{cm} \)? (Etiketa není umístěna na podstavách válce.) Výsledek uveďte s přesností na \( 1 \) desetinné místo.}
{\( 628{,}3\,\mathrm{cm}^2 \)}{\( 1256{,}6\,\mathrm{cm}^2 \)}{\( 314{,}2\,\mathrm{cm}^2 \)}{\( 785{,}4\,\mathrm{cm}^2 \)}{}{}}
\LANGsk{
\Question{
Koľko papiera potrebujeme k výrobe etikety na obal konzervy s hráškom tvaru valca s priemerom \( 10\,\mathrm{cm} \) a výškou \( 20\,\mathrm{cm} \)? (Etiketa nie je umiestená na podstavách valca.) Výsledok uveďte s presnosťou na \( 1 \) desatinné miesto.}
{\( 628{,}3\,\mathrm{cm}^2 \)}{\( 1256{,}6\,\mathrm{cm}^2 \)}{\( 314{,}2\,\mathrm{cm}^2 \)}{\( 785{,}4\,\mathrm{cm}^2 \)}{}{}}
\LANGes{
\Question{
¿Qué cantidad de papel necesitamos para producir una etiqueta (área lateral) de una lata de guisantes cilíndrica, cuyo diámetro es \( 10\,\mathrm{cm} \) y cuya altura mide \( 20\,\mathrm{cm} \)? (La etiqueta no está en las bases del cilindro.) Expresa el resultado con exactitud de \( 1 \) cifra decimal.}
{\( 628{,}3\,\mathrm{cm}^2 \)}{\( 1256{,}6\,\mathrm{cm}^2 \)}{\( 314{,}2\,\mathrm{cm}^2 \)}{\( 785{,}4\,\mathrm{cm}^2 \)}{}{}}
\End

\Data\ID{1103165906}
\Updated{2022-01-07 08:01:04}
\Level{B}
\SubArea{Volume and surface area formulas}
\IMG{1103165906.pdf}
\LANGen{
\Question{
The volume of a cylinder with the height of \( 12\,\mathrm{cm} \) is \( 60\,\mathrm{cm}^3 \). Find the surface area of this cylinder. Round your result to \( 2 \) decimal places.}
{\( 105.12\,\mathrm{cm}^2 \)}{\( 52.56\,\mathrm{cm}^2 \)}{\( 135.54\,\mathrm{cm}^2 \)}{\( 210.24\,\mathrm{cm}^2 \)}{}{}}
\LANGpl{
\Question{
Objętość walca o wysokości równej  \( 12\,\mathrm{cm} \) wynosi \( 60\,\mathrm{cm}^3 \). Oblicz pole powierzchni walca. Zaokrągli wynik do \( 2 \) dwóch miejsc po przecinku.}
{\( 105{,}12\,\mathrm{cm}^2 \)}{\( 52{,}56\,\mathrm{cm}^2 \)}{\( 135{,}54\,\mathrm{cm}^2 \)}{\( 210{,}24\,\mathrm{cm}^2 \)}{}{}}
\LANGcs{
\Question{
Válec o výšce  \( 12\,\mathrm{cm} \) má objem \( 60\,\mathrm{cm}^3 \). Určete jeho povrch. Výsledek uveďte s přesností na \( 2 \) desetinná místa.}
{\( 105{,}12\,\mathrm{cm}^2 \)}{\( 52{,}56\,\mathrm{cm}^2 \)}{\( 135{,}54\,\mathrm{cm}^2 \)}{\( 210{,}24\,\mathrm{cm}^2 \)}{}{}}
\LANGsk{
\Question{
Valec s výškou  \( 12\,\mathrm{cm} \) má objem \( 60\,\mathrm{cm}^3 \). Určte jeho povrch. Výsledok uveďte s presnosťou na \( 2 \) desatinné miesta.}
{\( 105{,}12\,\mathrm{cm}^2 \)}{\( 52{,}56\,\mathrm{cm}^2 \)}{\( 135{,}54\,\mathrm{cm}^2 \)}{\( 210{,}24\,\mathrm{cm}^2 \)}{}{}}
\LANGes{
\Question{
Un cílindro de altura  \( 12\,\mathrm{cm} \) tiene volumen \( 60\,\mathrm{cm}^3 \). Averigua su superficie. Expresa el resultado con exactitud de \( 2 \) cifras decimales.}
{\( 105{,}12\,\mathrm{cm}^2 \)}{\( 52{,}56\,\mathrm{cm}^2 \)}{\( 135{,}54\,\mathrm{cm}^2 \)}{\( 210{,}24\,\mathrm{cm}^2 \)}{}{}}
\End

\Data\ID{1103170701}
\Updated{2022-04-23 05:04:57}
\Level{B}
\SubArea{Volume and surface area formulas}
\IMG{1103170701.pdf}
\LANGen{
\Question{
Let there be a cone with the base radius of \( 6\,\mathrm{cm} \) and the perpendicular height of \( 8\,\mathrm{cm} \). Find the volume and surface area of the cone. Leave your answer in terms of \( \pi \).}
{\( V=96\pi\,\mathrm{cm}^3 \), \( S=96\pi\,\mathrm{cm}^2 \)}{\( V=96\pi\,\mathrm{cm}^3 \), \( S=84\pi\,\mathrm{cm}^2 \)}{\( V=288\pi\,\mathrm{cm}^3 \), \( S=84\pi\,\mathrm{cm}^2 \)}{\( V=16\pi\,\mathrm{cm}^3 \), \( S=96\pi\,\mathrm{cm}^2 \)}{}{}}
\LANGpl{
\Question{
Dany jest stożek, promień jego podstawy jest równy \( 6\,\mathrm{cm} \) natomiast prostopadła do niej wysokość wynosi \( 8\,\mathrm{cm} \). Oblicz objętość i pole powierzchni stożka. Podaj wynik jako wielokrotność \( \pi \).}
{\( V=96\pi\,\mathrm{cm}^3 \), \( S=96\pi\,\mathrm{cm}^2 \)}{\( V=96\pi\,\mathrm{cm}^3 \), \( S=84\pi\,\mathrm{cm}^2 \)}{\( V=288\pi\,\mathrm{cm}^3 \), \( S=84\pi\,\mathrm{cm}^2 \)}{\( V=16\pi\,\mathrm{cm}^3 \), \( S=96\pi\,\mathrm{cm}^2 \)}{}{}}
\LANGcs{
\Question{
Určete objem a povrch rotačního kužele o výšce \( 8\,\mathrm{cm} \), jehož podstava má poloměr \( 6\,\mathrm{cm} \). Výsledek vyjádřete jako násobek \( \pi \).}
{\( V=96\pi\,\mathrm{cm}^3 \), \( S=96\pi\,\mathrm{cm}^2 \)}{\( V=96\pi\,\mathrm{cm}^3 \), \( S=84\pi\,\mathrm{cm}^2 \)}{\( V=288\pi\,\mathrm{cm}^3 \), \( S=84\pi\,\mathrm{cm}^2 \)}{\( V=16\pi\,\mathrm{cm}^3 \), \( S=96\pi\,\mathrm{cm}^2 \)}{}{}}
\LANGsk{
\Question{
Určte objem a povrch rotačného kužeľa s výškou \( 8\,\mathrm{cm} \), ktorého podstava má polomer \( 6\,\mathrm{cm} \). Výsledok vyjadrite ako násobok \( \pi \).}
{\( V=96\pi\,\mathrm{cm}^3 \), \( S=96\pi\,\mathrm{cm}^2 \)}{\( V=96\pi\,\mathrm{cm}^3 \), \( S=84\pi\,\mathrm{cm}^2 \)}{\( V=288\pi\,\mathrm{cm}^3 \), \( S=84\pi\,\mathrm{cm}^2 \)}{\( V=16\pi\,\mathrm{cm}^3 \), \( S=96\pi\,\mathrm{cm}^2 \)}{}{}}
\LANGes{
\Question{
Averigua el volumen y la superficie de un cono sabiendo que su altura es de \( 8\,\mathrm{cm} \) y el radio de su base es de \( 6\,\mathrm{cm} \). Expresa el resultado como múltiplo de \( \pi \).}
{\( V=96\pi\,\mathrm{cm}^3 \), \( S=96\pi\,\mathrm{cm}^2 \)}{\( V=96\pi\,\mathrm{cm}^3 \), \( S=84\pi\,\mathrm{cm}^2 \)}{\( V=288\pi\,\mathrm{cm}^3 \), \( S=84\pi\,\mathrm{cm}^2 \)}{\( V=16\pi\,\mathrm{cm}^3 \), \( S=96\pi\,\mathrm{cm}^2 \)}{}{}}
\End

\Data\ID{1103170702}
\Updated{2022-01-07 08:01:13}
\Level{B}
\SubArea{Volume and surface area formulas}
\IMG{1103170702.pdf}
\LANGen{
\Question{
Let there be a cone with the base diameter of \( 8\,\mathrm{cm} \) and the slant height of \( 5\,\mathrm{cm} \). Find the volume and surface area of the cone. Leave your answer in terms of \( \pi \).}
{\( V=16\pi\,\mathrm{cm}^3 \), \( S=36\pi\,\mathrm{cm}^2 \)}{\( V=64\pi\,\mathrm{cm}^3 \), \( S=104\pi\,\mathrm{cm}^2 \)}{\( V=64\pi\,\mathrm{cm}^3 \), \( S=104\pi\,\mathrm{cm}^2 \)}{\( V=16\pi\,\mathrm{cm}^3 \), \( S=28\pi\,\mathrm{cm}^2 \)}{}{}}
\LANGpl{
\Question{
Dany jest stożek, średnica jego podstawy jest równa \( 8\,\mathrm{cm} \) a długość tworzącej to \( 5\,\mathrm{cm} \).  Oblicz objętość i pole powierzchni stożka. Podaj wynik jako wielokrotność  \( \pi \).}
{\( V=16\pi\,\mathrm{cm}^3 \), \( S=36\pi\,\mathrm{cm}^2 \)}{\( V=64\pi\,\mathrm{cm}^3 \), \( S=104\pi\,\mathrm{cm}^2 \)}{\( V=64\pi\,\mathrm{cm}^3 \), \( S=104\pi\,\mathrm{cm}^2 \)}{\( V=16\pi\,\mathrm{cm}^3 \), \( S=28\pi\,\mathrm{cm}^2 \)}{}{}}
\LANGcs{
\Question{
Určete objem a povrch rotačního kužele s podstavou o průměru \( 8\,\mathrm{cm} \), strana kužele má délku  \( 5\,\mathrm{cm} \). Výsledek vyjádřete jako násobek \( \pi \).}
{\( V=16\pi\,\mathrm{cm}^3 \), \( S=36\pi\,\mathrm{cm}^2 \)}{\( V=64\pi\,\mathrm{cm}^3 \), \( S=104\pi\,\mathrm{cm}^2 \)}{\( V=64\pi\,\mathrm{cm}^3 \), \( S=104\pi\,\mathrm{cm}^2 \)}{\( V=16\pi\,\mathrm{cm}^3 \), \( S=28\pi\,\mathrm{cm}^2 \)}{}{}}
\LANGsk{
\Question{
Určte objem a povrch rotačného kužeľa s podstavou s priemerom \( 8\,\mathrm{cm} \), strana kužeľa má dĺžku  \( 5\,\mathrm{cm} \). Výsledok vyjadrite ako násobok \( \pi \).}
{\( V=16\pi\,\mathrm{cm}^3 \), \( S=36\pi\,\mathrm{cm}^2 \)}{\( V=64\pi\,\mathrm{cm}^3 \), \( S=104\pi\,\mathrm{cm}^2 \)}{\( V=64\pi\,\mathrm{cm}^3 \), \( S=104\pi\,\mathrm{cm}^2 \)}{\( V=16\pi\,\mathrm{cm}^3 \), \( S=28\pi\,\mathrm{cm}^2 \)}{}{}}
\LANGes{
\Question{
Averigua el volumen y la superficie de un cono sabiendo que el radio de su base es de \( 8\,\mathrm{cm} \) y el lado del cono mide \( 5\,\mathrm{cm} \). Expresa el resultado como múltiplo de \( \pi \).}
{\( V=16\pi\,\mathrm{cm}^3 \), \( S=36\pi\,\mathrm{cm}^2 \)}{\( V=64\pi\,\mathrm{cm}^3 \), \( S=104\pi\,\mathrm{cm}^2 \)}{\( V=64\pi\,\mathrm{cm}^3 \), \( S=104\pi\,\mathrm{cm}^2 \)}{\( V=16\pi\,\mathrm{cm}^3 \), \( S=28\pi\,\mathrm{cm}^2 \)}{}{}}
\End

\Data\ID{1103170703}
\Updated{2022-01-07 08:01:52}
\Level{B}
\SubArea{Volume and surface area formulas}
\IMG{1103170703.pdf}
\LANGen{
\Question{
Let there be a cone with the volume of \( 50\pi\,\mathrm{cm}^3 \) and the height of \( 6\,\mathrm{cm} \). Find its diameter.}
{\( 10\,\mathrm{cm} \)}{\( 25\,\mathrm{cm} \)}{\( 8.3\,\mathrm{cm} \)}{\( 5\,\mathrm{cm} \)}{}{}}
\LANGpl{
\Question{
Dany jest stożek, którego objętość jest równa \( 50\pi\,\mathrm{cm}^3 \) a wysokość wynosi \( 6\,\mathrm{cm} \). Oblicz jego średnicę.}
{\( 10\,\mathrm{cm} \)}{\( 25\,\mathrm{cm} \)}{\( 8{,}3\,\mathrm{cm} \)}{\( 5\,\mathrm{cm} \)}{}{}}
\LANGcs{
\Question{
Rotační kužel má objem \( 50\pi\,\mathrm{cm}^3 \) a výšku \( 6\,\mathrm{cm} \). Určete průměr podstavy.}
{\( 10\,\mathrm{cm} \)}{\( 25\,\mathrm{cm} \)}{\( 8{,}3\,\mathrm{cm} \)}{\( 5\,\mathrm{cm} \)}{}{}}
\LANGsk{
\Question{
Rotačný kužeľ má objem \( 50\pi\,\mathrm{cm}^3 \) a výšku \( 6\,\mathrm{cm} \). Určte priemer podstavy.}
{\( 10\,\mathrm{cm} \)}{\( 25\,\mathrm{cm} \)}{\( 8{,}3\,\mathrm{cm} \)}{\( 5\,\mathrm{cm} \)}{}{}}
\LANGes{
\Question{
Un cono tiene un volumen de \( 50\pi\,\mathrm{cm}^3 \) y una altura de \( 6\,\mathrm{cm} \). Averigua el diámetro de la base.}
{\( 10\,\mathrm{cm} \)}{\( 25\,\mathrm{cm} \)}{\( 8{,}3\,\mathrm{cm} \)}{\( 5\,\mathrm{cm} \)}{}{}}
\End

\Data\ID{1103170704}
\Updated{2022-04-23 05:04:11}
\Level{B}
\SubArea{Volume and surface area formulas}
\IMG{1103170704.pdf}
\LANGen{
\Question{
Find the lateral area of the cone with the diameter of \( 24\,\mathrm{cm} \) and the perpendicular height of \( 5\,\mathrm{cm} \).}
{\( 156\pi\,\mathrm{cm}^2 \)}{\( 60\pi\,\mathrm{cm}^2 \)}{\( 52\pi\,\mathrm{cm}^2 \)}{\( 624\pi\,\mathrm{cm}^2 \)}{}{}}
\LANGpl{
\Question{
Oblicz pole powierzchni bocznej stożka, którego średnica jest równa  \( 24\,\mathrm{cm} \), a długość prostopadłej wysokości to  \( 5\,\mathrm{cm} \).}
{\( 156\pi\,\mathrm{cm}^2 \)}{\( 60\pi\,\mathrm{cm}^2 \)}{\( 52\pi\,\mathrm{cm}^2 \)}{\( 624\pi\,\mathrm{cm}^2 \)}{}{}}
\LANGcs{
\Question{
Určete obsah pláště rotačního kužele s průměrem podstavy \( 24\,\mathrm{cm} \) a výškou \( 5\,\mathrm{cm} \).}
{\( 156\pi\,\mathrm{cm}^2 \)}{\( 60\pi\,\mathrm{cm}^2 \)}{\( 52\pi\,\mathrm{cm}^2 \)}{\( 624\pi\,\mathrm{cm}^2 \)}{}{}}
\LANGsk{
\Question{
Určte obsah plášťa rotačného kužeľa s priemerom podstavy \( 24\,\mathrm{cm} \) a výškou \( 5\,\mathrm{cm} \).}
{\( 156\pi\,\mathrm{cm}^2 \)}{\( 60\pi\,\mathrm{cm}^2 \)}{\( 52\pi\,\mathrm{cm}^2 \)}{\( 624\pi\,\mathrm{cm}^2 \)}{}{}}
\LANGes{
\Question{
Averigua la superficie lateral de un cono sabiendo que el radio de su base es de \( 24\,\mathrm{cm} \) y su altura mide \( 5\,\mathrm{cm} \).}
{\( 156\pi\,\mathrm{cm}^2 \)}{\( 60\pi\,\mathrm{cm}^2 \)}{\( 52\pi\,\mathrm{cm}^2 \)}{\( 624\pi\,\mathrm{cm}^2 \)}{}{}}
\End

\Data\ID{1103170705}
\Updated{2022-04-23 05:04:57}
\Level{B}
\SubArea{Volume and surface area formulas}
\IMG{1103170705.pdf}
\LANGen{
\Question{
A cone has the lateral area of \( 72\,\mathrm{dm}^2 \) and the slant height of \( 80\,\mathrm{cm} \). Find the volume of this cone. Round your answer to \( 2 \) decimal places.}
{\( 64.20\,\mathrm{dm}^3 \)}{\( 192.59\,\mathrm{dm}^3 \)}{\( 69.74\,\mathrm{dm}^3 \)}{\( 23.25\,\mathrm{dm}^3 \)}{}{}}
\LANGpl{
\Question{
Dany jest stożek o polu powierzchni bocznej równej  \( 72\,\mathrm{dm}^2 \), długość jego tworzącej wynosi \( 80\,\mathrm{cm} \). Oblicz objętość stożka. Zaokrągli wynik do  \( 2 \) miejsc po przecinku.}
{\( 64{,}20\,\mathrm{dm}^3 \)}{\( 192{,}59\,\mathrm{dm}^3 \)}{\( 69{,}74\,\mathrm{dm}^3 \)}{\( 23{,}25\,\mathrm{dm}^3 \)}{}{}}
\LANGcs{
\Question{
Rotační kužel má obsah pláště \( 72\,\mathrm{dm}^2 \) a stranu o velikosti  \( 80\,\mathrm{cm} \). Určete jeho objem. Výsledek vyjádřete s přesností na \( 2 \) desetinná místa.}
{\( 64{,}20\,\mathrm{dm}^3 \)}{\( 192{,}59\,\mathrm{dm}^3 \)}{\( 69{,}74\,\mathrm{dm}^3 \)}{\( 23{,}25\,\mathrm{dm}^3 \)}{}{}}
\LANGsk{
\Question{
Rotačný kužeľ má obsah plášťa \( 72\,\mathrm{dm}^2 \) a stranu s veľkosťou  \( 80\,\mathrm{cm} \). Určte jeho objem. Výsledok vyjadrite s presnosťou na \( 2 \) desatinné miesta.}
{\( 64{,}20\,\mathrm{dm}^3 \)}{\( 192{,}59\,\mathrm{dm}^3 \)}{\( 69{,}74\,\mathrm{dm}^3 \)}{\( 23{,}25\,\mathrm{dm}^3 \)}{}{}}
\LANGes{
\Question{
El área lateral del un cono es de \( 72\,\mathrm{dm}^2 \)y su generatriz es de \( 80\,\mathrm{cm} \). Averigua su volumen. Espresa el resultado con una exactitud de \( 2 \) cifras decimales.}
{\( 64{,}20\,\mathrm{dm}^3 \)}{\( 192{,}59\,\mathrm{dm}^3 \)}{\( 69{,}74\,\mathrm{dm}^3 \)}{\( 23{,}25\,\mathrm{dm}^3 \)}{}{}}
\End

\Data\ID{1103189201}
\Updated{2022-01-06 08:01:59}
\Level{B}
\SubArea{Volume and surface area formulas}
\IMG{1103189201.pdf}
\LANGen{
\Question{
Find the volume of a square pyramid (see the picture) with the base edge length of \( 6\,\mathrm{cm} \) and the height of \( 8\,\mathrm{cm} \).}
{\( 96\,\mathrm{cm}^3 \)}{\( 288\,\mathrm{cm}^3 \)}{\( 768\,\mathrm{cm}^3 \)}{\( 128\,\mathrm{cm}^3 \)}{}{}}
\LANGpl{
\Question{
Oblicz objętość ostrosłupa czworokątnego (spójrz na rysunek),  długość krawędzi podstawy jest równa \( 6\,\mathrm{cm} \), jego wysokość wynosi \( 8\,\mathrm{cm} \).}
{\( 96\,\mathrm{cm}^3 \)}{\( 288\,\mathrm{cm}^3 \)}{\( 768\,\mathrm{cm}^3 \)}{\( 128\,\mathrm{cm}^3 \)}{}{}}
\LANGcs{
\Question{
Určete objem pravidelného čtyřbokého jehlanu (viz obrázek) s podstavnou hranou délky \( 6\,\mathrm{cm} \) a výškou \( 8\,\mathrm{cm} \).}
{\( 96\,\mathrm{cm}^3 \)}{\( 288\,\mathrm{cm}^3 \)}{\( 768\,\mathrm{cm}^3 \)}{\( 128\,\mathrm{cm}^3 \)}{}{}}
\LANGsk{
\Question{
Určte objem pravidelného štvorbokého ihlana (viď obrázok) s  hranou podstavy dĺžky \( 6\,\mathrm{cm} \) a výškou \( 8\,\mathrm{cm} \).}
{\( 96\,\mathrm{cm}^3 \)}{\( 288\,\mathrm{cm}^3 \)}{\( 768\,\mathrm{cm}^3 \)}{\( 128\,\mathrm{cm}^3 \)}{}{}}
\LANGes{
\Question{
Averigua el volumen de una pirámide rectangular. (mira el dibujo) sabiendo que el borde de la base es de\( 6\,\mathrm{cm} \) y su altura mide \( 8\,\mathrm{cm} \).}
{\( 96\,\mathrm{cm}^3 \)}{\( 288\,\mathrm{cm}^3 \)}{\( 768\,\mathrm{cm}^3 \)}{\( 128\,\mathrm{cm}^3 \)}{}{}}
\End

\Data\ID{1103189202}
\Updated{2022-01-07 08:01:10}
\Level{B}
\SubArea{Volume and surface area formulas}
\IMG{1103189202.pdf}
\LANGen{
\Question{
Find the surface area of a square pyramid (see the picture) with the base edge length of \( 6\,\mathrm{cm} \) and the height of \( 8\,\mathrm{cm} \).}
{\( 12\left(3+\sqrt{73}\right)\,\mathrm{cm}^2 \)}{\( 132\,\mathrm{cm}^2 \)}{\( 156\,\mathrm{cm}^2 \)}{\( 12\left(3+2\sqrt{73}\right)\,\mathrm{cm}^2 \)}{}{}}
\LANGpl{
\Question{
Oblicz pole powierzchni  ostrosłupa czworokątnego (spójrz na rysunek), długość krawędzi podstawy jest równa \( 6\,\mathrm{cm} \), a wysokość wynosi \( 8\,\mathrm{cm} \).}
{\( 12\left(3+\sqrt{73}\right)\,\mathrm{cm}^2 \)}{\( 132\,\mathrm{cm}^2 \)}{\( 156\,\mathrm{cm}^2 \)}{\( 12\left(3+2\sqrt{73}\right)\,\mathrm{cm}^2 \)}{}{}}
\LANGcs{
\Question{
Určete povrch pravidelného čtyřbokého jehlanu (viz obrázek) s podstavnou hranou délky  \( 6\,\mathrm{cm} \) a výškou \( 8\,\mathrm{cm} \).}
{\( 12\left(3+\sqrt{73}\right)\,\mathrm{cm}^2 \)}{\( 132\,\mathrm{cm}^2 \)}{\( 156\,\mathrm{cm}^2 \)}{\( 12\left(3+2\sqrt{73}\right)\,\mathrm{cm}^2 \)}{}{}}
\LANGsk{
\Question{
Určte povrch pravidelného štvorbokého ihlana (viď obrázok) s  hranou podstavy dĺžky  \( 6\,\mathrm{cm} \) a výškou \( 8\,\mathrm{cm} \).}
{\( 12\left(3+\sqrt{73}\right)\,\mathrm{cm}^2 \)}{\( 132\,\mathrm{cm}^2 \)}{\( 156\,\mathrm{cm}^2 \)}{\( 12\left(3+2\sqrt{73}\right)\,\mathrm{cm}^2 \)}{}{}}
\LANGes{
\Question{
Averigua la superficie de un prisma cuadrangular (observa el dibujo). La longitud del lado de la base es \( 6\,\mathrm{cm} \) y su altura es \( 8\,\mathrm{cm} \).}
{\( 12\left(3+\sqrt{73}\right)\,\mathrm{cm}^2 \)}{\( 132\,\mathrm{cm}^2 \)}{\( 156\,\mathrm{cm}^2 \)}{\( 12\left(3+2\sqrt{73}\right)\,\mathrm{cm}^2 \)}{}{}}
\End

\Data\ID{1103189203}
\Updated{2022-01-07 08:01:20}
\Level{B}
\SubArea{Volume and surface area formulas}
\IMG{1103189203.pdf}
\LANGen{
\Question{
Find the volume of a square pyramid (see the picture) with the base edge length of \( 6\,\mathrm{cm} \) and the slant height of \( 5\,\mathrm{cm} \).}
{\( 48\,\mathrm{cm}^3 \)}{\( 144\,\mathrm{cm}^3 \)}{\( 60\,\mathrm{cm}^3 \)}{\( 180\,\mathrm{cm}^3 \)}{}{}}
\LANGpl{
\Question{
Oblicz objętość ostrosłupa czworokątnego (spójrz na rysunek), długość krawędzi podstawy jest równa \( 6\,\mathrm{cm} \), a długość tworzącej \( 5\,\mathrm{cm} \).}
{\( 48\,\mathrm{cm}^3 \)}{\( 144\,\mathrm{cm}^3 \)}{\( 60\,\mathrm{cm}^3 \)}{\( 180\,\mathrm{cm}^3 \)}{}{}}
\LANGcs{
\Question{
Určete objem pravidelného čtyřbokého jehlanu (viz obrázek) s délkou podstavné hrany  \( 6\,\mathrm{cm} \) a stěnovou výškou \( 5\,\mathrm{cm} \).}
{\( 48\,\mathrm{cm}^3 \)}{\( 144\,\mathrm{cm}^3 \)}{\( 60\,\mathrm{cm}^3 \)}{\( 180\,\mathrm{cm}^3 \)}{}{}}
\LANGsk{
\Question{
Určte objem pravidelného štvorbokého ihlana (viď obrázok) s dĺžkou hrany podstavy  \( 6\,\mathrm{cm} \) a stenovou výškou \( 5\,\mathrm{cm} \).}
{\( 48\,\mathrm{cm}^3 \)}{\( 144\,\mathrm{cm}^3 \)}{\( 60\,\mathrm{cm}^3 \)}{\( 180\,\mathrm{cm}^3 \)}{}{}}
\LANGes{
\Question{
Averigua el volumen de una pirámide cuadrangular (observa el dibujo) sabiendo que la longitud del lado de la base es de \( 6\,\mathrm{cm} \) y su apotema es de \( 5\,\mathrm{cm} \).}
{\( 48\,\mathrm{cm}^3 \)}{\( 144\,\mathrm{cm}^3 \)}{\( 60\,\mathrm{cm}^3 \)}{\( 180\,\mathrm{cm}^3 \)}{}{}}
\End

\Data\ID{1103189204}
\Updated{2022-01-07 08:01:57}
\Level{B}
\SubArea{Volume and surface area formulas}
\IMG{1103189204.pdf}
\LANGen{
\Question{
The edge length of the regular tetrahedron shown in the picture is \( 6\,\mathrm{cm} \). Find the volume of this tetrahedron.}
{\( 18\sqrt2\,\mathrm{cm}^3 \)}{\( 36\sqrt2\,\mathrm{cm}^3 \)}{\( 18\sqrt3\,\mathrm{cm}^3 \)}{\( 54\sqrt2\,\mathrm{cm}^3 \)}{}{}}
\LANGpl{
\Question{
Długość krawędzi czworościanu prawidłowego zamieszczonego na rysunku wynosi \( 6\,\mathrm{cm} \). Oblicz objętość czworościanu.}
{\( 18\sqrt2\,\mathrm{cm}^3 \)}{\( 36\sqrt2\,\mathrm{cm}^3 \)}{\( 18\sqrt3\,\mathrm{cm}^3 \)}{\( 54\sqrt2\,\mathrm{cm}^3 \)}{}{}}
\LANGcs{
\Question{
Vypočítejte objem pravidelného čtyřstěnu (viz obrázek) s délkou hrany  \( 6\,\mathrm{cm} \).}
{\( 18\sqrt2\,\mathrm{cm}^3 \)}{\( 36\sqrt2\,\mathrm{cm}^3 \)}{\( 18\sqrt3\,\mathrm{cm}^3 \)}{\( 54\sqrt2\,\mathrm{cm}^3 \)}{}{}}
\LANGsk{
\Question{
Vypočítajte objem pravidelného štvorstena (viď obrázok) s dĺžkou hrany  \( 6\,\mathrm{cm} \).}
{\( 18\sqrt2\,\mathrm{cm}^3 \)}{\( 36\sqrt2\,\mathrm{cm}^3 \)}{\( 18\sqrt3\,\mathrm{cm}^3 \)}{\( 54\sqrt2\,\mathrm{cm}^3 \)}{}{}}
\LANGes{
\Question{
Averigua el volumen del tetraedro regular (observa el dibujo) cuyo lado mide  \( 6\,\mathrm{cm} \).}
{\( 18\sqrt2\,\mathrm{cm}^3 \)}{\( 36\sqrt2\,\mathrm{cm}^3 \)}{\( 18\sqrt3\,\mathrm{cm}^3 \)}{\( 54\sqrt2\,\mathrm{cm}^3 \)}{}{}}
\End

\Data\ID{1103189205}
\Updated{2022-04-23 05:04:57}
\Level{B}
\SubArea{Volume and surface area formulas}
\IMG{1103189205.pdf}
\LANGen{
\Question{
The edge length of the regular tetrahedron shown in the picture is \( 6\,\mathrm{cm} \). Find the height of this tetrahedron.}
{\( 2\sqrt6\,\mathrm{cm} \)}{\( 6\sqrt6\,\mathrm{cm} \)}{\( 3\sqrt3\,\mathrm{cm} \)}{\( 6\sqrt3\,\mathrm{cm} \)}{}{}}
\LANGpl{
\Question{
Długość krawędzi czworościanu prawidłowego zamieszczonego na rysunku wynosi \( 6\,\mathrm{cm} \). Oblicz wysokość czworościanu.}
{\( 2\sqrt6\,\mathrm{cm} \)}{\( 6\sqrt6\,\mathrm{cm} \)}{\( 3\sqrt3\,\mathrm{cm} \)}{\( 6\sqrt3\,\mathrm{cm} \)}{}{}}
\LANGcs{
\Question{
Určete výšku pravidelného čtyřstěnu (viz obrázek) s délkou hrany  \( 6\,\mathrm{cm} \).}
{\( 2\sqrt6\,\mathrm{cm} \)}{\( 6\sqrt6\,\mathrm{cm} \)}{\( 3\sqrt3\,\mathrm{cm} \)}{\( 6\sqrt3\,\mathrm{cm} \)}{}{}}
\LANGsk{
\Question{
Určte výšku pravidelného štvorstena (viď obrázok) s dĺžkou hrany  \( 6\,\mathrm{cm} \).}
{\( 2\sqrt6\,\mathrm{cm} \)}{\( 6\sqrt6\,\mathrm{cm} \)}{\( 3\sqrt3\,\mathrm{cm} \)}{\( 6\sqrt3\,\mathrm{cm} \)}{}{}}
\LANGes{
\Question{
Averigua la altura del tetraedro regular (mira el dibujo) cuyo lado mide \( 6\,\mathrm{cm} \).}
{\( 2\sqrt6\,\mathrm{cm} \)}{\( 6\sqrt6\,\mathrm{cm} \)}{\( 3\sqrt3\,\mathrm{cm} \)}{\( 6\sqrt3\,\mathrm{cm} \)}{}{}}
\End

\Data\ID{1103189206}
\Updated{2022-02-04 09:02:06}
\Level{B}
\SubArea{Volume and surface area formulas}
\IMG{1103189206.pdf}
\LANGen{
\Question{
The edge length of the regular tetrahedron shown in the picture is \( 6\,\mathrm{cm} \).  Find the surface area of this tetrahedron.}
{\( 36\sqrt3\,\mathrm{cm}^2 \)}{\( 12\sqrt3\,\mathrm{cm}^2 \)}{\( 36\sqrt6\,\mathrm{cm}^2 \)}{\( 9\sqrt3\,\mathrm{cm}^2 \)}{}{}}
\LANGpl{
\Question{
Długość krawędzi czworościanu prawidłowego zamieszczonego na rysunku wynosi \( 6\,\mathrm{cm} \).  Oblicz pole powierzchni czworościanu.}
{\( 36\sqrt3\,\mathrm{cm}^2 \)}{\( 12\sqrt3\,\mathrm{cm}^2 \)}{\( 36\sqrt6\,\mathrm{cm}^2 \)}{\( 9\sqrt3\,\mathrm{cm}^2 \)}{}{}}
\LANGcs{
\Question{
Určete povrch pravidelného čtyřstěnu (viz obrázek) s délkou hrany  \( 6\,\mathrm{cm} \).}
{\( 36\sqrt3\,\mathrm{cm}^2 \)}{\( 12\sqrt3\,\mathrm{cm}^2 \)}{\( 36\sqrt6\,\mathrm{cm}^2 \)}{\( 9\sqrt3\,\mathrm{cm}^2 \)}{}{}}
\LANGsk{
\Question{
Určte povrch pravidelného štvorstena (viď obrázok) s dĺžkou hrany  \( 6\,\mathrm{cm} \).}
{\( 36\sqrt3\,\mathrm{cm}^2 \)}{\( 12\sqrt3\,\mathrm{cm}^2 \)}{\( 36\sqrt6\,\mathrm{cm}^2 \)}{\( 9\sqrt3\,\mathrm{cm}^2 \)}{}{}}
\LANGes{
\Question{
Calcula la superficie del tetraedro regular (observa el dibujo) si su lado mide \( 6\,\mathrm{cm} \).}
{\( 36\sqrt3\,\mathrm{cm}^2 \)}{\( 12\sqrt3\,\mathrm{cm}^2 \)}{\( 36\sqrt6\,\mathrm{cm}^2 \)}{\( 9\sqrt3\,\mathrm{cm}^2 \)}{}{}}
\End

\Data\ID{1103189207}
\Updated{2022-01-07 08:01:47}
\Level{B}
\SubArea{Volume and surface area formulas}
\IMG{1103189207.pdf}
\LANGen{
\Question{
The base of a triangular pyramid is an equilateral triangle with a side of \( 6\,\mathrm{cm} \) (see the picture). The height of the pyramid is \( 4\,\mathrm{cm} \). Find the volume of the pyramid.}
{\( 12\sqrt3\,\mathrm{cm}^3 \)}{\( 12\sqrt2\,\mathrm{cm}^3 \)}{\( 24\sqrt3\,\mathrm{cm}^3 \)}{\( 36\sqrt2\,\mathrm{cm}^3 \)}{}{}}
\LANGpl{
\Question{
Dany jest ostrosłup trójkątny, jego podstawa to trójkąt równoboczny o boku równym \( 6\,\mathrm{cm} \) (spójrz na rysunek). Wysokość ostrosłupa wynosi \( 4\,\mathrm{cm} \). Oblicz objętość ostrosłupa.}
{\( 12\sqrt3\,\mathrm{cm}^3 \)}{\( 12\sqrt2\,\mathrm{cm}^3 \)}{\( 24\sqrt3\,\mathrm{cm}^3 \)}{\( 36\sqrt2\,\mathrm{cm}^3 \)}{}{}}
\LANGcs{
\Question{
Podstavu trojbokého jehlanu s výškou \( 4\,\mathrm{cm} \) tvoří rovnostranný trojúhelník s délkou strany \( 6\,\mathrm{cm} \) (viz obrázek). Určete objem jehlanu.}
{\( 12\sqrt3\,\mathrm{cm}^3 \)}{\( 12\sqrt2\,\mathrm{cm}^3 \)}{\( 24\sqrt3\,\mathrm{cm}^3 \)}{\( 36\sqrt2\,\mathrm{cm}^3 \)}{}{}}
\LANGsk{
\Question{
Podstavu trojbokého ihlana s výškou \( 4\,\mathrm{cm} \) tvorí rovnostranný trojuholník s dĺžkou strany \( 6\,\mathrm{cm} \) (viď obrázok). Určte objem ihlana.}
{\( 12\sqrt3\,\mathrm{cm}^3 \)}{\( 12\sqrt2\,\mathrm{cm}^3 \)}{\( 24\sqrt3\,\mathrm{cm}^3 \)}{\( 36\sqrt2\,\mathrm{cm}^3 \)}{}{}}
\LANGes{
\Question{
Sea un pirámide triangular cuya altura es de \( 4\,\mathrm{cm} \). Su base es un triángulo equilátero cuyos lados miden \( 6\,\mathrm{cm} \) (observa el dibujo). Averigua el volumen de la pirámide.}
{\( 12\sqrt3\,\mathrm{cm}^3 \)}{\( 12\sqrt2\,\mathrm{cm}^3 \)}{\( 24\sqrt3\,\mathrm{cm}^3 \)}{\( 36\sqrt2\,\mathrm{cm}^3 \)}{}{}}
\End

\Data\ID{1103189208}
\Updated{2022-04-23 05:04:11}
\Level{B}
\SubArea{Volume and surface area formulas}
\IMG{1103189208.pdf}
\LANGen{
\Question{
The base of a triangular pyramid is an equilateral triangle with a side of \( 6\,\mathrm{cm} \) (see the picture). The volume of the pyramid is \( 24\sqrt3\,\mathrm{cm}^3 \). Find the perpendicular height of the pyramid.}
{\( 8\,\mathrm{cm} \)}{\( 6\sqrt3\,\mathrm{cm} \)}{\( 12\,\mathrm{cm} \)}{\( 8\sqrt3\,\mathrm{cm} \)}{}{}}
\LANGpl{
\Question{
Podstawą ostrosłupa trójkątnego jest trójkąt równoboczny o boku równym \( 6\,\mathrm{cm} \) (spójrz na rysunek). Objętość ostrosłupa wynosi \( 24\sqrt3\,\mathrm{cm}^3 \).  Oblicz wysokość prostopadłą ostrosłupa.}
{\( 8\,\mathrm{cm} \)}{\( 6\sqrt3\,\mathrm{cm} \)}{\( 12\,\mathrm{cm} \)}{\( 8\sqrt3\,\mathrm{cm} \)}{}{}}
\LANGcs{
\Question{
Podstavu trojbokého jehlanu o objemu  \( 24\sqrt3\,\mathrm{cm}^3 \)  tvoří rovnostranný trojúhelník s délkou strany \( 6\,\mathrm{cm} \) (viz obrázek). Určete výšku jehlanu.}
{\( 8\,\mathrm{cm} \)}{\( 6\sqrt3\,\mathrm{cm} \)}{\( 12\,\mathrm{cm} \)}{\( 8\sqrt3\,\mathrm{cm} \)}{}{}}
\LANGsk{
\Question{
Podstavu trojbokého ihlana s objemom  \( 24\sqrt3\,\mathrm{cm}^3 \)  tvorí rovnostranný trojuholník s dĺžkou strany \( 6\,\mathrm{cm} \) (viď obrázok). Určte výšku ihlana.}
{\( 8\,\mathrm{cm} \)}{\( 6\sqrt3\,\mathrm{cm} \)}{\( 12\,\mathrm{cm} \)}{\( 8\sqrt3\,\mathrm{cm} \)}{}{}}
\LANGes{
\Question{
El volumen de un pirámide triangular es  \( 24\sqrt3\,\mathrm{cm}^3 \). Su base es un triángulo equilátero cuyo lado mide \( 6\,\mathrm{cm} \) (vea el dibujo). Averigua la altura del pirámide.}
{\( 8\,\mathrm{cm} \)}{\( 6\sqrt3\,\mathrm{cm} \)}{\( 12\,\mathrm{cm} \)}{\( 8\sqrt3\,\mathrm{cm} \)}{}{}}
\End

\Data\ID{2000003301}
\Updated{2022-06-11 07:06:35}
\Level{B}
\SubArea{Volume and surface area formulas}
\LANGen{
\Question{
The axial section of a cylinder is a square with the diagonal length of  \( 5\sqrt{2}\,\mathrm{cm} \). The lateral surface area of the cylinder is equal to:}
{\( 25\pi\,\mathrm{cm}^2 \)}{\( 25\,\mathrm{cm}^2 \)}{\( 25\sqrt{2}\,\mathrm{cm}^2 \)}{\( 25\sqrt{2}\pi\,\mathrm{cm}^2 \)}{}{}}
\LANGpl{
\Question{
Przekrój osiowy walca to kwadrat o przekątnej długości  \( 5\sqrt{2}\,\mathrm{cm} \). Powierzchnia boczna tego walca wynosi:}
{\( 25\pi\,\mathrm{cm}^2 \)}{\( 25\,\mathrm{cm}^2 \)}{\( 25\sqrt{2}\,\mathrm{cm}^2 \)}{\( 25\sqrt{2}\pi\,\mathrm{cm}^2 \)}{}{}}
\LANGcs{
\Question{
Osovým řezem válce je čtverec s délkou úhlopříčky \( 5\sqrt{2}\,\mathrm{cm} \). Vypočtěte obsah pláště válce.}
{\( 25\pi\,\mathrm{cm}^2 \)}{\( 25\,\mathrm{cm}^2 \)}{\( 25\sqrt{2}\,\mathrm{cm}^2 \)}{\( 25\sqrt{2}\pi\,\mathrm{cm}^2 \)}{}{}}
\LANGsk{
\Question{
Osovým rezom valca je štvorec s dĺžkou uhopriečky \( 5\sqrt{2}\,\mathrm{cm} \). Vypočítate obsah plášťa valca.}
{\( 25\pi\,\mathrm{cm}^2 \)}{\( 25\,\mathrm{cm}^2 \)}{\( 25\sqrt{2}\,\mathrm{cm}^2 \)}{\( 25\sqrt{2}\pi\,\mathrm{cm}^2 \)}{}{}}
\LANGes{
\Question{
La sección axial de un cilindro es un cuadrado cuya diagonal tiene una longitud de \( 5\sqrt{2}\,\mathrm{cm} \). La superficie lateral del cilindro es igual a:}
{\( 25\pi\,\mathrm{cm}^2 \)}{\( 25\,\mathrm{cm}^2 \)}{\( 25\sqrt{2}\,\mathrm{cm}^2 \)}{\( 25\sqrt{2}\pi\,\mathrm{cm}^2 \)}{}{}}
\End

\Data\ID{2000003302}
\Updated{2022-01-07 08:01:55}
\Level{B}
\SubArea{Volume and surface area formulas}
\LANGen{
\Question{
Find the volume of a regular quadrilateral pyramid if the bottom edge is \(3\,cm\)  long and the pyramid height is \(5\,cm\).}
{\( 15\,\mathrm{cm}^3 \)}{\( 75\,\mathrm{cm}^3 \)}{\( 25\,\mathrm{cm}^3 \)}{\( 45\,\mathrm{cm}^3 \)}{}{}}
\LANGpl{
\Question{
Znajdź objętość ostrosłupa prawidłowego czworokątnego o krawędzi podstawy \(3\,cm\) i wysokości \(5\,cm\).}
{\( 15\,\mathrm{cm}^3 \)}{\( 75\,\mathrm{cm}^3 \)}{\( 25\,\mathrm{cm}^3 \)}{\( 45\,\mathrm{cm}^3 \)}{}{}}
\LANGcs{
\Question{
Vypočítejte objem čtyřbokého jehlanu, jehož podstava má stranu o délce \(3\,cm\) a jehož výška je \(5\,cm\). .}
{\( 15\,\mathrm{cm}^3 \)}{\( 75\,\mathrm{cm}^3 \)}{\( 25\,\mathrm{cm}^3 \)}{\( 45\,\mathrm{cm}^3 \)}{}{}}
\LANGsk{
\Question{
Vypočítajte objem štvorbokého ihlanu, ktorého podstava má stranu s dĺžkou \(3\,cm\) a jeho výška je \(5\,cm\).}
{\( 15\,\mathrm{cm}^3 \)}{\( 75\,\mathrm{cm}^3 \)}{\( 25\,\mathrm{cm}^3 \)}{\( 45\,\mathrm{cm}^3 \)}{}{}}
\LANGes{
\Question{
Calcula el volumen de un pirámide cuya base es un cuadrado de lado \(3\,cm\) y cuya altura es de \(5\,cm\).}
{\( 15\,\mathrm{cm}^3 \)}{\( 75\,\mathrm{cm}^3 \)}{\( 25\,\mathrm{cm}^3 \)}{\( 45\,\mathrm{cm}^3 \)}{}{}}
\End

\Data\ID{2000003303}
\Updated{2022-01-07 08:01:37}
\Level{B}
\SubArea{Volume and surface area formulas}
\LANGen{
\Question{
The volume of a regular quadrilateral pyramid is \( 432\,\mathrm{cm} ^3\) and the base edge of the pyramid has the length equal to \( 12\,\mathrm{cm} \). The height of the pyramid is:}
{\( 9\,\mathrm{cm} \)}{\( 3\,\mathrm{cm} \)}{\( 36\,\mathrm{cm} \)}{\( 27\,\mathrm{cm} \)}{}{}}
\LANGpl{
\Question{
Objętość ostrosłupa prawidłowego czworokątnego wynosi  \( 432\,\mathrm{cm} ^3\),  a krawędź podstawy ma długość \( 12\,\mathrm{cm} \). Wysokość tego ostrosłupa wynosi:}
{\( 9\,\mathrm{cm} \)}{\( 3\,\mathrm{cm} \)}{\( 36\,\mathrm{cm} \)}{\( 27\,\mathrm{cm} \)}{}{}}
\LANGcs{
\Question{
Je dán čtyřboký jehlan, jehož objem je \( 432\,\mathrm{cm} ^3\) a velikost strany podstavy \( 12\,\mathrm{cm} \). Jaká je výška tohoto jehlanu?}
{\( 9\,\mathrm{cm} \)}{\( 3\,\mathrm{cm} \)}{\( 36\,\mathrm{cm} \)}{\( 27\,\mathrm{cm} \)}{}{}}
\LANGsk{
\Question{
Je daný štvorboký ihlan, ktorého objem je \( 432\,\mathrm{cm} ^3\) a veľkosť strany podstavy \( 12\,\mathrm{cm} \). Aká je výška tohto ihlanu?}
{\( 9\,\mathrm{cm} \)}{\( 3\,\mathrm{cm} \)}{\( 36\,\mathrm{cm} \)}{\( 27\,\mathrm{cm} \)}{}{}}
\LANGes{
\Question{
Tenemos una pirámide de base cuadrada cuyo volumen es de \( 432\,\mathrm{cm} ^3\) y cuya base tiene lado de \( 12\,\mathrm{cm} \). ¿Cuánto mide la altura del pirámide?}
{\( 9\,\mathrm{cm} \)}{\( 3\,\mathrm{cm} \)}{\( 36\,\mathrm{cm} \)}{\( 27\,\mathrm{cm} \)}{}{}}
\End

\Data\ID{2000003304}
\Updated{2022-01-07 08:01:14}
\Level{B}
\SubArea{Volume and surface area formulas}
\LANGen{
\Question{
The radii of the bases of a cylinder and a cone are equal and the height of the cylinder is twice the height of the cone. The volume of the cylinder and the volume of the cone are in the ratio:}
{\( 6\ \colon 1\)}{\( 3\ \colon 1\)}{\( 2\ \colon 1\)}{\( 12\ \colon 1\)}{}{}}
\LANGpl{
\Question{
Promienie podstaw walca i stożka są równe, a wysokość walca jest dwukrotnie większa od wysokości stożka. Stosunek objętości walca do objętości stożka wynosi:}
{\( 6\ \colon 1\)}{\( 3\ \colon 1\)}{\( 2\ \colon 1\)}{\( 12\ \colon 1\)}{}{}}
\LANGcs{
\Question{
Je dán válec a kužel. Poloměry jejich podstav jsou shodné a výška válce je dvojnásobkem výšky kužele. V jakém poměru je objem válce a objem kužele?}
{\( 6\ \colon 1\)}{\( 3\ \colon 1\)}{\( 2\ \colon 1\)}{\( 12\ \colon 1\)}{}{}}
\LANGsk{
\Question{
Je daný valec a kužeľ. Polomery ich podstáv sú zhodné a výška valca je dvojnásobkom výšky kužeľa. V akom pomere je objem valca a objem kužeľa?}
{\( 6\ \colon 1\)}{\( 3\ \colon 1\)}{\( 2\ \colon 1\)}{\( 12\ \colon 1\)}{}{}}
\LANGes{
\Question{
Tenemos un cilindro y un cono. Los radios de sus bases son iguales y la altura del cilindro es doble la altura del cono. ¿En qué proporción están los volumenes del cilindro y del cono?}
{\( 6\ \colon 1\)}{\( 3\ \colon 1\)}{\( 2\ \colon 1\)}{\( 12\ \colon 1\)}{}{}}
\End

\Data\ID{2000003305}
\Updated{2022-06-11 07:06:29}
\Level{B}
\SubArea{Volume and surface area formulas}
\LANGen{
\Question{
The area of the axial section of a cylinder is \( 12\,\mathrm{cm}^2 \). The lateral surface area of the cylinder is:}
{\( 12\pi\,\mathrm{cm}^2 \)}{\( 36\pi\,\mathrm{cm}^2 \)}{\( 12\,\mathrm{cm}^2 \)}{\( 36\,\mathrm{cm}^2 \)}{}{}}
\LANGpl{
\Question{
Pole przekroju osiowego walca wynosi \( 12\,\mathrm{cm}^2 \). Pole powierzchni bocznej walca wynosi:}
{\( 12\pi\,\mathrm{cm}^2 \)}{\( 36\pi\,\mathrm{cm}^2 \)}{\( 12\,\mathrm{cm}^2 \)}{\( 36\,\mathrm{cm}^2 \)}{}{}}
\LANGcs{
\Question{
Obsah osového řezu válce je \( 12\,\mathrm{cm}^2 \). Jaký je obsah pláště tohoto válce?}
{\( 12\pi\,\mathrm{cm}^2 \)}{\( 36\pi\,\mathrm{cm}^2 \)}{\( 12\,\mathrm{cm}^2 \)}{\( 36\,\mathrm{cm}^2 \)}{}{}}
\LANGsk{
\Question{
Obsah osového rezu valca je \( 12\,\mathrm{cm}^2 \). Aký je obsah plášťa tohto valca?}
{\( 12\pi\,\mathrm{cm}^2 \)}{\( 36\pi\,\mathrm{cm}^2 \)}{\( 12\,\mathrm{cm}^2 \)}{\( 36\,\mathrm{cm}^2 \)}{}{}}
\LANGes{
\Question{
La superficie de la sección axial de un cilindro es \( 12\,\mathrm{cm}^2 \). La superficie lateral del cilindro es:}
{\( 12\pi\,\mathrm{cm}^2 \)}{\( 36\pi\,\mathrm{cm}^2 \)}{\( 12\,\mathrm{cm}^2 \)}{\( 36\,\mathrm{cm}^2 \)}{}{}}
\End

\Data\ID{2000003306}
\Updated{2022-06-11 07:06:59}
\Level{B}
\SubArea{Volume and surface area formulas}
\LANGen{
\Question{
A rectangle with the sides of \( 4\,\mathrm{cm} \) and \( 6\,\mathrm{cm} \) is rotated around its longer side thus giving rise to a solid. Find the volume of such a solid?}
{\( 96\pi\,\mathrm{cm}^3 \)}{\( 48\pi\,\mathrm{cm}^3 \)}{\( 96\,\mathrm{cm}^3 \)}{\( 144\pi\,\mathrm{cm}^3 \)}{}{}}
\LANGpl{
\Question{
Prostokąt o bokach  \( 4\,\mathrm{cm} \) i \( 6\,\mathrm{cm} \) obraca się wokół dłuższego boku, tworząc w ten sposób bryłę.  Jaka jest objętość tej bryły?}
{\( 96\pi\,\mathrm{cm}^3 \)}{\( 48\pi\,\mathrm{cm}^3 \)}{\( 96\,\mathrm{cm}^3 \)}{\( 144\pi\,\mathrm{cm}^3 \)}{}{}}
\LANGcs{
\Question{
Obdélník se stranami o délce \( 4\,\mathrm{cm} \) a \( 6\,\mathrm{cm} \) otočíme kolem delší strany, čímž získáme těleso. Jaký je objem tohoto tělesa?}
{\( 96\pi\,\mathrm{cm}^3 \)}{\( 48\pi\,\mathrm{cm}^3 \)}{\( 96\,\mathrm{cm}^3 \)}{\( 144\pi\,\mathrm{cm}^3 \)}{}{}}
\LANGsk{
\Question{
Obdĺžnik so stranami s dĺžkou \( 4\,\mathrm{cm} \) a \( 6\,\mathrm{cm} \) otočíme okolo dlhšej strany, čím získame teleso. Aký je objem tohto telesa?}
{\( 96\pi\,\mathrm{cm}^3 \)}{\( 48\pi\,\mathrm{cm}^3 \)}{\( 96\,\mathrm{cm}^3 \)}{\( 144\pi\,\mathrm{cm}^3 \)}{}{}}
\LANGes{
\Question{
Un rectángulo cuyos lados miden \( 4\,\mathrm{cm} \) y \( 6\,\mathrm{cm} \) gira alrededor de su lado más largo generando un cuerpo. Halla el volumen de dicho cuerpo.}
{\( 96\pi\,\mathrm{cm}^3 \)}{\( 48\pi\,\mathrm{cm}^3 \)}{\( 96\,\mathrm{cm}^3 \)}{\( 144\pi\,\mathrm{cm}^3 \)}{}{}}
\End

\Data\ID{1003191301}
\Updated{2022-01-06 08:01:10}
\Level{C}
\SubArea{Volume and surface area formulas}
\LANGen{
\Question{
A frustum of a pyramid has rectangular ends and the sides of the base are \( 8\,\mathrm{cm} \) and \( 6\,\mathrm{cm} \) long. Find the volume of the frustum knowing that the area of the top end is \( 12\,\mathrm{cm}^2 \) and the height of the frustum is \( 5\,\mathrm{cm} \).}
{\( 140\,\mathrm{cm}^3 \)}{\( 100\,\mathrm{cm}^3 \)}{\( 420\,\mathrm{cm}^3 \)}{\( 1060\,\mathrm{cm}^3 \)}{}{}}
\LANGpl{
\Question{
Dany jest ostrosłup ścięty, jego podstawy to czworokąty. Długości dolnej podstawy są równe \( 8\,\mathrm{cm} \) i \( 6\,\mathrm{cm} \). Oblicz objętość ostrosłupa  wiedząc, że pole powierzchni jego górnej podstawy jest równe \( 12\,\mathrm{cm}^2 \) a wysokość  \( 5\,\mathrm{cm} \).}
{\( 140\,\mathrm{cm}^3 \)}{\( 100\,\mathrm{cm}^3 \)}{\( 420\,\mathrm{cm}^3 \)}{\( 1060\,\mathrm{cm}^3 \)}{}{}}
\LANGcs{
\Question{
Čtyřboký komolý jehlan má výšku  \( 5\,\mathrm{cm} \). Dolní podstava tvaru obdélníku má rozměry \( 8\,\mathrm{cm} \) a \( 6\,\mathrm{cm} \), horní podstava má obsah \( 12\,\mathrm{cm}^2 \). Určete jeho objem.}
{\( 140\,\mathrm{cm}^3 \)}{\( 100\,\mathrm{cm}^3 \)}{\( 420\,\mathrm{cm}^3 \)}{\( 1060\,\mathrm{cm}^3 \)}{}{}}
\LANGsk{
\Question{
Štvorboký zrezaný ihlan má výšku  \( 5\,\mathrm{cm} \). Dolná podstava tvaru obdĺžnika má rozmery \( 8\,\mathrm{cm} \) a \( 6\,\mathrm{cm} \), horná podstava má obsah \( 12\,\mathrm{cm}^2 \). Určte jeho objem.}
{\( 140\,\mathrm{cm}^3 \)}{\( 100\,\mathrm{cm}^3 \)}{\( 420\,\mathrm{cm}^3 \)}{\( 1060\,\mathrm{cm}^3 \)}{}{}}
\LANGes{
\Question{
La altura de una pirámide cuadrangular truncada es de \( 5\,\mathrm{cm} \). La base inferior tiene forma de rectángulo y sus medidas son  de \( 8\,\mathrm{cm} \) y \( 6\,\mathrm{cm} \), la superficie de la base superior es de\( 12\,\mathrm{cm}^2 \). Averigua el volumen del pirámide.}
{\( 140\,\mathrm{cm}^3 \)}{\( 100\,\mathrm{cm}^3 \)}{\( 420\,\mathrm{cm}^3 \)}{\( 1060\,\mathrm{cm}^3 \)}{}{}}
\End

\Data\ID{1103191302}
\Updated{2022-01-07 08:01:41}
\Level{C}
\SubArea{Volume and surface area formulas}
\IMG{1103191302_0.pdf}
\LANGen{
\Question{
A frustum of a pyramid has square ends and the squares have sides \( 8\,\mathrm{cm} \) and \( 6\,\mathrm{cm} \) long, respectively. Calculate the volume of the frustum if the perpendicular distance between its ends is \( 12\,\mathrm{cm} \).}
{\( 592\,\mathrm{cm}^3 \)}{\( 9616\,\mathrm{cm}^3 \)}{\( 1776\,\mathrm{cm}^3 \)}{\( 248\,\mathrm{cm}^3 \)}{}{}}
\LANGpl{
\Question{
Dany jest ostrosłup ścięty, jego podstawy to kwadraty o bokach równych odpowiednio  \( 8\,\mathrm{cm} \) i \( 6\,\mathrm{cm} \), wysokość ostrosłupa wynosi \( 12\,\mathrm{cm} \). Oblicz objętość ostrosłupa.}
{\( 592\,\mathrm{cm}^3 \)}{\( 9616\,\mathrm{cm}^3 \)}{\( 1776\,\mathrm{cm}^3 \)}{\( 248\,\mathrm{cm}^3 \)}{}{}}
\LANGcs{
\Question{
Pravidelný čtyřboký komolý jehlan s délkou podstavných hran \( 8\,\mathrm{cm} \) a \( 6\,\mathrm{cm} \) má výšku \( 12\,\mathrm{cm} \). Určete jeho objem.}
{\( 592\,\mathrm{cm}^3 \)}{\( 9616\,\mathrm{cm}^3 \)}{\( 1776\,\mathrm{cm}^3 \)}{\( 248\,\mathrm{cm}^3 \)}{}{}}
\LANGsk{
\Question{
Pravidelný štvorboký zrezaný ihlan s dĺžkou podstavných hrán \( 8\,\mathrm{cm} \) a \( 6\,\mathrm{cm} \) má výšku \( 12\,\mathrm{cm} \). Určte jeho objem.}
{\( 592\,\mathrm{cm}^3 \)}{\( 9616\,\mathrm{cm}^3 \)}{\( 1776\,\mathrm{cm}^3 \)}{\( 248\,\mathrm{cm}^3 \)}{}{}}
\LANGes{
\Question{
Sea una pirámide cuadrangular truncada cuyos lados de la base miden \( 8\,\mathrm{cm} \) y \( 6\,\mathrm{cm} \). Su altura es de \( 12\,\mathrm{cm} \). Averigua su volumen.}
{\( 592\,\mathrm{cm}^3 \)}{\( 9616\,\mathrm{cm}^3 \)}{\( 1776\,\mathrm{cm}^3 \)}{\( 248\,\mathrm{cm}^3 \)}{}{}}
\End

\Data\ID{1103191303}
\Updated{2022-01-07 08:01:05}
\Level{C}
\SubArea{Volume and surface area formulas}
\IMG{1103191303_0.pdf}
\LANGen{
\Question{
A frustum of a pyramid has square ends and the squares have sides \( 18\,\mathrm{cm} \) and \( 6\,\mathrm{cm} \) long, respectively. Calculate the surface area of the frustum if the perpendicular distance between its ends is \( 8\,\mathrm{cm} \).}
{\( 840\,\mathrm{cm}^2 \)}{\( 360\,\mathrm{cm}^2 \)}{\( 480\,\mathrm{cm}^2 \)}{\( 804\,\mathrm{cm}^2 \)}{}{}}
\LANGpl{
\Question{
Dany jest ostrosłup ścięty, jego podstawy to kwadraty o bokach równych odpowiednio  \( 18\,\mathrm{cm} \) i \( 6\,\mathrm{cm} \), wysokość ostrosłupa wynosi \( 8\,\mathrm{cm} \). Oblicz pole powierzchni ostrosłupa.}
{\( 840\,\mathrm{cm}^2 \)}{\( 360\,\mathrm{cm}^2 \)}{\( 480\,\mathrm{cm}^2 \)}{\( 804\,\mathrm{cm}^2 \)}{}{}}
\LANGcs{
\Question{
Pravidelný čtyřboký komolý jehlan s délkou podstavných hran  \( 18\,\mathrm{cm} \) a \( 6\,\mathrm{cm} \) má výšku \( 8\,\mathrm{cm} \). Určete jeho povrch.}
{\( 840\,\mathrm{cm}^2 \)}{\( 360\,\mathrm{cm}^2 \)}{\( 480\,\mathrm{cm}^2 \)}{\( 804\,\mathrm{cm}^2 \)}{}{}}
\LANGsk{
\Question{
Pravidelný štvorboký zrezaný ihlan s dĺžkou podstavných hrán  \( 18\,\mathrm{cm} \) a \( 6\,\mathrm{cm} \) má výšku \( 8\,\mathrm{cm} \). Určte jeho povrch.}
{\( 840\,\mathrm{cm}^2 \)}{\( 360\,\mathrm{cm}^2 \)}{\( 480\,\mathrm{cm}^2 \)}{\( 804\,\mathrm{cm}^2 \)}{}{}}
\LANGes{
\Question{
Sea un pirámide cuadrangular truncada cuyos lados de base miden \( 18\,\mathrm{cm} \) y \( 6\,\mathrm{cm} \). Su altura es de \( 8\,\mathrm{cm} \). Averigua su superficie.}
{\( 840\,\mathrm{cm}^2 \)}{\( 360\,\mathrm{cm}^2 \)}{\( 480\,\mathrm{cm}^2 \)}{\( 804\,\mathrm{cm}^2 \)}{}{}}
\End

\Data\ID{1103191304}
\Updated{2022-01-07 08:01:47}
\Level{C}
\SubArea{Volume and surface area formulas}
\IMG{1103191304.pdf}
\LANGen{
\Question{
A builders bucket is in the shape of a frustum of a right circular cone as shown in the picture. Find the volume of the bucket with the top and bottom diameter of \( 10\,\mathrm{cm} \) and \( 15\,\mathrm{cm} \) and with the height of \( 18\,\mathrm{cm} \).}
{\( 712.5\pi\,\mathrm{cm}^3 \)}{\( 350\pi\,\mathrm{cm}^3 \)}{\( 2023.5\pi\,\mathrm{cm}^3 \)}{\( 2850\pi\,\mathrm{cm}^3 \)}{}{}}
\LANGpl{
\Question{
Wiadro ma kształt stożka ściętego (patrz rysunek). Jaka jest objętość wiadra, jeśli wiemy, że jego dno ma średnicę  \( 10\,\mathrm{cm} \), natomiast górna podstawa wynosi \( 15\,\mathrm{cm} \), wysokość \( 18\,\mathrm{cm} \)?}
{\( 712{,}5\pi\,\mathrm{cm}^3 \)}{\( 350\pi\,\mathrm{cm}^3 \)}{\( 2023{,}5\pi\,\mathrm{cm}^3 \)}{\( 2850\pi\,\mathrm{cm}^3 \)}{}{}}
\LANGcs{
\Question{
Kbelík má tvar komolého kužele (viz obrázek). Jaký je objem kbelíku, pokud víme, že jeho dno má průměr \( 10\,\mathrm{cm} \), průměr horní části je \( 15\,\mathrm{cm} \) a výška je \( 18\,\mathrm{cm} \)?}
{\( 712{,}5\pi\,\mathrm{cm}^3 \)}{\( 350\pi\,\mathrm{cm}^3 \)}{\( 2023{,}5\pi\,\mathrm{cm}^3 \)}{\( 2850\pi\,\mathrm{cm}^3 \)}{}{}}
\LANGsk{
\Question{
Vedro má tvar zrezaného kužeľa (viď obrázok). Aký je objem vedierka, ak vieme, že jeho dno má priemer \( 10\,\mathrm{cm} \), priemer hornej časti je \( 15\,\mathrm{cm} \) a výška je \( 18\,\mathrm{cm} \)?}
{\( 712{,}5\pi\,\mathrm{cm}^3 \)}{\( 350\pi\,\mathrm{cm}^3 \)}{\( 2023{,}5\pi\,\mathrm{cm}^3 \)}{\( 2850\pi\,\mathrm{cm}^3 \)}{}{}}
\LANGes{
\Question{
Una cubeta tiene forma de cono truncado (observa el dibujo). ¿Cuál es el volumen de la cubeta si sabemos que el diámetro del fondo es de \( 10\,\mathrm{cm} \),  el diámetro de la parte superior es de \( 15\,\mathrm{cm} \) y su altura \( 18\,\mathrm{cm} \)?}
{\( 712{,}5\pi\,\mathrm{cm}^3 \)}{\( 350\pi\,\mathrm{cm}^3 \)}{\( 2023{,}5\pi\,\mathrm{cm}^3 \)}{\( 2850\pi\,\mathrm{cm}^3 \)}{}{}}
\End

\Data\ID{1103191305}
\Updated{2022-01-07 08:01:25}
\Level{C}
\SubArea{Volume and surface area formulas}
\IMG{1103191305.pdf}
\LANGen{
\Question{
What is the area of a metal plate needed to produce one bucket? The bucket is in the shape of a frustum of a cone as shown in the picture. The top and bottom diameters are \( 23\,\mathrm{cm} \) and \( 18\,\mathrm{cm} \) and the slant height is \( 17\,\mathrm{cm} \). Round your result to \( 1 \) decimal place.}
{\( 1349.3\,\mathrm{cm}^2 \)}{\( 3207.6\,\mathrm{cm}^2 \)}{\( 2189.7\,\mathrm{cm}^2 \)}{\( 1623.2\,\mathrm{cm}^2 \)}{}{}}
\LANGpl{
\Question{
Oblicz pole powierzchni metalowej płyty potrzebnej do wyprodukowania wiadra. Wiadro ma kształt ściętego stożka, jak pokazano na rysunku. Średnice górnej  i dolnej podstawy są równe  \( 23\,\mathrm{cm} \) i \( 18\,\mathrm{cm} \), wysokość ściany bocznej wynosi \( 17\,\mathrm{cm} \). Zaokrągli wynik do  \( 1 \) miejsca po przecinku.}
{\( 1349{,}3\,\mathrm{cm}^2 \)}{\( 3207{,}6\,\mathrm{cm}^2 \)}{\( 2189{,}7\,\mathrm{cm}^2 \)}{\( 1623{,}2\,\mathrm{cm}^2 \)}{}{}}
\LANGcs{
\Question{
Kolik materiálu potřebujeme na výrobu jedné nádoby tvaru komolého kužele (viz obrázek), jsou-li průměry podstav  \( 23\,\mathrm{cm} \), \( 18\,\mathrm{cm} \) a délka strany je \( 17\,\mathrm{cm} \)? Výsledek zaokrouhlete na \( 1 \) desetinné místo.}
{\( 1349{,}3\,\mathrm{cm}^2 \)}{\( 3207{,}6\,\mathrm{cm}^2 \)}{\( 2189{,}7\,\mathrm{cm}^2 \)}{\( 1623{,}2\,\mathrm{cm}^2 \)}{}{}}
\LANGsk{
\Question{
Koľko materiálu potrebujeme na výrobu jednej nádoby tvaru zrezaného kužeľa (viď obrázok), ak sú priemery podstáv  \( 23\,\mathrm{cm} \), \( 18\,\mathrm{cm} \) a dĺžka strany je \( 17\,\mathrm{cm} \)? Výsledok zaokrúhlite na \( 1 \) desatinné miesto.}
{\( 1349{,}3\,\mathrm{cm}^2 \)}{\( 3207{,}6\,\mathrm{cm}^2 \)}{\( 2189{,}7\,\mathrm{cm}^2 \)}{\( 1623{,}2\,\mathrm{cm}^2 \)}{}{}}
\LANGes{
\Question{
¿Qué cantidad de material necesitamos para producir un recipiente en forma de cono truncado (observa el dibujo) sabiendo que los diámetros de las bases son \( 23\,\mathrm{cm} \), \( 18\,\mathrm{cm} \) y la longitud del lado es de \( 17\,\mathrm{cm} \)? Expresa el resultado con una exactitud de \( 1 \) cifra decimal.}
{\( 1349{,}3\,\mathrm{cm}^2 \)}{\( 3207{,}6\,\mathrm{cm}^2 \)}{\( 2189{,}7\,\mathrm{cm}^2 \)}{\( 1623{,}2\,\mathrm{cm}^2 \)}{}{}}
\End

\Data\ID{1103191306}
\Updated{2022-01-07 08:01:37}
\Level{C}
\SubArea{Volume and surface area formulas}
\IMG{1103191305_0.pdf}
\LANGen{
\Question{
Find the volume (in liters) of a bucket. The bucket is in the shape of frustum of a cone (see the picture) with the top and bottom diameter of \( 23\,\mathrm{cm} \) and \( 18\,\mathrm{cm} \) and the slant height of \( 17\,\mathrm{cm} \). Round your answer to \( 2 \) decimal places.}
{\( 5.58\,\mathrm{l} \)}{\( 5.65\,\mathrm{l} \)}{\( 22.32\,\mathrm{l} \)}{\( 22.56\,\mathrm{l} \)}{}{}}
\LANGpl{
\Question{
Oblicz objętość  (w litrach) wiadra. Wiadro ma kształt ściętego stożka (patrz obrazek), średnica górnej i dolnej podstawy wynosi odpowiednio \( 23\,\mathrm{cm} \) i \( 18\,\mathrm{cm} \), wysokość ściany bocznej to  \( 17\,\mathrm{cm} \). Zaokrągli wynik do  \( 2 \) po przecinku.}
{\( 5{,}58\,\mathrm{l} \)}{\( 5{,}65\,\mathrm{l} \)}{\( 22{,}32\,\mathrm{l} \)}{\( 22{,}56\,\mathrm{l} \)}{}{}}
\LANGcs{
\Question{
Určete objem nádoby tvaru komolého kužele (viz obrázek), jsou-li průměry podstav  \( 23\,\mathrm{cm} \), \( 18\,\mathrm{cm} \) a délka strany je \( 17\,\mathrm{cm} \). Výsledek zaokrouhlete na \( 2 \) desetinná místa.}
{\( 5{,}58\,\mathrm{l} \)}{\( 5{,}65\,\mathrm{l} \)}{\( 22{,}32\,\mathrm{l} \)}{\( 22{,}56\,\mathrm{l} \)}{}{}}
\LANGsk{
\Question{
Určte objem nádoby tvaru zrezaného kužeľa (viď obrázok), ak sú priemery podstáv  \( 23\,\mathrm{cm} \), \( 18\,\mathrm{cm} \) a dĺžka strany je \( 17\,\mathrm{cm} \). Výsledok zaokrúhlite na \( 2 \) desatinné miesta.}
{\( 5{,}58\,\mathrm{l} \)}{\( 5{,}65\,\mathrm{l} \)}{\( 22{,}32\,\mathrm{l} \)}{\( 22{,}56\,\mathrm{l} \)}{}{}}
\LANGes{
\Question{
Averigua el volumen de un recipiente en forma de cono truncado (observa el dibujo) sabiendo que los diámetros de las bases son \( 23\,\mathrm{cm} \), \( 18\,\mathrm{cm} \) y la longitud del lado es de \( 17\,\mathrm{cm} \). Expresa el resultado con una exactitud de \( 2 \) cifras decimales.}
{\( 5{,}58\,\mathrm{l} \)}{\( 5{,}65\,\mathrm{l} \)}{\( 22{,}32\,\mathrm{l} \)}{\( 22{,}56\,\mathrm{l} \)}{}{}}
\End

\Data\ID{1103235601}
\Updated{2022-01-07 08:01:44}
\Level{C}
\SubArea{Volume and surface area formulas}
\IMG{1103235601.pdf}
\LANGen{
\Question{
Find the volume of a right pyramid whose base is a regular hexagon each side of which is \( 6\,\mathrm{cm} \) and height \( 8\,\mathrm{cm} \) (see the picture).}
{\( 144\sqrt3\,\mathrm{cm}^3 \)}{\( 72\sqrt3\,\mathrm{cm}^3 \)}{\( 48\sqrt3\,\mathrm{cm}^3 \)}{\( 24\sqrt3\,\mathrm{cm}^3 \)}{}{}}
\LANGpl{
\Question{
Oblicz objętość ostrosłupa prawidłowego, którego podstawą jest sześciokąt foremny o boku równym  \( 6\,\mathrm{cm} \), wysokość jest równa \( 8\,\mathrm{cm} \) (spójrz na rysunek).}
{\( 144\sqrt3\,\mathrm{cm}^3 \)}{\( 72\sqrt3\,\mathrm{cm}^3 \)}{\( 48\sqrt3\,\mathrm{cm}^3 \)}{\( 24\sqrt3\,\mathrm{cm}^3 \)}{}{}}
\LANGcs{
\Question{
Vypočítejte povrch pravidelného šestibokého jehlanu s délkou podstavné hrany  \( 6\,\mathrm{cm} \) a výškou \( 8\,\mathrm{cm} \) (viz obrázek).}
{\( 144\sqrt3\,\mathrm{cm}^3 \)}{\( 72\sqrt3\,\mathrm{cm}^3 \)}{\( 48\sqrt3\,\mathrm{cm}^3 \)}{\( 24\sqrt3\,\mathrm{cm}^3 \)}{}{}}
\LANGsk{
\Question{
Vypočítajte povrch pravidelného šesťbokého ihlanu s dĺžkou hrany podstavy  \( 6\,\mathrm{cm} \) a výškou \( 8\,\mathrm{cm} \) (viď obrázok).}
{\( 144\sqrt3\,\mathrm{cm}^3 \)}{\( 72\sqrt3\,\mathrm{cm}^3 \)}{\( 48\sqrt3\,\mathrm{cm}^3 \)}{\( 24\sqrt3\,\mathrm{cm}^3 \)}{}{}}
\LANGes{
\Question{
Calcula la superficie de una pirámide hexagonal regular sabiendo que la longitud del lado de la base es \( 6\,\mathrm{cm} \) y su altura es de \( 8\,\mathrm{cm} \) (observa el dibujo).}
{\( 144\sqrt3\,\mathrm{cm}^3 \)}{\( 72\sqrt3\,\mathrm{cm}^3 \)}{\( 48\sqrt3\,\mathrm{cm}^3 \)}{\( 24\sqrt3\,\mathrm{cm}^3 \)}{}{}}
\End

\Data\ID{1103235602}
\Updated{2022-01-07 08:01:25}
\Level{C}
\SubArea{Volume and surface area formulas}
\IMG{1103235602.pdf}
\LANGen{
\Question{
Find the surface area of a right pyramid whose base is a regular hexagon each side of which is \( 6\,\mathrm{cm} \) and height \( 9\,\mathrm{cm} \) (see the picture).}
{\( 162\sqrt3\,\mathrm{cm}^2 \)}{\( 15\sqrt3\,\mathrm{cm}^2 \)}{\( 9\left(\sqrt3+6\sqrt{13}\right)\,\mathrm{cm}^2 \)}{\( 117\sqrt3\,\mathrm{cm}^2 \)}{}{}}
\LANGpl{
\Question{
Oblicz pole powierzchni ostrosłupa prawidłowego, którego podstawa to sześciokąt foremny o boku równym \( 6\,\mathrm{cm} \), wysokość wynosi \( 9\,\mathrm{cm} \) (spójrz na rysunek).}
{\( 162\sqrt3\,\mathrm{cm}^2 \)}{\( 15\sqrt3\,\mathrm{cm}^2 \)}{\( 9\left(\sqrt3+6\sqrt{13}\right)\,\mathrm{cm}^2 \)}{\( 117\sqrt3\,\mathrm{cm}^2 \)}{}{}}
\LANGcs{
\Question{
Vypočítejte povrch pravidelného šestibokého jehlanu s délkou podstavné hrany  \( 6\,\mathrm{cm} \) a výškou \( 9\,\mathrm{cm} \) (viz obrázek).}
{\( 162\sqrt3\,\mathrm{cm}^2 \)}{\( 15\sqrt3\,\mathrm{cm}^2 \)}{\( 9\left(\sqrt3+6\sqrt{13}\right)\,\mathrm{cm}^2 \)}{\( 117\sqrt3\,\mathrm{cm}^2 \)}{}{}}
\LANGsk{
\Question{
Vypočítajte povrch pravidelného šesťbokého ihlanu s dĺžkou hrany podstavy  \( 6\,\mathrm{cm} \) a výškou \( 9\,\mathrm{cm} \) (viď obrázok).}
{\( 162\sqrt3\,\mathrm{cm}^2 \)}{\( 15\sqrt3\,\mathrm{cm}^2 \)}{\( 9\left(\sqrt3+6\sqrt{13}\right)\,\mathrm{cm}^2 \)}{\( 117\sqrt3\,\mathrm{cm}^2 \)}{}{}}
\LANGes{
\Question{
Calcula la superficie de una pirámide hexagonal regular sabiendo que el lado de su base mide \( 6\,\mathrm{cm} \) y su altura es de \( 9\,\mathrm{cm} \) (observa el dibujo).}
{\( 162\sqrt3\,\mathrm{cm}^2 \)}{\( 15\sqrt3\,\mathrm{cm}^2 \)}{\( 9\left(\sqrt3+6\sqrt{13}\right)\,\mathrm{cm}^2 \)}{\( 117\sqrt3\,\mathrm{cm}^2 \)}{}{}}
\End

\Data\ID{1103235603}
\Updated{2022-01-07 08:01:30}
\Level{C}
\SubArea{Volume and surface area formulas}
\IMG{1103235603.pdf}
\LANGen{
\Question{
The base of a right pyramid is a regular hexagon of side \( 4\,\mathrm{m} \) and its slant surfaces are inclined to the horizontal at an angle of \( 30^{\circ} \) (see the picture). Find the volume.}
{\( 16\sqrt3\,\mathrm{m}^3 \)}{\( 72\sqrt3\,\mathrm{m}^3 \)}{\( 48\sqrt3\,\mathrm{m}^3 \)}{\( 24\sqrt3\,\mathrm{m}^3 \)}{}{}}
\LANGpl{
\Question{
Podstawą ostrosłupa prawidłowego jest sześciokąt foremny o boku równym \( 4\,\mathrm{m} \), powierzchnie boczne są nachylone do poziomej pod kątem  \( 30^{\circ} \) (spójrz na rysunek). Oblicz objętość.}
{\( 16\sqrt3\,\mathrm{m}^3 \)}{\( 72\sqrt3\,\mathrm{m}^3 \)}{\( 48\sqrt3\,\mathrm{m}^3 \)}{\( 24\sqrt3\,\mathrm{m}^3 \)}{}{}}
\LANGcs{
\Question{
V pravidelném šestibokém jehlanu je délka podstavné hrany  \( 4\,\mathrm{m} \) a rovina boční stěny svírá s rovinou podstavy úhel  \( 30^{\circ} \) (viz obrázek). Určete jeho objem.}
{\( 16\sqrt3\,\mathrm{m}^3 \)}{\( 72\sqrt3\,\mathrm{m}^3 \)}{\( 48\sqrt3\,\mathrm{m}^3 \)}{\( 24\sqrt3\,\mathrm{m}^3 \)}{}{}}
\LANGsk{
\Question{
V pravidelnom šesťbokom ihlane je dĺžka hrany podstavy \( 4\,\mathrm{m} \) a rovina bočnej steny zviera s rovinou podstavy uhol  \( 30^{\circ} \) (viď obrázok). Určte jeho objem.}
{\( 16\sqrt3\,\mathrm{m}^3 \)}{\( 72\sqrt3\,\mathrm{m}^3 \)}{\( 48\sqrt3\,\mathrm{m}^3 \)}{\( 24\sqrt3\,\mathrm{m}^3 \)}{}{}}
\LANGes{
\Question{
Sea una pirámide hexagonal regular cuyo lado de la base mide \( 4\,\mathrm{m} \). El ángulo entre el plano lateral y el plano de la base es de \( 30^{\circ} \) (observa el dibujo). Averigua su volumen.}
{\( 16\sqrt3\,\mathrm{m}^3 \)}{\( 72\sqrt3\,\mathrm{m}^3 \)}{\( 48\sqrt3\,\mathrm{m}^3 \)}{\( 24\sqrt3\,\mathrm{m}^3 \)}{}{}}
\End

\Data\ID{1103235604}
\Updated{2022-01-07 08:01:53}
\Level{C}
\SubArea{Volume and surface area formulas}
\IMG{1103235604.pdf}
\LANGen{
\Question{
The area of the base of a regular hexagonal pyramid is \( 54\sqrt3\,\mathrm{cm}^2 \) and the lateral edge is two times the length of the base edge (see the picture). Find the volume of the pyramid.}
{\( 324\,\mathrm{cm}^3 \)}{\( 108\sqrt3\,\mathrm{cm}^3 \)}{\( 972\,\mathrm{cm}^3 \)}{\( 216\,\mathrm{cm}^3 \)}{}{}}
\LANGpl{
\Question{
Pole powierzchni podstawy ostrosłupa prawidłowego sześciokątnego wynosi \( 54\sqrt3\,\mathrm{cm}^2 \), jego boczna krawędź jest dwa razy dłuższa niż krawędź podstawy (spójrz na rysunek). Oblicz objętość ostrosłupa.}
{\( 324\,\mathrm{cm}^3 \)}{\( 108\sqrt3\,\mathrm{cm}^3 \)}{\( 972\,\mathrm{cm}^3 \)}{\( 216\,\mathrm{cm}^3 \)}{}{}}
\LANGcs{
\Question{
Podstava pravidelného šestibokého jehlanu má obsah  \( 54\sqrt3\,\mathrm{cm}^2 \) a délka boční hrany je dvakrát větší než délka podstavné hrany (viz obrázek).  Vypočítejte objem jehlanu.}
{\( 324\,\mathrm{cm}^3 \)}{\( 108\sqrt3\,\mathrm{cm}^3 \)}{\( 972\,\mathrm{cm}^3 \)}{\( 216\,\mathrm{cm}^3 \)}{}{}}
\LANGsk{
\Question{
Podstava pravidelného šesťbokého ihlanu má obsah  \( 54\sqrt3\,\mathrm{cm}^2 \) a dĺžka bočnej hrany je dvakrát väčšia ako dĺžka podstavnej hrany (viď obrázok).  Vypočítajte objem ihlanu.}
{\( 324\,\mathrm{cm}^3 \)}{\( 108\sqrt3\,\mathrm{cm}^3 \)}{\( 972\,\mathrm{cm}^3 \)}{\( 216\,\mathrm{cm}^3 \)}{}{}}
\LANGes{
\Question{
Sea una pirámide hexagonal regular cuya superficie de la base es de \( 54\sqrt3\,\mathrm{cm}^2 \). La longitud de la altura es dos veces mayor que la longitud del lado de la base (observa el dibujo).  Calcula el volumen del pirámide.}
{\( 324\,\mathrm{cm}^3 \)}{\( 108\sqrt3\,\mathrm{cm}^3 \)}{\( 972\,\mathrm{cm}^3 \)}{\( 216\,\mathrm{cm}^3 \)}{}{}}
\End

\Data\ID{1103235605}
\Updated{2022-01-07 08:01:34}
\Level{C}
\SubArea{Volume and surface area formulas}
\IMG{1103235605.pdf}
\LANGen{
\Question{
A regular hexagonal pyramid has the perimeter of its base \( 12\sqrt3\,\mathrm{cm} \) and its slant height is \( 5\,\mathrm{cm} \). Find the surface area.}
{\( 48\sqrt3\,\mathrm{cm}^2 \)}{\( 72\sqrt3\,\mathrm{cm}^2 \)}{\( 30\sqrt3\,\mathrm{cm}^2 \)}{\( 96\sqrt3\,\mathrm{cm}^2 \)}{}{}}
\LANGpl{
\Question{
Obwód podstawy ostrosłupa prawidłowego sześciokątnego wynosi \( 12\sqrt3\,\mathrm{cm} \), a długość jego tworzącej jest równa \( 5\,\mathrm{cm} \). Oblicz pole powierzchni.}
{\( 48\sqrt3\,\mathrm{cm}^2 \)}{\( 72\sqrt3\,\mathrm{cm}^2 \)}{\( 30\sqrt3\,\mathrm{cm}^2 \)}{\( 96\sqrt3\,\mathrm{cm}^2 \)}{}{}}
\LANGcs{
\Question{
Pravidelný šestiboký jehlan má obvod podstavy  \( 12\sqrt3\,\mathrm{cm} \) a stěnovou výšku délky \( 5\,\mathrm{cm} \). Určete povrch jehlanu.}
{\( 48\sqrt3\,\mathrm{cm}^2 \)}{\( 72\sqrt3\,\mathrm{cm}^2 \)}{\( 30\sqrt3\,\mathrm{cm}^2 \)}{\( 96\sqrt3\,\mathrm{cm}^2 \)}{}{}}
\LANGsk{
\Question{
Pravidelný šesťboký ihlan má obvod podstavy  \( 12\sqrt3\,\mathrm{cm} \) a stenovú výšku dĺžky \( 5\,\mathrm{cm} \). Určte povrch ihlanu.}
{\( 48\sqrt3\,\mathrm{cm}^2 \)}{\( 72\sqrt3\,\mathrm{cm}^2 \)}{\( 30\sqrt3\,\mathrm{cm}^2 \)}{\( 96\sqrt3\,\mathrm{cm}^2 \)}{}{}}
\LANGes{
\Question{
La base de un prisma hexagonal regular tiene un perímetro de \( 12\sqrt3\,\mathrm{cm} \) y su apotema mide \( 5\,\mathrm{cm} \). Averigua la superficie de la pirámide.}
{\( 48\sqrt3\,\mathrm{cm}^2 \)}{\( 72\sqrt3\,\mathrm{cm}^2 \)}{\( 30\sqrt3\,\mathrm{cm}^2 \)}{\( 96\sqrt3\,\mathrm{cm}^2 \)}{}{}}
\End

\Data\ID{1103235606}
\Updated{2022-01-07 08:01:10}
\Level{C}
\SubArea{Volume and surface area formulas}
\IMG{1103235606.pdf}
\LANGen{
\Question{
Find the volume of a regular hexagonal prism with lateral edge length of \( 12\,\mathrm{cm} \) and a base edge length of \( 9\,\mathrm{cm} \) (see the picture).}
{\( 1458\sqrt3\,\mathrm{cm}^3 \)}{\( 243\sqrt3\,\mathrm{cm}^3 \)}{\( 1944\sqrt3\,\mathrm{cm}^3 \)}{\( 729\sqrt3\,\mathrm{cm}^3 \)}{}{}}
\LANGpl{
\Question{
Oblicz objętość graniastosłupa prawidłowego sześciokątnego, którego krawędź boczna jest równa \( 12\,\mathrm{cm} \), długość krawędzi podstawy wynosi \( 9\,\mathrm{cm} \) (spójrz na rysunek).}
{\( 1458\sqrt3\,\mathrm{cm}^3 \)}{\( 243\sqrt3\,\mathrm{cm}^3 \)}{\( 1944\sqrt3\,\mathrm{cm}^3 \)}{\( 729\sqrt3\,\mathrm{cm}^3 \)}{}{}}
\LANGcs{
\Question{
Pravidelný šestiboký hranol má délku boční hrany \( 12\,\mathrm{cm} \) a délku podstavné hrany \( 9\,\mathrm{cm} \) (viz obrázek). Určete objem hranolu.}
{\( 1458\sqrt3\,\mathrm{cm}^3 \)}{\( 243\sqrt3\,\mathrm{cm}^3 \)}{\( 1944\sqrt3\,\mathrm{cm}^3 \)}{\( 729\sqrt3\,\mathrm{cm}^3 \)}{}{}}
\LANGsk{
\Question{
Pravidelný šesťboký hranol má dĺžku bočnej hrany \( 12\,\mathrm{cm} \) a dĺžku hrany podstavy \( 9\,\mathrm{cm} \) (viď obrázok). Určte objem hranola.}
{\( 1458\sqrt3\,\mathrm{cm}^3 \)}{\( 243\sqrt3\,\mathrm{cm}^3 \)}{\( 1944\sqrt3\,\mathrm{cm}^3 \)}{\( 729\sqrt3\,\mathrm{cm}^3 \)}{}{}}
\LANGes{
\Question{
La altura de un prisma hexagonal regular es de \( 12\,\mathrm{cm} \) y el lado de su base mide \( 9\,\mathrm{cm} \) (observa el dibujo). Averigua su volumen.}
{\( 1458\sqrt3\,\mathrm{cm}^3 \)}{\( 243\sqrt3\,\mathrm{cm}^3 \)}{\( 1944\sqrt3\,\mathrm{cm}^3 \)}{\( 729\sqrt3\,\mathrm{cm}^3 \)}{}{}}
\End

\Data\ID{1103235607}
\Updated{2022-01-07 08:01:15}
\Level{C}
\SubArea{Volume and surface area formulas}
\IMG{1103235607.pdf}
\LANGen{
\Question{
Find the surface area of a regular hexagonal prism with lateral edge length of \( 10\sqrt3\,\mathrm{cm} \) and a base edge length of \( 6\,\mathrm{cm} \) (see the picture).}
{\( 468\sqrt3\,\mathrm{cm}^2 \)}{\( 414\sqrt3\,\mathrm{cm}^2 \)}{\( 168\sqrt3\,\mathrm{cm}^2 \)}{\( 408\sqrt3\,\mathrm{cm}^2 \)}{}{}}
\LANGpl{
\Question{
Oblicz pole powierzchni graniastosłupa prawidłowego sześciokątnego, którego krawędź boczna jest równa  \( 10\sqrt3\,\mathrm{cm} \), długość krawędzi podstawy to \( 6\,\mathrm{cm} \) (spójrz na rysunek).}
{\( 468\sqrt3\,\mathrm{cm}^2 \)}{\( 414\sqrt3\,\mathrm{cm}^2 \)}{\( 168\sqrt3\,\mathrm{cm}^2 \)}{\( 408\sqrt3\,\mathrm{cm}^2 \)}{}{}}
\LANGcs{
\Question{
Vypočítejte povrch pravidelného šestibokého hranolu s délkou boční hrany \( 10\sqrt3\,\mathrm{cm} \) a délkou podstavné hrany \( 6\,\mathrm{cm} \) (viz obrázek).}
{\( 468\sqrt3\,\mathrm{cm}^2 \)}{\( 414\sqrt3\,\mathrm{cm}^2 \)}{\( 168\sqrt3\,\mathrm{cm}^2 \)}{\( 408\sqrt3\,\mathrm{cm}^2 \)}{}{}}
\LANGsk{
\Question{
Vypočítajte povrch pravidelného šesťbokého hranolu s dĺžkou bočnej hrany \( 10\sqrt3\,\mathrm{cm} \) a dĺžkou hrany podstavy \( 6\,\mathrm{cm} \) (viď obrázok).}
{\( 468\sqrt3\,\mathrm{cm}^2 \)}{\( 414\sqrt3\,\mathrm{cm}^2 \)}{\( 168\sqrt3\,\mathrm{cm}^2 \)}{\( 408\sqrt3\,\mathrm{cm}^2 \)}{}{}}
\LANGes{
\Question{
Calcula la superficie del prisma hexagonal sabiendo que su altura mide \( 10\sqrt3\,\mathrm{cm} \) y la longitud del lado de la base es \( 6\,\mathrm{cm} \) (observa el dibujo).}
{\( 468\sqrt3\,\mathrm{cm}^2 \)}{\( 414\sqrt3\,\mathrm{cm}^2 \)}{\( 168\sqrt3\,\mathrm{cm}^2 \)}{\( 408\sqrt3\,\mathrm{cm}^2 \)}{}{}}
\End

\Data\ID{1103235608}
\Updated{2022-01-07 08:01:46}
\Level{C}
\SubArea{Volume and surface area formulas}
\IMG{1103235608.pdf}
\LANGen{
\Question{
The regular hexagonal prism has a volume of \( 324\sqrt3\,\mathrm{cm}^3 \) and the base edge length is equal to the prism’s height. Find the height.}
{\( 6\,\mathrm{cm} \)}{\( 6\sqrt6\,\mathrm{cm} \)}{\( 36\sqrt6\,\mathrm{cm} \)}{\( 6\sqrt3\,\mathrm{cm} \)}{}{}}
\LANGpl{
\Question{
Objętość graniastosłupa prawidłowego sześciokątnego jest równa \( 324\sqrt3\,\mathrm{cm}^3 \), długość krawędzi jego podstawy jest równa wysokości graniastosłupa. Oblicz wysokość.}
{\( 6\,\mathrm{cm} \)}{\( 6\sqrt6\,\mathrm{cm} \)}{\( 36\sqrt6\,\mathrm{cm} \)}{\( 6\sqrt3\,\mathrm{cm} \)}{}{}}
\LANGcs{
\Question{
Pravidelný šestiboký hranol má objem  \( 324\sqrt3\,\mathrm{cm}^3 \) a délka podstavné hrany je rovna výšce hranolu (viz obrázek). Určete výšku hranolu.}
{\( 6\,\mathrm{cm} \)}{\( 6\sqrt6\,\mathrm{cm} \)}{\( 36\sqrt6\,\mathrm{cm} \)}{\( 6\sqrt3\,\mathrm{cm} \)}{}{}}
\LANGsk{
\Question{
Pravidelný šesťboký hranol má objem  \( 324\sqrt3\,\mathrm{cm}^3 \) a dĺžka  hrany podstavy je rovná výške hranola (viď obrázok). Určte výšku hranola.}
{\( 6\,\mathrm{cm} \)}{\( 6\sqrt6\,\mathrm{cm} \)}{\( 36\sqrt6\,\mathrm{cm} \)}{\( 6\sqrt3\,\mathrm{cm} \)}{}{}}
\LANGes{
\Question{
El volumen de un prisma hexagonal regular es de \( 324\sqrt3\,\mathrm{cm}^3 \) y la longitud del lado de su base es igual a su altura (observa el dibujo). Averigua la altura del prisma.}
{\( 6\,\mathrm{cm} \)}{\( 6\sqrt6\,\mathrm{cm} \)}{\( 36\sqrt6\,\mathrm{cm} \)}{\( 6\sqrt3\,\mathrm{cm} \)}{}{}}
\End

\Data\ID{9000120308}
\Updated{2022-01-07 08:01:02}
\Level{C}
\SubArea{Volume and surface area formulas}
\IMG{001203_08-figure0.pdf}
\LANGen{
\Question{
The height \(v\) 
of a regular hexagonal prism is a double of its side
\(a\). The volume of
the prism is \(648\sqrt{3}\, \mathrm{cm}^{3}\).
Use this information to find the length of the longest solid diagonal in the
prism.}
{\(12\sqrt{2}\, \mathrm{cm}\)}{\(10\sqrt{6}\, \mathrm{cm}\)}{\(12\sqrt{6}\, \mathrm{cm}\)}{\(6\sqrt{10}\, \mathrm{cm}\)}{\(\sqrt{432}\, \mathrm{cm}\)}{}}
\LANGpl{
\Question{
Wysokość  \(v\) graniastosłupa prawidłowego sześciokątnego jest dwukrotnie większa od długości boku
\(a\). Objętość graniastosłupa jest równa
 \(648\sqrt{3}\, \mathrm{cm}^{3}\). Oblicz długość  jego najdłuższej przekątnej.}
{\(12\sqrt{2}\, \mathrm{cm}\)}{\(10\sqrt{6}\, \mathrm{cm}\)}{\(12\sqrt{6}\, \mathrm{cm}\)}{\(6\sqrt{10}\, \mathrm{cm}\)}{\(\sqrt{432}\, \mathrm{cm}\)}{}}
\LANGcs{
\Question{
Pravidelný šestiboký hranol o objemu
\(648\sqrt{3}\, \mathrm{cm}^{3}\) má
výšku dvakrát větší než délka podstavné hrany. Nejdelší tělesová
úhlopříčka má délku:}
{\(12\sqrt{2}\, \mathrm{cm}\)}{\(10\sqrt{6}\, \mathrm{cm}\)}{\(12\sqrt{6}\, \mathrm{cm}\)}{\(6\sqrt{10}\, \mathrm{cm}\)}{\(\sqrt{432}\, \mathrm{cm}\)}{}}
\LANGsk{
\Question{
Pravidelný šesťboký hranol s objemom
\(648\sqrt{3}\, \mathrm{cm}^{3}\) má
výšku dvakrát väčšiu ako dĺžka hrany podstavy. Najdlhšia telesová
uhlopriečka má dĺžku:}
{\(12\sqrt{2}\, \mathrm{cm}\)}{\(10\sqrt{6}\, \mathrm{cm}\)}{\(12\sqrt{6}\, \mathrm{cm}\)}{\(6\sqrt{10}\, \mathrm{cm}\)}{\(\sqrt{432}\, \mathrm{cm}\)}{}}
\LANGes{
\Question{
Sea un prisma hexagonal regular, cuyo volumen es
\(648\sqrt{3}\, \mathrm{cm}^{3}\). Su altura es dos veces más larga que el lado de su base. La diagonal interiorl más larga mide:}
{\(12\sqrt{2}\, \mathrm{cm}\)}{\(10\sqrt{6}\, \mathrm{cm}\)}{\(12\sqrt{6}\, \mathrm{cm}\)}{\(6\sqrt{10}\, \mathrm{cm}\)}{\(\sqrt{432}\, \mathrm{cm}\)}{}}
\End
